%=========================================================================
\chapter{The WIR Hardware Model}
\label{chap:hwmodel}
\renewcommand{\sectiontitle}{}
%=========================================================================

This chapter describes and specifies the structure and \gls{API} of all classes that serve to model hardware in some \gls{WIR} system.
Thus, the subsequent \crefrange{sec:processors}{sec:buses} focus in detail on the complete hierarchy of system components from processors via memory regions to system buses as depicted in \cref{fig:hwmodel1}.

Class \texttt{WIR\_SystemComponent} depicted in \cref{fig:hwmodel1} serves as abstract base class for all the hardware components described in the following.
In analogy to parameters, the method \texttt{getType} (cf. \cref{lst:componenttype}) can be used to quickly check whether some actual object within the \gls{WIR} hardware model is, e.g., a processor core or a cache.

\begin{lstlisting}[style=code,float=ht,caption={Generic Type Handling of Class \texttt{WIR\_SystemComponent}},label=lst:componenttype]
// Get a system component's type.
virtual WIR_SystemComponentType WIR_SystemComponent::getType@( void )@ const @=@ §0§;Ü\label{lst:componenttype1}Ü

// Enum WIR_SystemComponentType represents different types of WIR parameters.
enum class WIR_SystemComponentType : @char@Ü\label{lst:componenttype2}Ü
{
  bus,      // A system bus.
  cache,    // A cache.
  core,     // A processor core.
  memory,   // A physical memory region.
  section,  // An ELF executable section.
  system    // An entire system.
};
\end{lstlisting}


%-------------------------------------------------------------------------
\section{Processors}
\label{sec:processors}
%-------------------------------------------------------------------------

\gls{WIR} processors model both the software and the hardware characteristics of actual processor cores, i.e., their \glspl{ISA}, register files, clock frequencies, logical memory sections etc.
The \gls{WIR} library is designed to support an arbitrary number of different processor models at the same time, without the need to rebuild the library.
The rationale behind this design is to enable the modelling of heterogeneous multi-core systems consisting of various processor cores of different type.
As already described in \cref{sec:keyfeatures}, \gls{WIR} features models for a basic ARMv4 core, MIPS SPIM, RISC-V RV32I, and TriCore TC1.3, plus all the derived models listed in \cref{tab:AnaOptISA} that inherit from these base models.
The simplified inheritance diagram of this basic class hierarchy is depicted in \cref{fig:wirprocessors}.

\begin{figure}[ht]
  \begin{center}
    \includegraphics[width=0.84\textwidth]{processors.png}
  \end{center}
  \caption{Simplified Inheritance Diagram of WIR Processor Models}
  \label{fig:wirprocessors}
\end{figure}

The following \cref{ssec:procmodel} presents the general concepts of the base classes \texttt{WIR\_BaseProcessor} and \texttt{WIR\_\-Processor} in more detail.
Hereafter, \crefrange{ssec:armmodel}{ssec:tcmodel} provide a short overview about the actual models for the MIPS, TriCore, RISC-V and ARM architectures.


% - - - - - - - - - - - - - - - - - - - - - - - - - - - - - - - - - - - -
\subsection{Generic Processor Model}
\label{ssec:procmodel}
% - - - - - - - - - - - - - - - - - - - - - - - - - - - - - - - - - - - -

Class \texttt{WIR\_BaseProcessor} comes with the usual set of constructors as shown in \crefrange{lst:procconstruction1}{lst:procconstruction3} of \cref{lst:procconstruction}.
Furthermore, processor models are characterized by a few descriptive attributes storing its names, clock frequency and supply voltage.
Methods \texttt{getProcessorName} (cf. \cref{lst:procconstruction4}) provides a string containing a processor model's general name (e.g., ``RISC-V''), while the string returned by \texttt{getISAName} (cf. \cref{lst:procconstruction5}) describes its specific \gls{ISA} version (e.g., ``RV32I V2.0'').
A processor's clock frequency can be retrieved using \texttt{getClockFrequency}, and \texttt{getVoltage} provides its supply voltage.
Both values can be relevant for \gls{WIR} based timing or energy analyses.

\begin{lstlisting}[style=code,float=h,caption={Constructors and Methods for General Attribute Handling of Class \texttt{WIR\_BaseProcessor}},label=lst:procconstruction]
// Default constructor creating an empty processor.
WIR_BaseProcessor::WIR_BaseProcessor@( void )@;Ü\label{lst:procconstruction1}Ü

// Copy constructor.
WIR_BaseProcessor::WIR_BaseProcessor@(@ const WIR_BaseProcessor æ&æp @)@;Ü\label{lst:procconstruction2}Ü

// Move constructor.
WIR_BaseProcessor::WIR_BaseProcessor@(@ WIR_BaseProcessor æ&&æp @)@;Ü\label{lst:procconstruction3}Ü

// Get the processor's descriptive name.
ästd::stringä WIR_BaseProcessor::getProcessorName@( void )@ const;Ü\label{lst:procconstruction4}Ü

// Get the ISA's descriptive name.
ästd::stringä WIR_BaseProcessor::getISAName@( void )@ const;Ü\label{lst:procconstruction5}Ü

// Get the processor's clock frequency in Hertz.
@unsigned long long@ WIR_BaseProcessor::getClockFrequency@( void )@ const;Ü\label{lst:procconstruction6}Ü

// Get the processor's supply voltage.
@float@ WIR_BaseProcessor::getVoltage@( void )@ const;Ü\label{lst:procconstruction7}Ü
\end{lstlisting}

As motivated earlier, \gls{WIR} is designed to cope with various processor types simultaneously.
For this sole purpose, the template class \texttt{WIR\_Processor<typename T>} has been introduced.
It inherits from class \texttt{WIR\_BasePro\-cessor}, features exactly the same constructors, and extends this base class by type handling mechanisms which allow to distinguish, e.g., an ARM processor within \gls{WIR} from a RISC-V one.
In order to do so, the type names of the derived, actual processors are made explicit using the class' template argument \texttt{T} (cf. \cref{fig:wirprocessors}).
Then, the methods \texttt{getProcessorTypeID} or \texttt{getProcessorType} (cf. \cref{lst:proctypes}) can be used to obtain the unique, numerical identifier of the processor model.
Both methods behave absolutely identically, they only differ in the sense that \texttt{getProcessorTypeID} is static (cf. \cref{lst:proctypes1}), while \texttt{getProcessorType} (cf. \cref{lst:proctypes2}) can only be used for actual instances of processor models.

\begin{lstlisting}[style=code,float=t,caption={Methods for Type Handling of Class \texttt{WIR\_Processor}},label=lst:proctypes]
// Get the unique ID of a processor type.
static WIR_id_t WIR_Processor<T>::getProcessorTypeID@( void )@;Ü\label{lst:proctypes1}Ü

// Get the unique ID of a processor type.
virtual WIR_id_t WIR_Processor<T>::getProcessorType@( void )@ const;Ü\label{lst:proctypes2}Ü
\end{lstlisting}

A \gls{WIR} processor model consists of the specification of a core's register file(s), its available constant immediates, the supported opcodes and operation formats, addressing modes and condition fields, plus the handling of logical sections (e.g., for code, data or constants) in memory.
All these aspects are discussed in detail in the following.


% -  -  -  -  -  -  -  -  -  -  -  -  -  -  -  -  -  -  -  -  -  -  -  -  -
\subsubsection{Register Types}
\label{sssec:regtypes}
% -  -  -  -  -  -  -  -  -  -  -  -  -  -  -  -  -  -  -  -  -  -  -  -  -

A processor core features a certain amount of physical hardware registers which might be grouped into several register files and might serve for dedicated purposes.
\cref{fig:hwmodel2} already showed the example of the \texttt{TC13} model including different types of address, data and extended data registers.
Since registers of such different types typically cannot be used interchangably, the \gls{WIR} hardware model must be able to distinguish between different register types.
This is realized by the generic base class \texttt{RegisterType} which is a sub-class of \texttt{WIR\_BaseProcessor} as shown in \cref{lst:regtypes}.

\begin{lstlisting}[style=code,float=h,caption={Generic Register Type Handling of Class \texttt{WIR\_BaseProcessor}},label=lst:regtypes]
class WIR_BaseProcessor : public WIR_SystemComponent
{
  public:

    [...]

    class RegisterType : public WIR_InheritableEnumÜ\label{lst:regtypes1}Ü
    {
      public:

        // Get a register type's bit width.
        @unsigned int@ getBitWidth@( void )@ const;Ü\label{lst:regtypes2}Ü

        // Get a register type's prefix strings used to create a register's name.
        const ästd::stringä ( æ&æ getPrefixes@( void )@ const )æ[æ §2§ æ]æ;Ü\label{lst:regtypes3}Ü

        // Get a register type's suffix strings used to create a register's name.
        const ästd::stringä ( æ&æ getSuffixes@( void )@ const )æ[æ §2§ æ]æ;Ü\label{lst:regtypes4}Ü

        [...]

      protected:

        // Default constructor creating a register type with given prefix and suffix strings and bit width.
        RegisterType@(@ const ästd::stringä æ&æpp, const ästd::stringä æ&æpv,                                                             const ästd::stringä æ&æsp, const ästd::stringä æ&æsv, @unsigned int@ w @)@;Ü\label{lst:regtypes5}Ü
    };
};
\end{lstlisting}

One crucial aspect of registers of different types are their potentially differing bit widths.
For the \texttt{TC13} example from \cref{fig:hwmodel2}, address and data registers are 32 bits wide each, while extended data registers have a width of 64 bits.
The bit width of a certain register type can thus be retrieved using method \texttt{getBitWidth} (cf. \cref{lst:regtypes2}).

Regularly, all the registers of a certain type follow a given naming scheme where some characters or strings are used to specify the register type, together with the actual unique number of a register.
Coming back to the running \texttt{TC13} example, the names of TriCore address registers always begin with ``\texttt{a}'', the prefix ``\texttt{d}'' characterizes data registers, and ``\texttt{e}'' is used for extended data registers.
For the proper and \gls{ISA}-compliant output of \gls{WIR} assembly code, these naming schemes of register types also need to be supported.
For this purpose, method \texttt{getPrefixes} (cf. \cref{lst:regtypes3}) returns a reference to an array of two strings.
These strings are both prefixes that are put before an actual register's number.
The first string inside the returned array denotes the prefix to be used for physical registers of some register type, while the second string is used as prefix for virtual registers.
This way, physical and virtual registers can clearly be recognized in \gls{WIR} assembly output.
Similarly, method \texttt{getSuffixes} (cf. \cref{lst:regtypes4}) provides two strings used as suffixes for physical or virtual register names, resp.

This characteristic information about register types is provided during their construction.
As can be seen from \cref{lst:regtypes5} of \cref{lst:regtypes}, the constructor expects all the strings for prefixes and suffixes of both physical and virtual registers as arguments (e.g., argument \texttt{sp} stands for the suffix string of physical registers), as well as a register type's bit width.

\begin{lstlisting}[style=code,float=t,caption={Concrete Register Type Handling (here for Class \texttt{TC13})},label=lst:tcregtypes]
class TC13 : public WIR_Processor<TC13>
{
  public:

    [...]

    class RegisterType : public WIR_Processor::RegisterTypeÜ\label{lst:tcregtypes1}Ü
    {
      public:

        // Declaration of address registers.
        static const RegisterType aReg;Ü\label{lst:tcregtypes2}Ü

        // Declaration of data registers.
        static const RegisterType dReg;Ü\label{lst:tcregtypes3}Ü

        // Declaration of extended data registers.
        static const RegisterType eReg;Ü\label{lst:tcregtypes4}Ü

        [...]
    };

    [...]
};

// Definition of register types.
const TC13::RegisterType TC13::RegisterType::aReg { "a", "a_", "", "", §32§ };Ü\label{lst:tcregtypes5}Ü
const TC13::RegisterType TC13::RegisterType::dReg { "d", "d_", "", "", §32§ };Ü\label{lst:tcregtypes6}Ü
const TC13::RegisterType TC13::RegisterType::eReg { "e", "e_", "", "", §64§ };Ü\label{lst:tcregtypes7}Ü
\end{lstlisting}

The class \texttt{RegisterType} itself inherits from a simple \gls{WIR}-internal support class named \texttt{WIR\_InheritableEnum} (cf. \cref{lst:regtypes1}).
The sole purpose of this support class is to enable the simple enumeration of all available register types in a \gls{WIR} processor model, while simultaneously being able to inherit and extend such an enumeration in derived classes (which is impossible with standard C/C++ \texttt{enum} types).
This procedure is illustrated in \cref{lst:tcregtypes}, again for the \texttt{TC13} model as example.
As can be seen from \cref{lst:tcregtypes1}, the TriCore register types are derived from the base class described earlier.
Class \texttt{TC13::RegisterType} then contains the static enumeration and declaration of actual register types (cf. \crefrange{lst:tcregtypes2}{lst:tcregtypes4}).
These declarations are finally defined and implemented in \crefrange{lst:tcregtypes5}{lst:tcregtypes7} with appropriate calls of the inherited constructor of class \texttt{WIR\_BaseProcessor::RegisterType}.
In the given example, physical data registers use prefix ``\texttt{d}'', while the names of virtual data registers begin with ``\texttt{d\_}''.


% -  -  -  -  -  -  -  -  -  -  -  -  -  -  -  -  -  -  -  -  -  -  -  -  -
\subsubsection{Registers}
\label{sssec:regs}
% -  -  -  -  -  -  -  -  -  -  -  -  -  -  -  -  -  -  -  -  -  -  -  -  -

Any register within \gls{WIR} -- irrespective of some actual processor model, or be it a virtual or physical register -- is generically modeled by class \texttt{WIR\_BaseRegister}.
This class comes with the constructors depicted in \cref{lst:regconstruction}.
As can be seen from \cref{lst:regconstr1}, the construction of a register always requires a reference to its register type as introduced in \cref{sssec:regtypes}.

\begin{lstlisting}[style=code,float=ht,caption={Constructors of Class \texttt{WIR\_BaseRegister}},label=lst:regconstruction]
// Default constructor creating a register of a given type.
explicit WIR_BaseRegister::WIR_BaseRegister@(@ const WIR_BaseProcessor::RegisterType æ&æt @)@;Ü\label{lst:regconstr1}Ü

// Copy constructor.
WIR_BaseRegister::WIR_BaseRegister@(@ const WIR_BaseRegister æ&ær @)@;Ü\label{lst:regconstr2}Ü

// Move constructor.
WIR_BaseRegister::WIR_BaseRegister@(@ WIR_BaseRegister æ&&ær @)@;Ü\label{lst:regconstr4}Ü
\end{lstlisting}

The attributes of the given register type are then used by the methods from \cref{lst:regtype}. \texttt{getType} (cf. \cref{lst:regtype1}) provides a reference to a register's type, while \texttt{getBitWidth} (\cref{lst:regtype2}) and \texttt{getName} (\cref{lst:regtype3}) return the bit width and name as discussed earlier.

\begin{lstlisting}[style=code,float=ht,caption={Type and Name Handling of Class \texttt{WIR\_BaseRegister}},label=lst:regtype]
// Get a register's actual type.
WIR_BaseProcessor::RegisterType æ&æWIR_BaseRegister::getType@( void )@;Ü\label{lst:regtype1}Ü

// Get a register's bit width.
@unsigned int@ WIR_BaseRegister::getBitWidth@( void )@;Ü\label{lst:regtype2}Ü

// Get a register's name.
ästd::stringä WIR_BaseRegister::getName@( void )@;Ü\label{lst:regtype3}Ü
\end{lstlisting}

Many processor architectures feature so-called hierarchical registers.
Here, two (or more) registers of a type with $n$ bits each are grouped together so that they form a large register of, say, $2 \ast n$ bits.
The two basic registers thus become part of a hierarchy where the larger parent register spans over its child registers.
In principle, several of these parent registers could also be grouped together to form an even larger hierarchy of registers, even though this is not common in real architectures.
Anyways, \gls{WIR} registers support the modelling of arbitrary such hierarchies of registers.

\cref{lst:reghier} shows the basic methods for the handling of such register hierarchies provided by class \texttt{WIR\_Base\-Register}.
\texttt{isChild} (cf. \cref{lst:reghier1}) returns whether a register is child of some register at all, while \texttt{isChildOf} (\cref{lst:reghier2}) can be used to check the child relation between two specific registers.
Methods \texttt{hasChilds} (\cref{lst:reghier3}) and \texttt{containsChild} (\cref{lst:reghier4}) can be used to obtain the reverse information, i.e., whether a register contains child registers at all, or whether one register contains a specific other register given by its unique numerical ID as child.
Furthermore, \texttt{isHierarchical} (\cref{lst:reghier5}) returns whether a register is part of some hierarchy, i.e., if it is a child register or if it has childs.
Finally, \texttt{getLeafs} (\cref{lst:reghier6}) can be used to inspect all leaf registers underneath a given register.
It returns a vector with references to all lower-level registers in the hierarchy that do not have child registers in turn.

\begin{lstlisting}[style=code,float=t,caption={Register Hierarchy Handling of Class \texttt{WIR\_BaseRegister}},label=lst:reghier]
// Determine whether a register is child of some register, i.e., if it is part of a register hierarchy.
@bool@ WIR_BaseRegister::isChild@( void )@ const;Ü\label{lst:reghier1}Ü

// Determine whether this register is a child of a given register in a register hierarchy.
@bool@ WIR_BaseRegister::isChildOf@(@ const WIR_BaseRegister æ&ær @)@ const;Ü\label{lst:reghier2}Ü

// Determine whether a register has child registers.
virtual @bool@ WIR_BaseRegister::hasChilds@( void )@ const;Ü\label{lst:reghier3}Ü

// Determine whether a register contains a child register with the given ID.
virtual @bool@ WIR_BaseRegister::containsChild@(@ WIR_id_t id @)@ const;Ü\label{lst:reghier4}Ü

// Determine whether a register is part of a register hierarchy, i.e., if it has some childs or is a child.
@bool@ WIR_BaseRegister::isHierarchical@( void )@ const;Ü\label{lst:reghier5}Ü

// Get all leafs of the hierarchy underneath this register.
const ästd::vectoräæ<æästd::reference_wrapperäæ<æWIR_BaseRegisteræ>>æ WIR_BaseRegister::getLeafs@( void )@ const;Ü\label{lst:reghier6}Ü
\end{lstlisting}

As already introduced earlier, \gls{WIR} generally differentiates between physical and virtual registers.
For this purpose, class \texttt{WIR\_BaseRegister} furthermore provides methods to distinguish between physical and virtual registers (cf. \cref{lst:phvreg}).
As their names suggest, method \texttt{isPhysical} (\cref{lst:phvreg1}) returns \texttt{true} only for physical registers, while the opposite holds for \texttt{isVirtual} (\cref{lst:phvreg2}).

\begin{lstlisting}[style=code,float=ht,caption={Distinction Between Physical and Virtual Registers in Class \texttt{WIR\_BaseRegister}},label=lst:phvreg]
// Determine whether a register is physical.
@bool@ WIR_BaseRegister::isPhysical@( void )@ const;Ü\label{lst:phvreg1}Ü

// Determine whether a register is virtual.
virtual @bool@ WIR_BaseRegister::isVirtual@( void )@ const;Ü\label{lst:phvreg2}Ü
\end{lstlisting}

In addition, \gls{WIR} models physical and virtual registers as specialized classes \texttt{WIR\_PhysicalRegister} and \texttt{WIR\_VirtualRegister}, which both inherit from \texttt{WIR\_BaseRegister}.
Everything that has been discussed about registers up to this point holds for both of these derived classes.
In their specific details, however, both physical and virtual registers slightly differ from each other.

The characteristic features of physical registers are represented by the methods shown in \cref{lst:phreg}.
The most relevant aspect of physical registers is that they are tightly connected to some processor model within \gls{WIR}.
Thus, method \texttt{getProcessor} (\cref{lst:phreg1}) can be used to retrieve a reference to the base processor containing that physical register.
One physical register that plays an exceptional role for almost all processor architectures is the stack pointer.
In order to quickly determine whether some physical register is an architecture's stack pointer, method \texttt{isStackPointer} (\cref{lst:phreg2}) can be used.

Register hierarchies as introduced earlier in this section can only consist of either physical or virtual registers.
Thus, the generic \gls{API} for register hierarchy handling from \cref{lst:reghier} requires some extensions specifically for physical registers.
Method \texttt{getChilds} (cf. \cref{lst:phreg3} in \cref{lst:phreg}) returns the list of references to all physical child registers that a physical register has.
This list comes in little-endian order, i.e., the least-significant child register is the first element in the returned list.
The \texttt{containsChild} variant from \cref{lst:phreg4} can be used to check whether a given physical register is child of another one.
In analogy to the \gls{WIR} list \gls{API} desribed in \cref{sec:listlookup}, the two \texttt{find} methods from \cref{lst:phreg5,lst:phreg6} can be used to find some physical register in the list of child registers of some other physical register.
In order to navigate upwards within a hierarchy of physical registers, class \texttt{WIR\_PhysicalRegister} provides the methods \texttt{getParent} (\cref{lst:phreg7}) and \texttt{getRoot} (\cref{lst:phreg8}).
The former returns a reference to a register's immediate parent, i.e., to that physical register that contains this register as direct child.
If \texttt{getParent} is used for a register that has no such parent, it returns a reference to the very same register.
The latter provides a reference to the overall root of a register hierarchy, i.e., to that physical register that has no more parent.

\begin{lstlisting}[style=code,float=t,caption={Characteristic Features of Class \texttt{WIR\_PhysicalRegister}},label=lst:phreg]
// Get the processor model to which a physical register belongs.
WIR_BaseProcessor æ&æWIR_PhysicalRegister::getProcessor@( void )@ const;Ü\label{lst:phreg1}Ü

// Determine whether a register is a processor's stack pointer.
@bool@ WIR_PhysicalRegister::isStackPointer@( void )@ const;Ü\label{lst:phreg2}Ü

// Get the list of all child registers, in little-endian order.
const ästd::listäæ<æästd::reference_wrapperäæ<æWIR_PhysicalRegisteræ>> &æWIR_PhysicalRegister::getChilds@( void )@ const;Ü\label{lst:phreg3}Ü

// Determine whether a register contains a given physical register as child.
@bool@ WIR_PhysicalRegister::containsChild@(@ const WIR_PhysicalRegister æ&ær @)@ const;Ü\label{lst:phreg4}Ü

// Find the child register with the given ID.
ästd::listäæ<æästd::reference_wrapperäæ<æWIR_PhysicalRegisteræ>>::const_iteratoræ WIR_PhysicalRegister::findChild@(@ WIR_id_t id @)@ const;Ü\label{lst:phreg5}Ü

// Find the given child register.
ästd::listäæ<æästd::reference_wrapperäæ<æWIR_PhysicalRegisteræ>>::const_iteratoræ WIR_PhysicalRegister::findChild@(@ const WIR_PhysicalRegister æ&ær @)@ const;Ü\label{lst:phreg6}Ü

// Get a reference to a register's immediate parent in the hierarchy, or a reference to itself.
WIR_PhysicalRegister æ&æWIR_PhysicalRegister::getParent@( void )@ const;Ü\label{lst:phreg7}Ü

// Get a reference to a register's top-level root in the hierarchy, or a reference to itself.
WIR_PhysicalRegister æ&æWIR_PhysicalRegister::getRoot@( void )@ const;Ü\label{lst:phreg8}Ü
\end{lstlisting}

\cref{sec:functions} has already discussed the connection between virtual registers and \gls{WIR} functions, as well as the concept of precoloring virtual with physical registers.
Thus, virtual registers are conceptually owned by functions, and not by processor models, unlike physical registers.
\cref{lst:vreg} finally shows the characteristic features of virtual registers, and method \texttt{getFunction} (\cref{lst:vreg1}) can be used to retrieve the function to which a virtual register belongs.
Methods \texttt{isPrecolored} (\cref{lst:vreg2}) and \texttt{getPrecolor} (\cref{lst:vreg3}) help to check if a virtual register is precolored and, if so, to obtain the physical register serving as precolor of a virtual register.
If a virtual register is not precolored, \texttt{getPrecolor} fails with a runtime error.
The methods to handle hierarchies of virtual registers from \crefrange{lst:vreg4}{lst:vreg9} are equivalent to those of class \texttt{WIR\_PhysicalRegister}.

\begin{lstlisting}[style=code,float=t,caption={Characteristic Features of Class \texttt{WIR\_VirtualRegister}},label=lst:vreg]
// Get the function to which a virtual register belongs.
WIR_Function æ&æWIR_VirtualRegister::getFunction@( void )@ const;Ü\label{lst:vreg1}Ü

// Determine whether a virtual register is precolored.
@bool@ WIR_VirtualRegister::isPrecolored@( void )@ const;Ü\label{lst:vreg2}Ü

// Get the physical register associated with a virtual register.
WIR_PhysicalRegister æ&æWIR_VirtualRegister::getPrecolor@( void )@ const;Ü\label{lst:vreg3}Ü

// Get the list of all child registers, in little-endian order.
const ästd::listäæ<æästd::reference_wrapperäæ<æWIR_VirtualRegisteræ>> &æWIR_VirtualRegister::getChilds@( void )@ const;Ü\label{lst:vreg4}Ü

// Determine whether a register contains a given virtual register as child.
@bool@ WIR_VirtualRegister::containsChild@(@ const WIR_VirtualRegister æ&ær @)@ const;Ü\label{lst:vreg5}Ü

// Find the child register with the given ID.
ästd::listäæ<æästd::reference_wrapperäæ<æWIR_VirtualRegisteræ>>::const_iteratoræ WIR_VirtualRegister::findChild@(@ WIR_id_t id @)@ const;Ü\label{lst:vreg6}Ü

// Find the given child register.
ästd::listäæ<æästd::reference_wrapperäæ<æWIR_VirtualRegisteræ>>::const_iteratoræ WIR_VirtualRegister::findChild@(@ const WIR_VirtualRegister æ&ær @)@ const;Ü\label{lst:vreg7}Ü

// Get a reference to a register's immediate parent in the hierarchy, or a reference to itself.
WIR_VirtualRegister æ&æWIR_VirtualRegister::getParent@( void )@ const;Ü\label{lst:vreg8}Ü

// Get a reference to a register's top-level root in the hierarchy, or a reference to itself.
WIR_VirtualRegister æ&æWIR_VirtualRegister::getRoot@( void )@ const;Ü\label{lst:vreg9}Ü
\end{lstlisting}

After having presented the generic concepts of both physical and virtual \gls{WIR} registers up to this point, only the connection between these generic base classes -- \texttt{WIR\_PhysicalRegister} and \texttt{WIR\_VirtualRegister} -- and actual \gls{WIR} processor models remains to be discussed.
For each register type of a \gls{WIR} processor model, dedicated classes exist that model both physical and virtual registers of that processor-specific register type.
Naturally, these processor-specific register classes are inherited from the base classes \texttt{WIR\_PhysicalRegister} and \texttt{WIR\_VirtualRegister}, resp.

We again fall back to the \texttt{TC13} example already used earlier in this section.
Within the TriCore architecture, two 32-bit data registers are tied together and form a 64-bit extended data register.
Thus, \gls{WIR} on the one hand comes with the class \texttt{TC\_DRegV} modelling virtual data registers of the TriCore architecture, cf. \cref{lst:tcdregv}.
The specification of such an actual virtual register class is extremely simple and basically boils down to the constructor from \cref{lst:tcdregv1} that requires no arguments at all and simply invokes the constructor of \texttt{WIR\_VirtualRegister} with the TriCore-specific register type \texttt{dReg}.

\begin{lstlisting}[style=code,float=ht,caption={Concrete Modelling of Virtual Data Registers (here for TriCore architecture)},label=lst:tcdregv]
class TC_DRegV final : public WIR_VirtualRegister
{
  public:

    // Default constructor for virtual TriCore data registers.
    TC_DRegV@( void )@ : WIR_VirtualRegister { TC13::RegisterType::dReg }Ü\label{lst:tcdregv1}Ü
    {};

    [...]
};
\end{lstlisting}

On the other hand, the TriCore model requires a class \texttt{TC\_ERegV} modelling virtual extended data registers, which is conceptually shown in \cref{lst:tceregv}.
Its structure is as simple as before, with the key difference that its constructor on the fly generates two virtual data registers that serve as childs of the extended data register to be created (cf. \cref{lst:tceregv2,lst:tceregv3}).

\begin{lstlisting}[style=code,float=t,caption={Concrete Modelling of Virtual Extended Data Registers (here for TriCore architecture)},label=lst:tceregv]
class TC_ERegV final : public WIR_VirtualRegister
{
  public:

    // Default constructor for virtual TriCore extended data registers.
    TC_ERegV@( void )@ : WIR_VirtualRegister { TC13::RegisterType::eReg }Ü\label{lst:tceregv1}Ü
    {
      // Create two virtual child data registers.
      pushBackChild@(@ TC_DRegV@() )@;Ü\label{lst:tceregv2}Ü
      pushBackChild@(@ TC_DRegV@() )@;Ü\label{lst:tceregv3}Ü
    };

    [...]
};
\end{lstlisting}

Very similarly, class \texttt{TC\_DRegP} modelling physical data registers is derived as depicted in \cref{lst:tcdregp}.
Its constructor (cf. \cref{lst:tcdregp2}) gets a string characterizing the TriCore's register name as argument, as well as a Boolean flag denoting whether this will be the stack pointer register.
This information is then simply delegated to class \texttt{WIR\_PhysicalRegister}, together with the TriCore register type \texttt{dReg}.
The only difference between processor-specific classes for virtual and physical registers is that the constructor is declared private in the latter case.
This has the effect that application code using the \gls{WIR} library is unable to create physical registers itself, while this of course is possible for virtual registers.
The rationale behind this design is that only processor models within \gls{WIR} shall be able to create processor-specific physical registers.
For this reason, the \texttt{TC13} model is declared as friend in \cref{lst:tcdregp1} of \cref{lst:tcdregp}.

\begin{lstlisting}[style=code,float=ht,caption={Concrete Modelling of Physical Data Registers (here for TriCore architecture)},label=lst:tcdregp]
class TC_DRegP final : public WIR_PhysicalRegister
{
  private:

    friend class TC13;Ü\label{lst:tcdregp1}Ü

    // Default constructor for physical TriCore data registers.
    explicit TC_DRegP@(@ const ästd::stringä æ&æn, @bool@ sp æ=æ false @)@ :                                                WIR_PhysicalRegister { TC13::RegisterType::dReg, n, sp }Ü\label{lst:tcdregp2}Ü
    {};

    [...]
};
\end{lstlisting}

\cref{lst:tcphregs} finally shows the overall connection between a \gls{WIR} processor model and its various classes modelling physical registers.
Whenever a new instance of a processor model is constructed, this constructor is responsible to also create all the physical registers that the modeled architecture actually features.
\cref{lst:tcphregs1} of \cref{lst:tcphregs} shows the declaration of the constructor of class \texttt{TC13} which simply receives no arguments at all and arranges all aspects of the TriCore model silently under the hood.
Thus, exactly this constructor \texttt{TC13::TC13} repetitively calls the constructor of class \texttt{TC\_DRegP} from \cref{lst:tcdregp2} of \cref{lst:tcdregp} in order to set up all physical data registers D0, \dots, D15 of the TriCore architecture.
Likewise, \texttt{TC13::TC13} constructs the TriCore's physical extended data registers E0, \dots, E14.
It furthermore ensures that the correct pairs of physical data registers are added as child registers so that all register hierarchies are properly represented.

\begin{lstlisting}[style=code,float=t,caption={Physical Register Handling (here for Class \texttt{TC13})},label=lst:tcphregs]
class TC13 : public WIR_Processor<TC13>
{
  public:

    // Default constructor for TC13 processor architectures.
    TC13@( void )@;Ü\label{lst:tcphregs1}Ü

    [...]

    // Access to physical data register D0.
    const TC_DRegP æ&æD0@( void )@ const;Ü\label{lst:tcphregs2}Ü

    // Access to physical data register D1.
    const TC_DRegP æ&æD1@( void )@ const;Ü\label{lst:tcphregs3}Ü

    [...]

    // Access to physical extended data register E0.
    const TC_ERegP æ&æE0@( void )@ const;Ü\label{lst:tcphregs4}Ü

    // Access to physical extended data register E2.
    const TC_ERegP æ&æE2@( void )@ const;Ü\label{lst:tcphregs5}Ü

    [...]
};
\end{lstlisting}

For the convenient access to its various physical registers, a \gls{WIR} processor model provides specific methods shown in \crefrange{lst:tcphregs2}{lst:tcphregs5} of \cref{lst:tcphregs}.
Here, \texttt{TC13::D1} (\cref{lst:tcphregs3}) returns a reference to the physical data register D1 that has been created by the constructor of class \texttt{TC13} as described above.
Similarly, a reference to extended data register E0 can be obtained using method \texttt{TC13::E0} (\cref{lst:tcphregs4}).
The register hierarchy underneath this register E0 can finally be explored using the methods already presented in \cref{lst:phreg}.
In particular, method \texttt{getChilds} from \cref{lst:phreg3} returns a list consisting of those physical registers D0 and D1 that are also provided by methods \texttt{TC13::D0} and \texttt{TC13::D1} (cf. \cref{lst:tcphregs2,lst:tcphregs3} of \cref{lst:tcphregs}).
This way, all physical registers and their hierarchies are properly modeled within \gls{WIR} processor models.

Structured access to all physical registers of a processor model is provided by the first variant of \texttt{getPhRegs} of the generic base class \texttt{WIR\_Base\-Processor} shown in \cref{lst:phregaccess1} of \cref{lst:phregaccess}
The second variant shown in \cref{lst:phregaccess2} can be used to conveniently obtain all physical registers of a given type.

\begin{lstlisting}[style=code,float=h,caption={Physical Register Access in Class \texttt{WIR\_BaseProcessor}},label=lst:phregaccess]
// Get a vector containing all physical registers of a processor architecture.
const ästd::vectoräæ<æästd::reference_wrapperäæ<æconst WIR_PhysicalRegisteræ>>æ WIR_BaseProcessor::getPhRegs@( void )@ const;Ü\label{lst:phregaccess1}Ü

// Get a vector containing all physical registers of a processor architecture of the specified type.
const ästd::vectoräæ<æästd::reference_wrapperäæ<æconst WIR_PhysicalRegisteræ>>æ WIR_BaseProcessor::getPhRegs@(@       const RegisterType æ&æt @)@ const;Ü\label{lst:phregaccess2}Ü
\end{lstlisting}


% -  -  -  -  -  -  -  -  -  -  -  -  -  -  -  -  -  -  -  -  -  -  -  -  -
\subsubsection{Immediates}
\label{sssec:immediates}
% -  -  -  -  -  -  -  -  -  -  -  -  -  -  -  -  -  -  -  -  -  -  -  -  -

\cref{fig:wirimmediates} from \cref{ssec:immparams} (see \cpageref{fig:wirimmediates}) already revealed that processor-specific classes for immediate parameters are derived from the generic base classes \texttt{WIR\_SignedImmediateParameter} and \texttt{WIR\_UnsignedImmedi\-ateParameter}, resp.
This process is as simple as depicted in \cref{lst:tcconst9signed} for signed 9-bit immediates of the TriCore processor model.
As can be seen from \cref{lst:tcconst9signed2}, an actual value \texttt{v} is simply forwarded to the constructor of the base class, together with the hard-coded width of 9 bits.
This design enables \gls{WIR} to enforce that actual values \texttt{v} are always within the ranges that can be represented with the immediate's bit width.

\begin{lstlisting}[style=code,float=h,caption={Derivation of Processor-Specific Immediates (here for TriCore architecture)},label=lst:tcconst9signed]
class TC_Const9_Signed final : public WIR_SignedImmediateParameter<TC_Const9_Signed>
{
  public:

    // Default constructor for TriCore signed const9 immediates.
    explicit TC_Const9_Signed@( signed long long@ v @)@ :Ü\label{lst:tcconst9signed1}Ü
      WIR_SignedImmediateParameter<TC_Const9_Signed> { v, §9§ }Ü\label{lst:tcconst9signed2}Ü
    {};

    [...]
};
\end{lstlisting}

Furthermore, it allows to add specific checks to an immediate class' constructor constraining the valid value range even further in a processor-specific way.
For example, some TriCore operations require 10-bit unsigned immediates that are multiples of 4.
Thus, the constructor of class \texttt{TC\_Const10\_Unsigned} asserts that the division of value \texttt{v} by 4 produces no remainder.


% -  -  -  -  -  -  -  -  -  -  -  -  -  -  -  -  -  -  -  -  -  -  -  -  -
\subsubsection{Opcodes}
\label{sssec:opcodes}
% -  -  -  -  -  -  -  -  -  -  -  -  -  -  -  -  -  -  -  -  -  -  -  -  -

As already discussed earlier in \cref{sssec:regtypes}, register types are modeled using \gls{WIR} classes resembling enumeration types which can be inherited into actual processor models.
The very same concept is exploited to represent a processor model's opcodes (cf. \cref{lst:opcodes}).

As can be seen, generic opcodes are relized by class \texttt{OpCode} which again is a sub-class of \texttt{WIR\_BaseProcessor}.
And in analogy to register types, this class \texttt{OpCode} is derived from \texttt{WIR\_InheritableEnum} (\cref{lst:opcodes1}).
Objects of this class \texttt{WIR\_BaseProcessor::OpCode} are the very first mandatory argument when \gls{WIR} operations are created, cf. \cref{lst:opconstruction} in \cref{sec:operations}.
An opcode is tightly coupled with a mnemonic string where such strings are used to produce \gls{ISA}-compliant assembly output of \gls{WIR} code.
The mnemonic string associated with an opcode can be retrieved via method \texttt{getName} (\cref{lst:opcodes2}).

\enlargethispage{5mm}

\cref{sec:operations} furthermore already discussed that operations can be classified into different categories, depending on their opcodes (cf. \cref{lst:operclass}).
Thus, an operation's opcode in general determines whether an operation is, e.g., a function call or something else.
As a consequence, the very same Boolean properties as already presented in \cref{lst:operclass} in the context of \gls{WIR} operations can directly be checked for opcodes, cf. \crefrange{lst:opcodes3}{lst:opcodes14} of \cref{lst:opcodes}.
In rare cases, there could be situations where the classification might not only depend on an opcode itself, but maybe on the specific context where and how an opcode occurs in some \gls{WIR} operation.
For this reason, all these methods from \cref{lst:opcodes} take a reference to an operation as argument so that the particular operation could be inspected by these classification methods.
In the methods' default implementation within class \texttt{WIR\_BaseProcessor::OpCode} as shown in \cref{lst:opcodes}, the operations \texttt{o} provided as arguments are simply ignored.
But since all these methods are virtual methods, they can be overloaded in inherited processor models so that a precise classification of some opcode in the context of its operation can be realized.

\begin{lstlisting}[style=code,float=!ht,caption={Generic Opcode Handling of Class \texttt{WIR\_BaseProcessor}},label=lst:opcodes]
class WIR_BaseProcessor : public WIR_SystemComponent
{
  public:

    [...]

    class OpCode : public WIR_InheritableEnumÜ\label{lst:opcodes1}Ü
    {
      public:

        // Get an opcode's specific mnemonic.
        ästd::stringä getName@( void )@ const;Ü\label{lst:opcodes2}Ü

        // Determine whether an opcode reads from memory.
        virtual @bool@ isMemoryLoad@(@ const WIR_Operation æ&æo @)@ const;Ü\label{lst:opcodes3}Ü

        // Determine whether an opcode writes to memory.
        virtual @bool@ isMemoryStore@(@ const WIR_Operation æ&æo @)@ const;Ü\label{lst:opcodes4}Ü

        // Determine whether an opcode implicitly accesses memory, without being an explicit load or store.
        virtual @bool@ isImplicitMemoryAccess@(@ const WIR_Operation æ&æo @)@ const;Ü\label{lst:opcodes5}Ü

        // Determine whether an opcode is a register-register move.
        virtual @bool@ isMove@(@ const WIR_Operation æ&æo @)@ const;Ü\label{lst:opcodes6}Ü

        // Determine whether an opcode calls a function.
        virtual @bool@ isCall@(@ const WIR_Operation æ&æo @)@ const;Ü\label{lst:opcodes7}Ü

        // Determine whether an opcode indirectly calls a function.
        virtual @bool@ isIndirectCall@(@ const WIR_Operation æ&æo @)@ const;Ü\label{lst:opcodes8}Ü

        // Determine whether an opcode returns from a function call.
        virtual @bool@ isReturn@(@ const WIR_Operation æ&æo @)@ const;Ü\label{lst:opcodes9}Ü

        // Determine whether an opcode performs a conditional jump.
        virtual @bool@ isConditionalJump@(@ const WIR_Operation æ&æo @)@ const;Ü\label{lst:opcodes10}Ü

        // Determine whether an opcode performs an unconditional jump.
        virtual @bool@ isUnconditionalJump@(@ const WIR_Operation æ&æo @)@ const;Ü\label{lst:opcodes11}Ü

        // Determine whether an opcode indirectly performs a jump.
        virtual @bool@ isIndirectJump@(@ const WIR_Operation æ&æo @)@ const;Ü\label{lst:opcodes12}Ü

        // Determine whether an opcode is an assembly data directive.
        virtual @bool@ isAsmDataDirective@(@ const WIR_Operation æ&æo @)@ const;Ü\label{lst:opcodes13}Ü

        // Determine whether an opcode has side effects.
        virtual @bool@ hasSideEffects@(@ const WIR_Operation æ&æo @)@ const;Ü\label{lst:opcodes14}Ü

        [...]

      protected:

        // Default constructor creating an opcode with given mnemonic string.
        explicit OpCode@(@ const ästd::stringä æ&æs @)@;Ü\label{lst:opcodes15}Ü

        // Variadic constructor creating an opcode with given mnemonic string and properties.
        templateæ<ætypename... Argsæ>æ OpCode@(@ const ästd::stringä æ&æs, WIR_OpCodeType t, Argsæ&&æ... __args @)@;Ü\label{lst:opcodes16}Ü
    };
};
\end{lstlisting}

\clearpage

All these characteristics of opcodes, their mnemonics and their classification, are set during construction.
\cref{lst:opcodes15,lst:opcodes16} show the declaration of the constructors of class \texttt{OpCode} which both expect the mnemonic string as mandatory argument.
The constructor from \cref{lst:opcodes16} furthermore accepts a variadic number of opcode types so that all necessary properties of opcodes can be specified conveniently.
The enumeration type \texttt{WIR\_OpCodeType} is defined in \cref{lst:opctypes}.

\begin{lstlisting}[style=code,float=t,caption={Classification of Opcodes by Type \texttt{WIR\_OpCodeType}},label=lst:opctypes]
// Enum WIR_OpCodeType represents different types of WIR opcodes.
enum class WIR_OpCodeType : @char@
{
  load,                     // An opcode that reads from memory.
  store,                    // An opcode that writes to memory.
  implicitMemoryAccess,     // An opcode that implicitly accesses memory.
  move,                     // An opcode that is a register-register move.
  call,                     // An opcode that calls a function.
  indirectCall,             // An opcode that indirectly calls a function.
  ret,                      // An opcode that returns from a function call.
  condJump,                 // An opcode that conditionally branches.
  uncondJump,               // An opcode that unconditionally branches.
  indirectJump,             // An opcode that indirectly branches.
  asmDirective,             // An opcode standing for an assembly data directive.
  sideEffect                // An opcode that has side effects.
};
\end{lstlisting}

The concretization of opcodes for some real processor architecture is as easy as already introduced earlier in the context of register types.
An excerpt of the opcode modelling for the running \texttt{TC13} example is shown in \cref{lst:tcopcodes}.
All opcodes simply need to be enumerated in the class declaration (cf. \crefrange{lst:tcopcodes2}{lst:tcopcodes4}).
The opcodes' definitions in \crefrange{lst:tcopcodes5}{lst:tcopcodes7} then construct all the required objects and associate all attributes to them.

This way, rather simple opcodes that do not belong to any specialized category captured in \cref{lst:opctypes} can be specified in a very lightweight manner.
See, e.g., the \texttt{ABS} and \texttt{ADD} opcodes which only require their respective mnemonic strings in \cref{lst:tcopcodes5,lst:tcopcodes6}, resp.
Using opcode types, the \texttt{TC13} opcode \texttt{CALL} is classified as both a function call and a memory access (cf. \cref{lst:tcopcodes7}).

\begin{lstlisting}[style=code,float=t,caption={Concrete Opcode Handling (here for Class \texttt{TC13})},label=lst:tcopcodes]
class TC13 : public WIR_Processor<TC13>
{
  public:

    [...]

    class OpCode : public WIR_Processor::OpCodeÜ\label{lst:tcopcodes1}Ü
    {
      public:

        // Declaration of Absolute Value opcode.
        static const OpCode ABS;Ü\label{lst:tcopcodes2}Ü

        // Declaration of Addition opcode.
        static const OpCode ADD;Ü\label{lst:tcopcodes3}Ü

        // Declaration of Function Call opcode.
        static const OpCode CALL;Ü\label{lst:tcopcodes4}Ü

        [...]
    };

    [...]
};

// Definition of opcodes.
const TC13::OpCode TC13::OpCode::ABS { "abs" };Ü\label{lst:tcopcodes5}Ü
const TC13::OpCode TC13::OpCode::ADD { "add" };Ü\label{lst:tcopcodes6}Ü
const TC13::OpCode TC13::OpCode::CALL { "call",                                                                                                   WIR_OpCodeType::call, WIR_OpCodeType::implicitMemoryAccess };Ü\label{lst:tcopcodes7}Ü
[...]
\end{lstlisting}

All the mechanisms described in this section allow the very simple but yet flexible and powerful specification and classification of opcodes for real-life processor architectures.
This flexbility is demonstrated in the following example that -- in contrast to the running example that bases on the TriCore architecture -- now refers to RISC-V.

\begin{wirExample}%[skin=enhancedfirst]
The RISC-V RV32 \gls{ISA}~\cite{riscv17} does not feature a dedicated opcode for register-register moves.
However, recognizing such operations that copy a value from one register to another is of particular interest during register allocation (cf. \cref{sec:optregalloc}).
Such moves are realized in RISC-V assembly code using operations of the style \texttt{addi \emph{tgtReg}, \emph{srcReg}, 0} where the constant 0 is added to a given source register, and the result of this (dummy) addition is put in the target register.

In order to identify such variants of \texttt{addi} as moves, it would be incorrect to simply classify the opcode \texttt{RV32I::OpCode::ADDI} as move during its construction, since this would hold for all kinds of \texttt{addi}'s, even for those where non-zero constants are added.

Instead, an overloaded method \\
\verb|  |\texttt{\textbf{virtual} \textcolor{blue}{bool} RV32I::OpCode::isMove\textcolor{blue}{(} \textbf{const} WIR\_Operation \textcolor{purple}{\&}o \textcolor{blue}{)}} \\
should be provided which returns \texttt{true} only if the mnemonic of \texttt{o} is \texttt{RV32I::OpCode::ADDI} and if its immediate value is 0.
% \end{wirExample}

% \begin{wirExample}[title=,skin=enhancedlast]
Alternatively, a pseudo-opcode that explicitly models a register-register move could be added to the processor model.
A pseudo-opcode is an opcode that is not directly part of the \gls{ISA} specification, but that is added to the processor model for convenience reasons: \\
\verb|  |\texttt{\textbf{const} RV32I::OpCode RV32I::OpCode::MOV \{ \textcolor{red}{"addi"}, WIR\_OpCodeType::move \}; } \\
This \texttt{MOV} pseudo-opcode is clearly classified as move, and it uses the same mnemonic string ``\texttt{addi}'' as the regular RISC-V opcode \texttt{RV32I::OpCode::ADDI} does.
If application code ensures that each operation using this \texttt{MOV} opcode always has a target register, a source register, and the immediate constant 0 as parameters, then the \gls{WIR} library will always produce \gls{ISA}-compliant RISC-V assembly code and will properly recognize such moves.
\end{wirExample}


% -  -  -  -  -  -  -  -  -  -  -  -  -  -  -  -  -  -  -  -  -  -  -  -  -
\subsubsection{Operation Formats}
\label{sssec:opformats}
% -  -  -  -  -  -  -  -  -  -  -  -  -  -  -  -  -  -  -  -  -  -  -  -  -

\begin{lstlisting}[style=code,float=t,caption={Generic Operation Format Handling of Class \texttt{WIR\_BaseProcessor}},label=lst:opformats]
class WIR_BaseProcessor : public WIR_SystemComponent
{
  public:

    [...]

    class OperationFormat : public WIR_InheritableEnumÜ\label{lst:opformats1}Ü
    {
      public:

        // Get an operation format's bit width.
        @unsigned int@ getBitWidth@( void )@ const;Ü\label{lst:opformats2}Ü

        // Get an operation format's byte size.
        @unsigned int@ getSize@( void )@ const;Ü\label{lst:opformats3}Ü

        // Get an operation format's specific name.
        ästd::stringä getName@( void )@ const;Ü\label{lst:opformats4}Ü

        [...]

      protected:

        // Default constructor creating an operation format with given bit width and name.
        OperationFormat@( unsigned int@ w, const ästd::stringä æ&æn @)@;Ü\label{lst:opformats5}Ü
    };
};
\end{lstlisting}

As discussed in \cref{sec:operations}, operation formats are -- besides opcodes -- the second critical attribute that is required during the creation of \gls{WIR} operations in order to ensure \gls{ISA}-safeness.
Operation formats specify the number and the types of an operation's explicit parameters, as well as their order.
For the enumeration of a processor model's operation format, exactly the same mechanisms as already introduced for opcodes are used (cf. \cref{lst:opformats}).
Here, sub-class \texttt{WIR\_BaseProcessor::OperationFormat} represents such operation formats.
Their most relevant attribute is their bit width, e.g., whether an operation of a specific format is 16 or 32 bits wide, and this information can be retrieved using method \texttt{getBitWidth} (cf. \cref{lst:opformats2}).
\texttt{getSize} (\cref{lst:opformats3}) uses this bit width, divides it by 8 and return this as a format's byte size.
Should a bit width not be a multiple of 8, \texttt{getSize} rounds up accordingly.
For the purpose of \gls{WIR} code inspection, the human-readable name of an operation format can be retrieved via \texttt{getName} (\cref{lst:opformats4}).
The constructor of class \texttt{OperationFormat} takes a format's bit width and its name as argument (\cref{lst:opformats5}).

Declaration and definition of actual operation formats for some concrete processor model -- here again for the \texttt{TC13} example -- follow the same lines as described in earlier sections, cf. \cref{lst:tcopformats}.
The naming scheme for the operation formats used in this code example mostly uses capital letters to denote which kind of parameter is used where by an operation format.
For the \texttt{TC13} example, ``\texttt{A}'' represents any address register parameter, while ``\texttt{D}'' stands for an arbitrary data register parameter.
Thus, operation format ``\texttt{AAA}'' (cf. \cref{lst:tcopformats2}) requires three registers as explicit parameters, all of them need to be address registers.
The comment attached to this format furthermore indicates that the first address register is defined while the following ones are supposed to be used.
Likewise, format ``\texttt{DDC9\_1}'' (\cref{lst:tcopformats3}) expects two data registers as first two parameters.
The ``\texttt{C9}'' here denotes a 9-bit constant immediate as third parameter.
Again, the format's comment furthermore reveals that the first data register is assumed to be written while the second one is read.
The constant immediate of \texttt{DDC9\_1} is assumed to be signed.
Since the TriCore \gls{ISA} features several operations using two data registers and 9-bit constants, but which slightly differ from each other, operation formats might have a suffix like, e.g., ``\texttt{\_1}''.
In \cref{lst:tcopformats4}, the prefix ``\texttt{S}'' denotes a small operation format, i.e., one that is 16 bits wide instead of 32 bits.
Furthermore, this particular operation format requires exactly one physical register as read and written parameter, which has to be the stack pointer.
This is indicated by ``\texttt{SP}'' in \cref{lst:tcopformats4}, and ``\texttt{C8}'' here denotes an unsigned 8-bit constant.

\begin{lstlisting}[style=code,float=t,caption={Concrete Operation Format Handling (here for Class \texttt{TC13})},label=lst:tcopformats]
class TC13 : public WIR_Processor<TC13>
{
  public:

    [...]

    class OperationFormat : public WIR_Processor::OperationFormatÜ\label{lst:tcopformats1}Ü
    {
      public:

        // Declaration of 32-bit operation format aReg (def), aReg (use), aReg (use).
        static const OperationFormat AAA;Ü\label{lst:tcopformats2}Ü

        // Declaration of 32-bit operation format dReg (def), dReg (use), signed const9.
        static const OperationFormat DDC9_1;Ü\label{lst:tcopformats3}Ü

        // Declaration of 16-bit operation format SP (defuse), unsigned const8.
        static const OperationFormat SSPC8;Ü\label{lst:tcopformats4}Ü

        [...]
    };

    [...]
};

// Definition of operation formats.
const TC13::OperationFormat TC13::OperationFormat::AAA { §32§, "AAA" };Ü\label{lst:tcopformats5}Ü
const TC13::OperationFormat TC13::OperationFormat::DDC9_1 { §32§, "DDC9_1" };Ü\label{lst:tcopformats6}Ü
const TC13::OperationFormat TC13::OperationFormat::SSPC8 { §16§, "SSPC8" };Ü\label{lst:tcopformats7}Ü
[...]
\end{lstlisting}

From the definition of these exemplary operation formats in \crefrange{lst:tcopformats5}{lst:tcopformats7} of \cref{lst:tcopformats}, it becomes clear that their bit widths and their names get properly assigned.
What remains to be discussed is how the \gls{WIR} library is able to enforce that, e.g., an operation of format \texttt{DDC9\_1} is always guaranteed to have one defined data register parameter first, a used data register next, finally followed by a signed 9-bit immediate.
Documenting such properties in a comment is nice, but not sufficient at all.

Information of this kind is registered within \gls{WIR} during its initialization.
\cref{chap:build} already mentioned -- without going into any details -- that the \gls{WIR} library must be initialized before any of its data structures can be used in application code.
This initialization is shown in \cref{lst:testprograminit} of \cref{lst:testprogram} (see \cpageref{lst:testprogram}) where the global function \texttt{WIR\_Init} is called.
Main purpose of \texttt{WIR\_Init} is to recursively trigger the initialization of all the \gls{WIR} processor models.
Thus, all processor models come with a static \texttt{init} method that gets invoked by \texttt{WIR\_Init}, below again for the running \texttt{TC13} example:

\begin{lstlisting}[style=code,float=h]
class TC13 : public WIR_Processor<TC13>
{
  public:

    [...]

    // Perform global initialization tasks for the TC13 processor model.
    static @void@ init@( void )@;

    [...]
};
\end{lstlisting}

An excerpt of this method \texttt{TC13::init} showing the registration of the exact shape of the operation formats from \cref{lst:tcopformats} is shown in \cref{lst:tcopfreg}.
In \cref{lst:tcopfreg1,lst:tcopfreg2}, virtual TriCore address and data registers are created first.
\cref{lst:tcopfreg3} creates a physical address register ``\texttt{A10}'' for the purpose of initialization, which is the TriCore's stack pointer.
These registers are afterwards used to specify the exact layout of each operation format.
For example, \crefrange{lst:tcopfreg4}{lst:tcopfreg5} show the registration of TriCore operation format \texttt{AAA}, which gets associated with three register parameters, each of them referring to the virtual address register from \cref{lst:tcopfreg1}.
Furthermore, these register parameters clearly indicate whether they are defined or used (cf. \cref{ssec:regparams}).
This registration now allows the \gls{WIR} library to reject any operation of format \texttt{AAA} that, e.g., has a data register as first parameter, or where an address register at the second position is defined instead of being used, or where only two address registers are provided.

Similarly, lines \crefrange{lst:tcopfreg6}{lst:tcopfreg7} show the precise specification of the TriCore operation format \texttt{DDC9\_1}.
It requires the virtual data register from \cref{lst:tcopfreg2}.
The fact that the TriCore \gls{ISA} makes use of immediates of exactly 9 bits width which are supposed to be signed is represented by the \texttt{TC\_Const9\_Signed} immediate in \cref{lst:tcopfreg7} (see also \cref{lst:tcconst9signed} on \cpageref{lst:tcconst9signed}).
This way, any attempt to create an operation of format \texttt{DDC9\_1} featuring, e.g., a constant that does not fit into the value range of 9 bits in two's complement format leads to a runtime error.

\begin{lstlisting}[style=code,float=t,caption={Operation Format Registration (here for Class \texttt{TC13})},label=lst:tcopfreg]
@void@ TC13::init@( void )@
{
  [...]

  // TC13 operation formats.
  const TC_ARegV *aV = new TC_ARegV;Ü\label{lst:tcopfreg1}Ü
  const TC_DRegV *dV = new TC_DRegV;Ü\label{lst:tcopfreg2}Ü
  const TC_ARegP *SP = new TC_ARegP( "10" );Ü\label{lst:tcopfreg3}Ü

  registerOperationFormat@(@Ü\label{lst:tcopfreg4}Ü
    OperationFormat::AAA,
    { new WIR_RegisterParameter@(@ *aV, WIR_Usage::def @)@, new WIR_RegisterParameter@(@ *aV, WIR_Usage::use @)@,
      new WIR_RegisterParameter@(@ *aV, WIR_Usage::use @)@ } @)@;Ü\label{lst:tcopfreg5}Ü

  registerOperationFormat@(@Ü\label{lst:tcopfreg6}Ü
    OperationFormat::DDC9_1,
    { new WIR_RegisterParameter@(@ *dV, WIR_Usage::def @)@, new WIR_RegisterParameter@(@ *dV, WIR_Usage::use @)@,
      new TC_Const9_Signed@(@ §0§ @)@ } @)@;Ü\label{lst:tcopfreg7}Ü

  registerOperationFormat@(@Ü\label{lst:tcopfreg8}Ü
    OperationFormat::SSPC8,
    { new WIR_RegisterParameter@(@ *SP, WIR_Usage::defuse @)@, new TC_Const8_Unsigned@(@ §0§ @)@ } @)@;Ü\label{lst:tcopfreg9}Ü

  [...]

  // TC13 opcode to operation format mapping.
  registerOpCode@(@ OpCode::ADD_A, OperationFormat::AAA @)@;Ü\label{lst:tcopfreg10}Ü
  registerOpCode@(@ OpCode::ADD, OperationFormat::DDC9_1 @)@;Ü\label{lst:tcopfreg11}Ü
  registerOpCode@(@ OpCode::SUB_A, OperationFormat::SSPC8 @)@;Ü\label{lst:tcopfreg12}Ü

  [...]
};
\end{lstlisting}

Using virtual registers during registration has the effect that an operation format will accept any register as parameter that has the very same type as the provided virtual register.
However, there might be cases where an operation of a certain format requires one specific physical register as parameter and does not accept the whole class of, e.g., address registers.
This can be achieved by registering an operation format with exactly the needed physical register as parameter, as can be seen for \texttt{SSPC8} in \crefrange{lst:tcopfreg8}{lst:tcopfreg9}.
Here, the TriCore's physical stack pointer register is used as parameter that is both defined and used.
This way, \gls{WIR} permits to create operations of this format only if the provided register parameter either is a physical register that is identical to the stack pointer, or if it is a virtual address register that is precolored with the stack pointer.

Now that the registration of operation formats has been introduced, the last issue that remains to achieve \gls{ISA}-safeness is the specification of which opcodes can have which operation formats.
Similar to the registration of operation formats, this mapping of opcodes to formats is realized during the initialization of the \gls{WIR} library, again within the processor-specific \texttt{init} methods.
This is shown in \crefrange{lst:tcopfreg10}{lst:tcopfreg12} of \cref{lst:tcopfreg} for the running \texttt{TC13} example.
As can be seen, the \texttt{init} method simply registers each opcode with all supported formats as specified by an architecture's \gls{ISA}.
This way, any attempt to create an operation with an opcode and some non-registered operation format will be rejected by the \gls{WIR} library with a runtime error.


% -  -  -  -  -  -  -  -  -  -  -  -  -  -  -  -  -  -  -  -  -  -  -  -  -
\subsubsection{Addressing Modes}
\label{sssec:amodes}
% -  -  -  -  -  -  -  -  -  -  -  -  -  -  -  -  -  -  -  -  -  -  -  -  -

As already motivated in \cref{ssec:admodeparams}, memory load or store operations of some \glspl{ISA} might require the specification of some addressing mode, i.e., of the way how the machine computes the effective address to be loaded or stored, given the parameters provided by the operation.
Once more, the specification of such addressing modes is realized within \gls{WIR} using the enumeration class \texttt{WIR\_BaseProcessor::AddressingMode} as shown in \cref{lst:adrmode}.
Besides their human-readable name, addressing modes do not carry any further specific attributes.
The sole purpose of this class is to be able to use them explicitly as parameters of an operation using class \texttt{WIR\_AddressingModeParameter} (cf. \cref{ssec:admodeparams}), and to be able to distinguish between them during code generation, analysis, or optimization.

\begin{lstlisting}[style=code,float=t,caption={Generic Addressing Mode Handling of Class \texttt{WIR\_BaseProcessor}},label=lst:adrmode]
class WIR_BaseProcessor : public WIR_SystemComponent
{
  public:

    [...]

    class AddressingMode : public WIR_InheritableEnumÜ\label{lst:adrmode1}Ü
    {
      public:

        // Get an addressing mode's specific name.
        ästd::stringä getName@( void )@ const;Ü\label{lst:adrmode2}Ü

        [...]

      protected:

        // Default constructor creating an addressing mode with given name.
        explicit AddressingMode@(@ const ästd::stringä æ&æn @)@;Ü\label{lst:adrmode3}Ü
    };
};
\end{lstlisting}

For example, the \texttt{TC13} processor model features the so-called post-increment and pre-increment addressing modes, where a pointer residing in some address register is automatically incremented by some bytes either before or after a memory access.
This is reflected within \gls{WIR} in the straightforward way as depicted in \cref{lst:tcadrmode}, where the two addressing modes \texttt{pre} and \texttt{post} are declared and defined.

\begin{lstlisting}[style=code,float=t,caption={Concrete Addressing Mode Handling (here for Class \texttt{TC13})},label=lst:tcadrmode]
class TC13 : public WIR_Processor<TC13>
{
  public:

    [...]

    class AddressingMode : public WIR_Processor::AddressingModeÜ\label{lst:tcadrmode1}Ü
    {
      public:

        // Declaration of pre-increment addressing.
        static const AddressingMode pre;Ü\label{lst:tcadrmode2}Ü

        // Declaration of post-increment addressing.
        static const AddressingMode post;Ü\label{lst:tcadrmode3}Ü

        [...]
    };

    [...]
};

// Definition of addressing modes.
const TC13::AddressingMode TC13::AddressingMode::pre { "preincr" };Ü\label{lst:tcadrmode4}Ü
const TC13::AddressingMode TC13::AddressingMode::post { "postincr" };Ü\label{lst:tcadrmode5}Ü
[...]
\end{lstlisting}


% -  -  -  -  -  -  -  -  -  -  -  -  -  -  -  -  -  -  -  -  -  -  -  -  -
\subsubsection{Condition Fields}
\label{sssec:cond}
% -  -  -  -  -  -  -  -  -  -  -  -  -  -  -  -  -  -  -  -  -  -  -  -  -

The analogy between condition field parameters (cf. \cref{ssec:condparams}) and sub-class \texttt{WIR\_BaseProcessor::Condi\-tion} is exactly of the same nature as described for addressing modes in the previous section.
Condition fields in the \gls{WIR} hardware model also only carry their name string as their only attribute.
For the sake of brevity, we thus omit any code listings here, because they are identical to \cref{lst:adrmode,lst:tcadrmode} and would not provdide any additional insights.


% -  -  -  -  -  -  -  -  -  -  -  -  -  -  -  -  -  -  -  -  -  -  -  -  -
\subsubsection{Logical Sections}
\label{sssec:sections}
% -  -  -  -  -  -  -  -  -  -  -  -  -  -  -  -  -  -  -  -  -  -  -  -  -

The \gls{WIR} library assumes that binary code executed by some processor model has a structure resembling that of \gls{ELF} files~\cite{elf95}.
Thus, it is logically structured into individual sections that hold, e.g., program code, data, uninitialized data, or constant read-only data.
This section structure is orthogonal to that introduced in the \gls{WIR} software model (cf. \cref{chap:swmodel}), the relationship between the software model and the section structure of the \gls{WIR} hardware model will be explained in detail in \cref{chap:sysmodel}.

Such sections are represented by class \texttt{WIR\_Section}.
The construction of \gls{WIR} sections requires a non-empty string holding the section's name, as well as a reference to a physical memory region (cf. \cref{sec:memoryregions}) to which the section will be attached, see \cref{lst:sctconstr1,lst:sctconstr2} of \cref{lst:sctconstruction}.
As usual, method \texttt{getName} (\cref{lst:sctconstr5}) can be used to retrieve the section name.

\begin{lstlisting}[style=code,float=ht,caption={Constructors and Methods for Name Handling of Class \texttt{WIR\_Section}},label=lst:sctconstruction]
// Default constructor creating a named section attached to some physical memory region.
WIR_Section::WIR_Section@(@ const ästd::stringä æ&æs, const WIR_MemoryRegion æ&ær @)@;Ü\label{lst:sctconstr1}Ü

// Default constructor creating a named section attached to some physical memory region.
WIR_Section::WIR_Section@(@ ästd::stringä æ&&æs, const WIR_MemoryRegion æ&ær @)@;Ü\label{lst:sctconstr2}Ü

// Copy constructor.
WIR_Section::WIR_Section@(@ const WIR_Section æ&æs @)@;Ü\label{lst:sctconstr3}Ü

// Move constructor.
WIR_Section::WIR_Section@(@ WIR_Section æ&&æs @)@;Ü\label{lst:sctconstr4}Ü

// Get the section name.
ästd::stringä WIR_Section::getName@( void )@ const;Ü\label{lst:sctconstr5}Ü
\end{lstlisting}

A \gls{WIR} section is characterized by various attributes which again resemble those of \gls{ELF} file sections.
First, there is a section's external alignment that determines the shape of a section's start and end addresses in memory.
The alignment can be set and retrieved using methods \texttt{setAlignment} and \texttt{getAlignment}, resp. (cf. \cref{lst:sctattribs1,lst:sctattribs2} of \cref{lst:sctattribs}).
Such alignments use the standard data type \texttt{std::size\_t}, and they denote how many least-significant bits of the section's start and end address must be 0.
E.g., setting a section's alignment to 3 means that the least-significant 3 bits of the start and end addresses are 0.
In other words, the start and end addresses must be multiples of 8.
A default external alignment of 10 bits is applied by the constructors of class \texttt{WIR\_Section}.

In a similar fashion, the internal alignment of a section can be specified using methods \texttt{setBlock} and \texttt{getBlock} (\cref{lst:sctattribs3,lst:sctattribs4}).
The internal alignment influences which start and end addresses in memory individual symbols get that are packed into one section.
The constructors of class \texttt{WIR\_Section} apply a default internal alignment of 3 bits.

Next, \texttt{getLength} (\cref{lst:sctattribs5}) returns the byte size of a section which includes the byte sizes of all symbols grouped into that section, plus maybe required additional padding bytes resulting from the section's internal alignment.

As already mentioned above, sections are attached to physical memory regions during construction.
Thus, various sections can be associated with one and the same memory region.
This has the consequence that sections are kind of floating around within their respective region so that the start address of one section usually depends on the sizes of preceding sections within their memory region.
In certain cases, it might, however, be necessary to fix a section at some explicit memory address within its region.
The fact whether such an explicit start address has been provided can be checked using method \texttt{isStartSet} (cf. \cref{lst:sctattribs6}).
Method \texttt{getStart} (\cref{lst:sctattribs7}) then provides this explicitly set start address, or 0 if no explicit start address has been defined.

To be more precise, the \gls{WIR} hardware model distinguishes between \emph{execution time} and \emph{load time} at this point.
Execution time refers to that moment when a processor model actively starts executing some binary code.
In contrast, load time takes place before and denotes that moment when a processor's memories are loaded with the contents from the binary executable code.
A section's explicit start address as returned by method \texttt{getStart} now refers to execution time.
For most sections, this start address at execution time is identical to that at load time.

\begin{lstlisting}[style=code,float=ht,caption={Characteristic Attributes of Class \texttt{WIR\_Section}},label=lst:sctattribs]
// Set a section's external alignment, i.e., the alignment of the section's start/end boundaries.
@void@ WIR_Section::setAlignment@(@ ästd::size_tä a @)@;Ü\label{lst:sctattribs1}Ü

// Get a section's external alignment, i.e., the alignment of the section's start/end boundaries.
ästd::size_tä WIR_Section::getAlignment@( void )@ const;Ü\label{lst:sctattribs2}Ü

// Set a section's internal alignment, i.e., the alignment of contents within a section.
@void@ WIR_Section::setBlock@(@ ästd::size_tä a @)@;Ü\label{lst:sctattribs3}Ü

// Get a section's internal alignment, i.e., the alignment of contents within a section.
ästd::size_tä WIR_Section::getBlock@( void )@ const;Ü\label{lst:sctattribs4}Ü

// Get a section's length in bytes.
ästd::size_tä WIR_Section::getLength@( void )@;Ü\label{lst:sctattribs5}Ü

// Determine whether an explicit start address of a section has been set.
@bool@ isStartSet@( void )@ const;Ü\label{lst:sctattribs6}Ü

// Get a section's explicit start address during execution time, i.e., its Virtual Memory Address (VMA).
WIR_MemoryAddress WIR_Section::getStart@( void )@ const;Ü\label{lst:sctattribs7}Ü

// Determine whether the load address of a section has been set previously.
@bool@ isAtSet@( void )@ const;Ü\label{lst:sctattribs8}Ü

// Get a section's load address, i.e., its Loaded Memory Address (LMA).
WIR_MemoryAddress WIR_Section::getAt@( void )@ const;Ü\label{lst:sctattribs9}Ü

// Get the memory region to which a section is attached.
WIR_MemoryRegion æ&æWIR_Section::getRegion@( void )@ const;Ü\label{lst:sctattribs10}Ü

// Get a section's load memory region.
WIR_MemoryRegion æ&æWIR_Section::getLoadRegion@( void )@ const;Ü\label{lst:sctattribs11}Ü

// Determine whether a load memory region has been set previously.
WIR_MemoryAddress WIR_Section::isLoadRegionSet@( void )@ const;Ü\label{lst:sctattribs12}Ü
\end{lstlisting}

An example where this is not the case are data sections mapped to some \gls{ROM}.
Since data is usually both read and written, such particular sections need to be copied from \gls{ROM} to some \gls{RAM} prior to execution time.
Such sections thus have different start addresses at load and at execution time.
For this reason, methods \texttt{isAtSet} (\cref{lst:sctattribs8}) and \texttt{getAt} (\cref{lst:sctattribs9}) can be used to determine whether a start address at load time has been specified and which address that is, resp.
Referring once more to the \gls{ELF} standard, the start address at execution time provided by \texttt{getStart} denotes the so-called \gls{VMA}, while the address returned by \texttt{getAt} stands for the \gls{LMA}.

Memory addresses are represented by type \texttt{WIR\_MemoryAddress}\label{page:memaddr} within \gls{WIR}, which in turn simply is defined as \texttt{WIR\_SaturatingInt<std::uint32\_t>}.
The \gls{WIR}-internal template class \texttt{WIR\_SaturatingInt} is a useful abstraction for memory addresses in that sense that overflows in address computations shall be avoided.
Instead, all computations applied to \gls{WIR} addresses provide the maximum or minimum value of a saturating integer's value range in the case of over- or underflows, resp.
Class \texttt{WIR\_SaturatingInt} basically provides all standard arithmetical, shift, logical and comparison operators with inherent saturation, its implementation can be found in file \texttt{analyses/generic/wirsaturatingint.h}.

Finally, method \texttt{getRegion} (cf. \cref{lst:sctattribs10}) can be used to get the physical memory region to which a section is attached.
If a section gets loaded from one position to another during boot time as described above, method \texttt{getLoadRegion} (\cref{lst:sctattribs11}) can be used to retrieve that memory region from which the section is loaded.
However, \texttt{getLoadRegion} fails with a runtime error if no such load section has been set before.
In order to determine whether a load region was specified or not, method \texttt{isLoadRegionSet} (\cref{lst:sctattribs10}) should be used.

Sections are finally owned and managed by \gls{WIR} processors.
\gls{WIR} mandates and enforces that each processor model always includes at least the following 4 default sections reserved for the indicated contents:
\begin{description}
  \item[\texttt{.text}:] Executable code
  \item[\texttt{.data}:] Initialized data with read/write access rights
  \item[\texttt{.bss}:] Uninitialized data with read/write access rights
  \item[\texttt{.rodata}:] Constant initialized data with read access rights only
\end{description}

In addition to the \gls{API} already shown in \cref{lst:procconstruction} (cf. \cpageref{lst:procconstruction}), class \texttt{WIR\_BaseProcessor} provides the methods depicted in \cref{lst:procsct} for the handling of \gls{WIR} sections.
A reference to the set of all sections managed by a processor model is returned by \texttt{getSections} (cf. \cref{lst:procsct1}), where the return type \texttt{WIR\_SectionSet} is basically declared as \texttt{std::set<std::reference\_wrapper<WIR\_Section>{}>} within \gls{WIR} (cf. \cref{chap:setsmaps}).
The abovementioned 4 default sections can conveniently be accessed using methods \texttt{getTextSection} to \texttt{getRODataSection} from \crefrange{lst:procsct2}{lst:procsct5}.

\begin{lstlisting}[style=code,float=t,caption={Methods for Section Handling of Class \texttt{WIR\_BaseProcessor}},label=lst:procsct]
// Get the set of all sections of a processor model.
const WIR_SectionSet æ&æWIR_BaseProcessor::getSections@( void )@ const;Ü\label{lst:procsct1}Ü

// Get a processor's default ".text" section.
WIR_Section æ&æWIR_BaseProcessor::getTextSection@( void )@ const;Ü\label{lst:procsct2}Ü

// Get a processor's default ".data" section.
WIR_Section æ&æWIR_BaseProcessor::getDataSection@( void )@ const;Ü\label{lst:procsct3}Ü

// Get a processor's default ".bss" section.
WIR_Section æ&æWIR_BaseProcessor::getBssSection@( void )@ const;Ü\label{lst:procsct4}Ü

// Get a processor's default ".rodata" section.
WIR_Section æ&æWIR_BaseProcessor::getRODataSection@( void )@ const;Ü\label{lst:procsct5}Ü

// Determine whether a processor contains a section with the given ID.
@bool@ WIR_BaseProcessor::containsSection@(@ WIR_id_t id @)@ const;Ü\label{lst:procsct6}Ü

// Determine whether a processor contains a section with the given name.
@bool@ WIR_BaseProcessor::containsSection@(@ const ästd::stringä æ&æn @)@ const;Ü\label{lst:procsct7}Ü

// Determine whether a processor contains a section with the given VMA.
@bool@ WIR_BaseProcessor::containsSectionStart@(@ const WIR_MemoryAddress æ&æa @)@ const;Ü\label{lst:procsct8}Ü

// Determine whether a processor contains a section with the given LMA.
@bool@ WIR_BaseProcessor::containsSectionAt@(@ const WIR_MemoryAddress æ&æa @)@ const;Ü\label{lst:procsct9}Ü

// Find the section with the given ID.
WIR_SectionSetæ::const_iteratoræ WIR_BaseProcessor::findSection@(@ WIR_id_t id @)@ const;Ü\label{lst:procsct10}Ü

// Find the section with the given name.
WIR_SectionSetæ::const_iteratoræ WIR_BaseProcessor::findSection@(@ ästd::stringä æ&æn @)@ const;Ü\label{lst:procsct11}Ü

// Find the section with the given VMA.
WIR_SectionSetæ::const_iteratoræ WIR_BaseProcessor::findSectionByStart@(@ const WIR_MemoryAddress æ&æa @)@ const;Ü\label{lst:procsct12}Ü

// Find the section with the given LMA.
WIR_SectionSetæ::const_iteratoræ WIR_BaseProcessor::findSectionByAt@(@ const WIR_MemoryAddress æ&æa @)@ const;Ü\label{lst:procsct13}Ü
\end{lstlisting}

Whether a certain section exists within a processor model can be queried via the various \texttt{contains} methods.
They allow to check the presence of a section with a given numerical ID (\cref{lst:procsct6}), with a given name string (\cref{lst:procsct7}), or with given \glspl{VMA} or \glspl{LMA}, resp. (\cref{lst:procsct8,lst:procsct9}).
The equivalent \texttt{find} methods (cf. \crefrange{lst:procsct10}{lst:procsct13}) locate a section within a processor's section set and provide an iterator to the found position.
If the set does not contain the section to be found, the \texttt{end()} iterator is returned.


% - - - - - - - - - - - - - - - - - - - - - - - - - - - - - - - - - - - -
\subsection{MIPS SPIM Processor Model}
\label{ssec:mipsmodel}
% - - - - - - - - - - - - - - - - - - - - - - - - - - - - - - - - - - - -

Within the \gls{WIR} library, the MIPS SPIM processor model just serves as a demonstrator for its various design concepts, and not for actual code generation or optimization.
For the sake of simplicity, this processor model thus only focuses on the MIPS Core Instruction Set~\cite{mips07} and the integer arithmetic instructions.
Class \texttt{MIPS} from file \texttt{WIR/arch/generic/mips.h} is the central element of this processor model, and -- as depicted in \cref{fig:wirprocessors} -- is inherited from base class \texttt{WIR\_Processor}.

The modeled \texttt{MIPS} architecture features a total of 32 general-purpose integer registers \texttt{\$0} -- \texttt{\$31}, each of them being 32 bits wide.
Thus, this processor model declares only one register type \texttt{MIPS::RegisterType::reg}, defined in file \texttt{WIR/arch/generic/mipsregistertypes.cc}.
The key characteristics of this register type and of the resulting physical and virtual \texttt{MIPS} registers are summarized in \cref{tab:mipsregtypes}.
As can be seen, the modeled registers are not hierarchical.
The physical registers are represented by class \texttt{MIPS\_RegP} (file \texttt{mipsregphysical.h}), and class \texttt{MIPS} provides the methods \texttt{r0()} to \texttt{r31()} to access them.
The physical register returned by \texttt{r29()} is marked as the architecture's stack pointer.
The model's virtual registers are defined by class \texttt{MIPS\_RegV} (file \texttt{mipsregvirtual.h}), and their names always begin with the prefix ``\texttt{\$r\_}'' in order to clearly distinguish them from the physical MIPS registers.

\begin{table}[h]
  \begin{center}
    \caption{Register Types, Classes and Methods of \texttt{MIPS} Processor Model}
    \label{tab:mipsregtypes}
    \begin{tabular}{|c|c|c|l|l|l|l|}
      \hline
      \multirow{2}{*}[-1mm]{\textbf{\shortstack[c]{Register\\ Type}}} & \multirow{2}{*}[-1mm]{\textbf{\shortstack[c]{Bit\\ Width}}} & \multirow{2}{*}[-1mm]{\textbf{\shortstack[c]{Has\\ Childs}}} & \multicolumn{2}{c|}{\textbf{Physical Registers}} & \multicolumn{2}{c|}{\textbf{Virtual Registers}} \\
      & & & Class & Prefix & Class & Prefix \\
      \hhline{|=|=|=|=|=|=|=|}
      \texttt{reg} & 32 & & \texttt{MIPS\_RegP} & ``\texttt{\$}'' & \texttt{MIPS\_RegV} & ``\texttt{\$r\_}'' \\
                   &    & & \multicolumn{2}{c|}{\texttt{r0()} -- \texttt{r31()}} & & \\
      \hline
    \end{tabular}
  \end{center}
\end{table}

Code examples how to create virtual \texttt{MIPS} registers and how to check their various properties can be found in the unit tests provided in directory \texttt{WIR/arch/generic/tests}, more specifically in test programs \texttt{mipsVirtualReg\-ister1.cc} and \texttt{genericmips.cc}.

The classes of immediate parameters of the \texttt{MIPS} model are summarized in \cref{tab:mipsimmtypes}.
Due to the architecture's simplicity, basically only 16-bit immediates are used by MIPS operations in their signed and unsigned variants.
The MIPS' shift operations allow to shift a register's contents by up to 31 bits so that they use class \texttt{MIPS\_Immediate5\_Shamt} as immediates specifying such shift amounts.

\begin{table}[h]
  \begin{center}
    \caption{Classes for Immediates of \texttt{MIPS} Processor Model}
    \label{tab:mipsimmtypes}
    \begin{tabular}{|c|c|l|l|p{42mm}|}
      \hline
      \multirow{2}{*}[-1mm]{\textbf{\shortstack[c]{Immediate\\ Type}}} & \multirow{2}{*}[-1mm]{\textbf{\shortstack[c]{Bit\\ Width}}} & \multirow{2}{*}{\textbf{Class}} & \multirow{2}{*}{\textbf{Value Range}} & \multirow{2}{*}{\textbf{Comment}} \\
      & & & & \\
      \hhline{|=|=|=|=|=|}
      Signed 16-bit & 16 & \texttt{MIPS\_Immediate16\_Signed} & $[-2^{15}$, $+2^{15} - 1]$ & \\
      \hline
      Unsigned 16-bit & 16 & \texttt{MIPS\_Immediate16\_Unsigned} & $[0$, $+2^{16} - 1]$ & \\
      \hline
      Unsigned 5-bit & 5 & \texttt{MIPS\_Immediate5\_Shamt} & $[0$, $+31]$ & Shift amounts for 32-bit architecture \\
      \hline
    \end{tabular}
  \end{center}
\end{table}

All \texttt{MIPS} unit tests in \texttt{WIR/arch/generic/tests/mipsImmediate*} feature concrete code examples how to create and use immediates of the classes shown in \cref{tab:mipsimmtypes}.
These unit tests also cover scenarios where values that are beyond an immediate's valid value range are used.
These tests demonstrate the behavior of the \gls{WIR} library in the presence of ill-formed immediates.

Since the \texttt{MIPS} processor model excludes all floating-point operations, the opcodes modeled by class \texttt{MIPS::Op\-Code} consist of the 45 remaining ones according to~\cite{mips07}, beginning with \texttt{ADD} and ending with \texttt{XORI} (cf. file \texttt{mipsopcodes.cc}).
None of the modeled opcodes is a pseudo-opcode, and their classification into opcode types is shown in \cref{tab:mipsopctypes}.
Opcodes not occurring in the table are regular integer arithmetic operations not carrying any specific type within \gls{WIR}.

\begin{table}[t]
  \begin{center}
    \caption{Opcode Types of \texttt{MIPS} Processor Model}
    \label{tab:mipsopctypes}
    \begin{tabular}{|c|p{95mm}|p{42mm}|}
      \hline
      \multirow{2}{*}[-1mm]{\textbf{\shortstack[c]{Opcode\\ Type}}} & \multirow{2}{*}{\textbf{Opcodes}} & \multirow{2}{*}{\textbf{Comment}} \\
      & & \\
      \hhline{|=|=|=|}
      \texttt{load} & \texttt{LB}, \texttt{LBU}, \texttt{LH}, \texttt{LHU}, \texttt{LW} & \\
      \hline
      \texttt{store} & \texttt{SB}, \texttt{SH}, \texttt{SW} & \\
      \hline
      \texttt{call} & \texttt{JAL}, \texttt{JALR} & \\
      \hline
      \texttt{indirectCall} & \texttt{JALR} & \\
      \hline
      \texttt{ret} & \texttt{JR} & \\
      \hline
      \texttt{condJump} & \texttt{BEQ}, \texttt{BNE} & \\
      \hline
      \texttt{uncondJump} & \texttt{J}, \texttt{JR} & \\
      \hline
      \texttt{indirectJump} & \texttt{JR} & \\
      \hline
    \end{tabular}
  \end{center}
\end{table}

Operation formats are defined in file \texttt{mipsoperationformats.cc}, and they are all 32 bits wide.
The naming scheme used for the \texttt{MIPS} operation formats distinguishes between signed and unsigned immediates, labels, registers, and shift amounts, see also \cref{tab:mipsopf}.
Using these specifiers, a total of 14 operation formats is defined, ranging from \texttt{MIPS::OperationFormat::L} (one used label only) to \texttt{MIPS::OperationFormat::RRS} (one defined and one used register, one shift amount).

\begin{table}[h]
  \begin{center}
    \caption{Naming Scheme for Operation Formats of \texttt{MIPS} Processor Model}
    \label{tab:mipsopf}
    \begin{tabular}{|c|p{136mm}|}
      \hline
      \multirow{2}{*}[-1mm]{\textbf{\shortstack[c]{Operation Format\\ Specifier}}} & \multirow{2}{*}{\textbf{Parameter Type}} \\
      & \\
      \hhline{|=|=|}
      ``\texttt{I}'' & Signed immediate \texttt{MIPS\_Immediate16\_Signed} \\
      \hline
      ``\texttt{IU}'' & Unsigned immediate \texttt{MIPS\_Immediate16\_Unsigned} \\
      \hline
      ``\texttt{L}'' & Label parameter \\
      \hline
      ``\texttt{R}'' & Register parameter of type \texttt{MIPS::RegisterType::reg} \\
      \hline
      ``\texttt{S}'' & Shift amount \\
      \hline
    \end{tabular}
  \end{center}
\end{table}

As already described in \cref{sssec:opformats}, the static method \texttt{MIPS::init} serves to register the exact layout of all operation formats and maps all MIPS opcodes to their respective operation formats.
Code examples showing the creation and use of all modeled \texttt{MIPS} opcodes in all supported operation formats can be found in unit test \texttt{WIR/arch/generic/tests/mipsOperations1.cc}.


% - - - - - - - - - - - - - - - - - - - - - - - - - - - - - - - - - - - -
\subsection{TriCore Processor Models}
\label{ssec:tcmodel}
% - - - - - - - - - - - - - - - - - - - - - - - - - - - - - - - - - - - -

Many of the fundamental concepts of \gls{WIR} processor models have been introduced in \cref{ssec:procmodel} with the TriCore architecture as running example.
We thus continue next with the description of this family of processor models.
More precisely, \gls{WIR} contains two TriCore models, one for \gls{ISA} version 1.3, and one for version 1.3.1~\cite{tricore08}.
As shown in \cref{fig:tricore}, the model \texttt{TC13} is inherited from base class \texttt{WIR\_Processor}, while \texttt{TC131} in turn inherits from \texttt{TC13}.
This shows how the generic concepts introduced in \cref{ssec:procmodel} enable the elegant and simple modelling of entire families of processors that all build upon each other.

\begin{figure}[t]
  \begin{center}
    \includegraphics[width=0.89\textwidth]{tricore.png}
  \end{center}
  \caption{Simplified Inheritance Diagram of TriCore Processor Models}
  \label{fig:tricore}
\end{figure}

As already mentioned in \cref{sssec:regtypes}, the TriCore architecture distinguishes between address registers which are reserved for pointer/address arithmetics and memory accesses on the one hand, and data registers that hold general integer or floating-point data on the other hand.
Thus, the \texttt{TC13} model of \gls{WIR} from file \texttt{WIR/arch/tri\-core/tc13.h} defines the register types \texttt{aReg} and \texttt{dReg} that are both 32 bits wide (cf. \cref{tab:tcregtypes}).
The classes realizing physical and virtual registers of these types are \texttt{TC\_ARegP} and \texttt{TC\_ARegV} (files \texttt{tcaregphysical.h} and \texttt{tcaregvirtual.h}) and \texttt{TC\_DRegP} and \texttt{TC\_DRegV} (files \texttt{tcdregphysical.h} and \texttt{tcdregvirtual.h}), resp.
For the simple access to these physical address and data registers, class \texttt{TC13} comes with methods \texttt{A0()} -- \texttt{A15()} and \texttt{D0()} -- \texttt{D15()}.
Physical address register \texttt{A10()} is the model's stack pointer.

\begin{table}[h]
  \begin{center}
    \caption{Register Types, Classes and Methods of \texttt{TC13} Processor Model}
    \label{tab:tcregtypes}
    \begin{tabular}{|c|c|c|l|l|l|l|}
      \hline
      \multirow{2}{*}[-1mm]{\textbf{\shortstack[c]{Register\\ Type}}} & \multirow{2}{*}[-1mm]{\textbf{\shortstack[c]{Bit\\ Width}}} & \multirow{2}{*}[-1mm]{\textbf{\shortstack[c]{Has\\ Childs}}} & \multicolumn{2}{c|}{\textbf{Physical Registers}} & \multicolumn{2}{c|}{\textbf{Virtual Registers}} \\
      & & & Class & Prefix & Class & Prefix \\
      \hhline{|=|=|=|=|=|=|=|}
      \texttt{aReg} & 32 & & \texttt{TC\_ARegP} & ``\texttt{a}'' & \texttt{TC\_ARegV} & ``\texttt{a\_}'' \\
                   &     & & \multicolumn{2}{c|}{\texttt{A0()} -- \texttt{A15()}} & & \\
      \hline
      \texttt{dReg} & 32 & & \texttt{TC\_DRegP} & ``\texttt{d}'' & \texttt{TC\_DRegV} & ``\texttt{d\_}'' \\
                   &     & & \multicolumn{2}{c|}{\texttt{D0()} -- \texttt{D15()}} & & \\
      \hline
      \texttt{eReg} & 64 & \checkmark    & \texttt{TC\_ERegP} & ``\texttt{e}'' & \texttt{TC\_ERegV} & ``\texttt{e\_}'' \\
                    &    & \texttt{dReg} & \multicolumn{2}{c|}{\texttt{E0()}, \texttt{E2()}, \dots, \texttt{E14()}} & & \\
      \hline
      \texttt{pReg} & 64 & \checkmark    & \texttt{TC\_PRegP} & ``\texttt{p}'' & \texttt{TC\_PRegV} & ``\texttt{p\_}'' \\
                    &    & \texttt{aReg} & \multicolumn{2}{c|}{\texttt{P0()}, \texttt{P2()}, \dots, \texttt{P14()}} & & \\
      \hline
      \texttt{pswBit} & 1 & & \texttt{TC\_PSWBit} & ``\texttt{PSW.}'' & & \\
                      &   & & \multicolumn{2}{c|}{\texttt{PSW\_C()}} & & \\
      \hline
    \end{tabular}
  \end{center}
\end{table}

According to the TriCore \gls{ISA}, neighboring pairs of such 32-bit address or data registers are hierarchically grouped to form so-called 64-bit extended registers.
For example, data registers \texttt{D4} and \texttt{D5} form the extended data register \texttt{E4}.
These extended registers serve to store 64-bit integer or double-precision floating-point data, or double addresses.
For these reasons, the \texttt{TC13} model introduces the hierarchical register types \texttt{eReg} for extended data registers and \texttt{pReg} for extended address registers, resp.
They all have exactly two child data or address registers.
Their implementations come with classes \texttt{TC\_ERegP} and \texttt{TC\_ERegV} for extended data registers (files \texttt{tceregphysical.h} and \texttt{tceregvirtual.h}), and with classes \texttt{TC\_PRegP} and \texttt{TC\_PRegV} for extended address registers (files \texttt{tcpregphysical.h} and \texttt{tcphregvirtual.h}), resp.
Again, convenient access methods \texttt{E0()} -- \texttt{E14()} and \texttt{P0()} -- \texttt{P14()} exist.
Data flow through such hierarchies of registers is fully automatically considered by the \gls{WIR} library.
This means that any modification of, e.g., data register \texttt{D5} automatically also leads to a re-definition of extended register \texttt{E4}.
In the other sense of direction, writing to \texttt{E4} is equivalent to a definition of both \texttt{D4} and \texttt{D5}.

The last \texttt{TC13} register type models the Carry bit of the TriCore's \gls{PSW} register.
Since the TriCore \gls{ISA} contains a few operations that implicitly define or use this internal Carry bit, this very special register type has been introduced for the proper modelling of data flow via this status bit.
There only exists class \texttt{TC\_PSWBit} as implementation of the physical \gls{PSW} bit, but no virtual register counterpart.
Class \texttt{TC13} offers the method \texttt{PSW\_C()} to access this physical register.

\begin{lstlisting}[style=code,float=t,caption={Testing of Physical Register Properties of Class \texttt{TC13}},label=lst:tcregtests]
// Determine whether a register is the TriCore's stack pointer A10.
static @bool@ TC13::isSP@(@ const WIR_BaseRegister æ&ær @)@;Ü\label{lst:tcregtests1}Ü

// Determine whether a register is the TriCore's return address register A11.
static @bool@ TC13::isRA@(@ const WIR_BaseRegister æ&ær @)@;Ü\label{lst:tcregtests2}Ü

// Determine whether a register is the TriCore's implicit address register A15.
static @bool@ TC13::isA15@(@ const WIR_BaseRegister æ&ær @)@;Ü\label{lst:tcregtests3}Ü

// Determine whether a register is the TriCore's implicit data register D15.
static @bool@ TC13::isD15@(@ const WIR_BaseRegister æ&ær @)@;Ü\label{lst:tcregtests4}Ü

// Determine whether a register is the TriCore's PSW carry bit.
static @bool@ TC13::isPSW_C@(@ const WIR_BaseRegister æ&ær @)@;Ü\label{lst:tcregtests5}Ü

// Determine whether a register is part of the TriCore's lower register context.
static @bool@ TC13::isLCReg@(@ const WIR_BaseRegister æ&ær @)@;Ü\label{lst:tcregtests6}Ü
\end{lstlisting}

Within this structure of the TriCore register file, certain registers play special roles.
The stack pointer \texttt{A10} has already been mentioned above, and address register \texttt{A11} is the TriCore's return address register.
Furthermore, registers \texttt{A15} and \texttt{D15} are the machine's implicit address and data registers, resp.
In order to easily determine whether some arbitrary register \texttt{r} is one of these special registers, class \texttt{TC13} offers the methods shown in \crefrange{lst:tcregtests1}{lst:tcregtests5} of \cref{lst:tcregtests}.
They all return \texttt{true} if the given register \texttt{r} either is exactly the physical special-purpose TriCore register, or if it is virtual but precolored with the according physical register.
Last but not least, the TriCore's register file is structured into an upper and a lower context, where the upper context registers automatically get saved during function calls.
Lower context registers\footnote{TriCore Lower Context: \texttt{A2}, \texttt{A3}, \texttt{A4}, \texttt{A5}, \texttt{A6}, \texttt{A7}; \texttt{D0}, \texttt{D1}, \texttt{D2}, \texttt{D3}, \texttt{D4}, \texttt{D5}, \texttt{D6}, \texttt{D7}; \texttt{E0}, \texttt{E2}, \texttt{E4}, \texttt{E6}; \texttt{P2}, \texttt{P4}, \texttt{P6}.} require special handling during function calls so that method \texttt{TC13::isLCReg} (cf. \cref{lst:tcregtests6}) can be used to determine whether some register \texttt{r} is from the lower context.

Since TriCore \gls{ISA} version 1.3.1 does not add more registers to the architecture, \gls{WIR} model \texttt{TC131} does not need to declare and implement additional register types and classes.
Instead, it simply inherits all the above concepts from \texttt{TC13} without any modification or extension.

The unit tests in files \texttt{WIR/arch/tricore/tests/tc*Reg*} contain various examples showing how to create and use both physical and virtual registers of the various TriCore register types.

\begin{table}[t]
  \begin{center}
    \caption{Classes for Immediates of \texttt{TC13} Processor Model}
    \label{tab:tcimmtypes}
    \begin{tabular}{|c|c|l|l|p{52mm}|}
      \hline
      \multirow{2}{*}[-1mm]{\textbf{\shortstack[c]{Immediate\\ Type}}} & \multirow{2}{*}[-1mm]{\textbf{\shortstack[c]{Bit\\ Width}}} & \multirow{2}{*}{\textbf{Class}} & \multirow{2}{*}{\textbf{Value Range}} & \multirow{2}{*}{\textbf{Comment}} \\
      & & & & \\
      \hhline{|=|=|=|=|=|}
      Unsigned 18-bit & 32 & \texttt{TC\_Const18\_Unsigned} & $[0$, $+0x\textrm{F0003FFF}]$ & Absolute addresses \\
      \hline
      Signed 16-bit & 16 & \texttt{TC\_Const16\_Signed} & $[-2^{15}$, $+2^{15} - 1]$ & \\
      \hline
      Unsigned 16-bit & 16 & \texttt{TC\_Const16\_Unsigned} & $[0$, $+2^{16} - 1]$ & \\
      \hline
      Signed 10-bit & 10 & \texttt{TC\_Const10\_Signed} & $[-512$, $+511]$ & \\
      \hline
      Unsigned 10-bit & 10 & \texttt{TC\_Const10\_Unsigned} & $[0$, $+1023]$ & Stack pointer-relative offsets,\newline must be multiples of 4 \\
      \hline
      Signed 9-bit & 9 & \texttt{TC\_Const9\_Signed} & $[-256$, $+255]$ & \\
      \hline
      Unsigned 9-bit & 9 & \texttt{TC\_Const9\_Unsigned} & $[0$, $+511]$ & \\
      \hline
      Unsigned 8-bit & 8 & \texttt{TC\_Const8\_Unsigned} & $[0$, $+255]$ & \\
      \hline
      Unsigned 5-bit & 5 & \texttt{TC\_Const5\_Unsigned} & $[0$, $+31]$ & Word-level bit indices and shift amounts \\
      \hline
      Signed 4-bit & 4 & \texttt{TC\_Const4\_Signed} & $[-8$, $+7]$ & \\
      \hline
      Unsigned 4-bit & 4 & \texttt{TC\_Const4\_Unsigned} & $[0$, $+15]$ & \\
      \hline
      Unsigned 3-bit & 3 & \texttt{TC\_Const3\_Unsigned} & $[0$, $+7]$ & Byte-level bit indices \\
      \hline
      Unsigned 2-bit & 2 & \texttt{TC\_Const2\_Unsigned} & $[0$, $+3]$ & Address factor of 1, 2, 4 or 8, encoded as logarithm of base 2 \\
      \hline
      Unsigned 1-bit & 1 & \texttt{TC\_Const1\_Unsigned} & $[0$, $+1]$ & Handling of corner cases by Q-format arithmetic operations \\
      \hline
    \end{tabular}
  \end{center}
\end{table}

In sharp contrast to the extremely simple \gls{ISA} design of the MIPS SPIM architecture, the TriCore designers seemed to enjoy the motto ``the more, the better''.
The TriCore \gls{ISA} is characterized by features stemming from various different domains of external computer architecture, amongst others \gls{DSP}, \gls{SIMD} computing, or network processing.
Due to this feature richness of the instruction set, a quite large number of different immediate parameters needed to be modeled within \gls{WIR}, cf. \cref{tab:tcimmtypes}.
Most of them are the straightforward realizations of signed or unsigned immediates of their given bit widths so that we omit their discussion here for the sake of brevity.
Some others are, however, rather special and need a few words of explanation.

Most notably, class \texttt{TC\_Const18\_Unsigned} represents absolute addresses in the context of memory access operations.
According to the TriCore \gls{ISA}, such absolute addresses are encoded in machine operations using 18 bits.
However, the machine's interpretation of these 18 bits differs fundamentally from that of classical unsigned binary numbers, requiring special provisions in class \texttt{TC\_Const18\_Unsigned}.
The TriCore architecture expands such 18-bit absolute addresses to full 32-bit addresses by splitting its 4 most-significant bits from the 14 least-significant bits, and putting a series of 14 zero bits in between:
\begin{displaymath}
  \textit{addr}_{32} := \texttt{const18}[17:14]\ \&\ \texttt{"00\,0000\,0000\,0000"}\ \&\ \texttt{const18}[13:0]
\end{displaymath}
where the $\&$ operator stands for the concatenation of bit strings.
The value range that can be represented with such 18-bit absolute addresses is thus completely different from that of unsigned numbers of 18 bits width (see \cref{tab:tcimmtypes}).
For this reason, class \texttt{TC\_Const18\_Unsigned} internally uses the full width of 32 bits in order to properly represent absolute addresses having the above scheme.
Any arbitrary unsigned 32-bit values can be used to construct immediates of type \texttt{TC\_Const18\_Unsigned}, while the class' constructor rejects those values whose inner bits at positions 14 -- 27 are non-zero.

\texttt{TC\_Const10\_Unsigned} immediates are used as offsets in special variants of load or store operations exploiting stack pointer-relative addressing.
Since the stack pointer has to be aligned to 4 bytes, constant offsets to be added to the stack pointer have to adhere to the same rule.
According to the TriCore \gls{ISA}, these special load and store operations use an unsigned 8-bit immediate that is internally multiplied by 4.
The \texttt{TC13} processor model of \gls{WIR} now uses 10-bit unsigned immediates for this purpose, such that the constructor of class \texttt{TC\_Const10\_Unsigned} issues a runtime error whenever actual values are used that are no even multiples of 4.

\begin{table}[t]
  \begin{center}
    \caption{Opcode Types of \texttt{TC13} Processor Model}
    \label{tab:tc13opctypes}
    \begin{tabular}{|c|p{79mm}|p{42mm}|}
      \hline
      \multirow{2}{*}[-1mm]{\textbf{\shortstack[c]{Opcode\\ Type}}} & \multirow{2}{*}{\textbf{Opcodes}} & \multirow{2}{*}{\textbf{Comment}} \\
      & & \\
      \hhline{|=|=|=|}
      \texttt{load} & \texttt{LD\_A}, \texttt{LD\_B}, \texttt{LD\_BU}, \texttt{LD\_D}, \texttt{LD\_DA}, \texttt{LD\_H}, \texttt{LD\_HU}, \texttt{LD\_Q}, \texttt{LD\_W}, \texttt{LDLCX}, \texttt{LDMST}, \texttt{LDUCX}, \texttt{RSLCX}, \texttt{ST\_T}, \texttt{SWAP\_W} & Standard load operations, context handling, load-\newline modify-store operations \\
      \hline
      \texttt{store} & \texttt{LDMST}, \texttt{ST\_A}, \texttt{ST\_B}, \texttt{ST\_D}, \texttt{ST\_DA}, \texttt{ST\_H}, \texttt{ST\_Q}, \texttt{ST\_T}, \texttt{ST\_W}, \texttt{STLCX}, \texttt{STUCX}, \texttt{SVLCX}, \texttt{SWAP\_W} & Standard store operations, context handling, load-\newline modify-store operations \\
      \hline
      \texttt{implicitMemoryAccess} & \texttt{BISR}, \texttt{CACHEA\_I}, \texttt{CACHEA\_W}, \texttt{CACHEA\_WI}, \texttt{CALL}, \texttt{CALLA}, \texttt{CALLI}, \texttt{RET}, \texttt{RFE}, \texttt{RFM} & \\
      \hline
      \texttt{move} & \texttt{MOV\_RR}, \texttt{MOV\_AA} & Pseudo-opcode \texttt{MOV\_RR} \\
      \hline
      \texttt{call} & \texttt{CALL}, \texttt{CALLA}, \texttt{JL}, \texttt{JLA} & \\
      \hline
      \texttt{indirectCall} & \texttt{CALLI}, \texttt{JLI} & \\
      \hline
      \texttt{ret} & \texttt{RET}, \texttt{RFE}, \texttt{RFM} & Returns, incl. from exceptions or monitors \\
      \hline
      \texttt{condJump} & \texttt{JEQ}, \texttt{JEQ\_A}, \texttt{JGE}, \texttt{JGE\_U}, \texttt{JGEZ}, \texttt{JGTZ}, \texttt{JLEZ}, \texttt{JLT},\newline \texttt{JLT\_U}, \texttt{JLTZ}, \texttt{JNE}, \texttt{JNE\_A}, \texttt{JNED}, \texttt{JNEI}, \texttt{JNZ},\newline \texttt{JNZ\_A}, \texttt{JNZ\_T}, \texttt{JZ}, \texttt{JZ\_A}, \texttt{JZ\_T}, \texttt{LOOP} & \\
      \hline
      \texttt{uncondJump} & \texttt{J}, \texttt{JA}, \texttt{JI}, \texttt{LOOPU} & \\
      \hline
      \texttt{indirectJump} & \texttt{JI} & \\
      \hline
      \texttt{sideEffect} & \texttt{BISR}, \texttt{CACHEA\_I}, \texttt{CACHEA\_W}, \texttt{CACHEA\_WI}, \texttt{DEBUG}, \texttt{DISABLE}, \texttt{DSYNC}, \texttt{ENABLE}, \texttt{ISYNC}, \texttt{LDLCX}, \texttt{LDUCX}, \texttt{MTCR}, \texttt{RSLCX}, \texttt{RSTV}, \texttt{SVLCX}, \texttt{SYSCALL}, \texttt{TLBDEMAP}, \texttt{TLBFLUSH\_A}, \texttt{TLBFLUSH\_B}, \texttt{TLBMAP}, \texttt{TLBPROBE\_A}, \texttt{TLBPROBE\_I}, \texttt{TRAPSV}, \texttt{TRAPV}, \texttt{UPDFL} & Interrupt, cache, Trans-\newline lation Look-Aside Buffer\newline (TLB) handling, debugging, context and core register handling, system call operations \\
      \hline
    \end{tabular}
  \end{center}
\end{table}

Class \texttt{TC\_Const5\_Unsigned} comes into play for word-level bit manipulation and shift operations where bit positions or shift amounts between 0 or 31 bits need to be specified.
3-bit unsigned immediates (class \texttt{TC\_Const3\_\-Unsigned}) are used only by TriCore opcode \texttt{ST\_T} that stores one single bit somewhere inside a byte.
For the specification of the bit's position within its byte, 3 bits are required.
Similarly, 2-bit unsigned immediates of type \texttt{TC\_Const2\_Unsigned} are only used by the special-purpose \texttt{ADDSC\_A} opcode for address computation.
Here, some address is multiplied by 1, 2, 4 or 8 bytes, and this small factor is encoded as its respective logarithm.
One bit immediates of class \texttt{TC\_Const1\_Unsigned} are widely used by the TriCore's operations for fractional arithmetics to handle certain arithmetic corner cases in the Q15 format.

For all abovementioned classes of TriCore immediates, dedicated unit tests exist showing how to use these immediates and demonstrating that invalid values are properly rejected.
These unit tests can be found in files \texttt{WIR/arch/tricore/tests/tcConst*}.

\begin{table}[b]
  \begin{center}
    \caption{Opcode Types of \texttt{TC131} Processor Model}
    \label{tab:tc131opctypes}
    \begin{tabular}{|c|p{79mm}|p{42mm}|}
      \hline
      \multirow{2}{*}[-1mm]{\textbf{\shortstack[c]{Opcode\\ Type}}} & \multirow{2}{*}{\textbf{Opcodes}} & \multirow{2}{*}{\textbf{Comment}} \\
      & & \\
      \hhline{|=|=|=|}
      \texttt{implicitMemoryAccess} & \texttt{CACHEI\_W}, \texttt{CACHEI\_WI} & Data cache index handling  \\
      \hline
      \texttt{sideEffect} & \texttt{CACHEI\_W}, \texttt{CACHEI\_WI} & \\
      \hline
    \end{tabular}
  \end{center}
\end{table}

The entire TriCore 1.3 instruction set~\cite{tricore08} is carefully modeled by \gls{WIR} class \texttt{TC13}, including all its \gls{CPU}, \gls{FPU} and \gls{MMU} instructions.
Due to the sheer complexity of this \gls{ISA}, the \texttt{TC13} model features a total of 332 opcodes, beginning with \texttt{ABS} and ending with \texttt{XOR\_T} (see file \texttt{tc13opcodes.cc}).
Their classification into opcode types is shown in \cref{tab:tc13opctypes}.
It is worth mentioning that there are some operations that are classified as both memory loads and stores.
These operations atomically load some data from memory, potentially modify it, and write it back to memory.
Various operations perform implicit memory accesses, especially those realizing function calls and returns from functions.
This comes from the TriCore's automatic saving and restoring of upper context registers in memory during function calls/returns.
Finally, quite many operations cause side effects to the processor's internal state and should thus be taken with care during \gls{WIR}-based code analyses and optimizations.
One pseudo-opcode named \texttt{MOV\_RR} has been added to the \texttt{TC13} processor model in order to clearly characterize register-register moves.
This opcode \texttt{MOV\_RR} uses the TriCore's regular assembly mnemonic ``\texttt{mov}''.
But since the \gls{ISA} knows various flavors of this \texttt{mov} operation, and only one of them realizes a register-register-move, a pseudo-opcode was added.

The TriCore 1.3.1 version extends the instruction set by 5 additional machine operations.
By exploiting inheritance again (cf. \cref{fig:tricore}), processor model \texttt{TC131} (file \texttt{tc131.h}) simply inherits all abovementioned opcodes from \texttt{TC13} and extends them by the novel opcodes \texttt{CACHEI\_W}, \texttt{CACHEI\_WI}, \texttt{FTOIZ}, \texttt{FTOQ31Z} and \texttt{FTOUZ}.
The opcode classification of these \texttt{TC131} opcodes is summarized in \cref{tab:tc131opctypes}, see also file \texttt{tc131opcodes.cc}.

\begin{table}[t]
  \begin{center}
    \caption{Naming Scheme for Operation Formats of \texttt{TC13} Processor Model}
    \label{tab:tcopf}
    \begin{tabular}{|c|p{136mm}|}
      \hline
      \multirow{2}{*}[-1mm]{\textbf{\shortstack[c]{Operation Format\\ Specifier}}} & \multirow{2}{*}{\textbf{Parameter Type}} \\
      & \\
      \hhline{|=|=|}
      ``\texttt{A}'' & Address register parameter of type \texttt{TC13::RegisterType::aReg} \\
      \hline
      ``\texttt{C1}'' & Unsigned 1-bit immediate constant \\
      \hline
      ``\texttt{C2}'' & Unsigned 2-bit immediate constant \\
      \hline
      ``\texttt{C3}'' & Unsigned 3-bit immediate constant \\
      \hline
      ``\texttt{C4}'' & Signed or unsigned 4-bit immediate constant \\
      \hline
      ``\texttt{C5}'' & Unsigned 5-bit immediate constant \\
      \hline
      ``\texttt{C8}'' & Unsigned 8-bit immediate constant \\
      \hline
      ``\texttt{C9}'' & Signed or unsigned 9-bit immediate constant \\
      \hline
      ``\texttt{C10}'' & Signed or unsigned 10-bit immediate constant \\
      \hline
      ``\texttt{C16}'' & Signed or unsigned 16-bit immediate constant \\
      \hline
      ``\texttt{C18}'' & Unsigned 18-bit immediate constant \\
      \hline
      ``\texttt{D}'' & Data register parameter of type \texttt{TC13::RegisterType::dReg} \\
      \hline
      ``\texttt{E}'' & Extended data register parameter of type \texttt{TC13::RegisterType::eReg} \\
      \hline
      ``\texttt{I}'' & Physical implicit data or address register parameter \texttt{A15} or \texttt{D15} \\
      \hline
      ``\texttt{L}'' & Label parameter \\
      \hline
      ``\texttt{P}'' & Extended address register parameter of type \texttt{TC13::RegisterType::pReg} \\
      \hline
      ``\texttt{PSW}'' & \gls{PSW} register parameter of type \texttt{TC13::RegisterType::pswBit} \\
      \hline
      ``\texttt{S\dots}'' & Prefix for short 16-bit operation formats \\
      \hline
      ``\texttt{SP}'' & Physical stack pointer register parameter \texttt{A10} \\
      \hline
      ``\texttt{SYS}'' & No parameters at all \\
      \hhline{|=|=|}
      ``\texttt{ABSA}'' & Absolute addressing mode \\
      \hline
      ``\texttt{BOA}'' & Base + offset addressing mode \\
      \hline
      ``\texttt{BRA}'' & Bit-reverse addressing mode \\
      \hline
      ``\texttt{CA}'' & Circular addressing mode \\
      \hline
      ``\texttt{PIA}'' & Post- or pre-increment addressing mode \\
      \hline
    \end{tabular}
  \end{center}
\end{table}

The TriCore operation formats are defined in file \texttt{tc13operationformats.cc}.
As the \texttt{TC131} processor model does not require any additional formats, there is no such file for this architecture.
A total of 176 32 bits wide formats are defined, while 54 operation formats are 16 bits wide.
According to the used naming scheme (cf. \cref{tab:tcopf}), the prefix ``\texttt{S}'' indicates such short operation formats.
This naming scheme uses different characters or strings to represent the various register parameter types, immediate types, or label parameters.
Furthermore, some operation formats require the physical TriCore registers \texttt{A15}, \texttt{D15} (implicit address or data registers) or the stack pointer \texttt{A10}.
The presence of these physical registers is indicated by the specifiers ``\texttt{I}'' (implicit register) or ``\texttt{SP}'' (stack pointer).
Operations of format ``\texttt{SYS}'' are 32 bits wide and possess no parameters at all -- this is the case for, e.g., the \texttt{NOP} (no operation) opcode.
Based on this naming scheme, it is easy to see that an operation of, e.g., format \texttt{TC13::OperationFormat::SSPC8} is a short operation of 16 bits width, requiring the physical stack pointer and an unsigned 8-bits immediate as parameters.

In order to see how to construct and use all TriCore opcodes in all supported operation formats, unit tests \texttt{WIR/arch/tricore/tests/tcOperations1.cc} and \texttt{tcOperations2.cc} can be inspected.
The former creates all operations in all flavors for both the \texttt{TC13} and \texttt{TC131} \glspl{ISA} and checks all specific operation properties.
The latter also creates all these operations, but afterwards prints them in TriCore-compliant assembly syntax and feeds this printout into the GNU TriCore assembler.
This way, unit test \texttt{tcOperations2.cc} ensures that \gls{WIR} definitely produces valid TriCore assembly output.

The last 5 rows of \cref{tab:tcopf} already draw the bow towards the modelling of the TriCore's various addressing modes within \gls{WIR}.
TriCore supports absolute, base + offset, bit-reverse, circular, and post- or pre-increment addressing.
In principle, all these addressing modes could be modeled explicitly by inheriting from class \texttt{WIR\_Base\-Processor::AddressingMode} and by exploiting the concepts sketched in \cref{sssec:amodes}.
When, however, having a closer look at the \gls{ISA} design, it turns out that the structure of these addressing modes is so different that most of them require their very own, special combinations of parameters.

\begin{wirExample}
Absolute addressing makes use of 18-bit immediates of type \texttt{TC\_Const18\_Unsigned}, cf. \cref{tab:tcimmtypes}.
In contrast, base + offset addressing needs one address register parameter for the base address, plus either a 10-bit or a 16-bit immediate.
Bit-reverse addressing bases on the provision of an extended address register.
Circular addressing also requires an extended address register, but furthermore needs a 10-bit immediate parameter.
Finally, post- or pre-increment addressing needs an address register parameter and a 10-bit immediate.
\end{wirExample}

These examples show that the modelling of each addressing mode inherently leads to dedicated, specialized operation formats that have to be modeled, too.
For the sake of simplicity, and to avoid modelling the same aspect twice within \gls{WIR}, only the post- and pre-increment addressing modes are modeled using class \texttt{AddressingMode}.
In contrast, the special structure of all other addressing modes is made explicit already in the underlying operation formats by introducing the specifiers ``\texttt{ABSA}'' to ``\texttt{PIA}'' from the bottom of \cref{tab:tcopf}.
The \texttt{TC13} modes for post- and pre-increment addressing have already been introduced earlier, see \cref{lst:tcadrmode} on \cpageref{lst:tcadrmode}.

As can be seen from the following example, only the structure of the two post- and pre-increment addressing modes is identical, justifying their modelling using class \texttt{WIR\_BaseProcessor::AddressingMode}.
The structures of all other addressing modes all differ from each other so that the remaining addressing modes are properly and uniquely specified by their underlying operation formats.

\begin{wirExample}
The \texttt{LD\_W} operation loading a word from memory to a data register supports all the abovementioned addressing modes.
Assuming that \texttt{b} is a \texttt{WIR\_BasicBlock} and that \texttt{a}, \texttt{d} and \texttt{p} are address, data and extended address registers, resp., the following lines of code add \texttt{LD\_W} operations with the various addressing modes:

{\small
\verb|  |\texttt{b.pushBackInstruction\textcolor{blue}{(} \textcolor{gray}{// Absolute Addressing}} \\
\verb|    |\texttt{\{ \{ TC13::OpCode::LD\_W, TC13::OperationFormat::DC18ABSA,} \\
\verb|        |\texttt{WIR\_RegisterParameter\textcolor{blue}{(} d, WIR\_Usage::def \textcolor{blue}{)}, TC\_Const18\_Unsigned\textcolor{blue}{(} \textcolor{orange}{16383} \textcolor{blue}{)} \} \} \textcolor{blue}{)};} \\
\verb|  |\texttt{b.pushBackInstruction\textcolor{blue}{(} \textcolor{gray}{// Base + Long Offset Addressing}} \\
\verb|    |\texttt{\{ \{ TC13::OpCode::LD\_W, TC13::OperationFormat::DAC16BOA,} \\
\verb|        |\texttt{WIR\_RegisterParameter\textcolor{blue}{(} d, WIR\_Usage::def \textcolor{blue}{)},} \\
\verb|        |\texttt{WIR\_RegisterParameter\textcolor{blue}{(} a, WIR\_Usage::use \textcolor{blue}{)}, TC\_Const16\_Signed\textcolor{blue}{(} \textcolor{orange}{-32768} \textcolor{blue}{)} \} \} \textcolor{blue}{)};} \\
\verb|  |\texttt{b.pushBackInstruction\textcolor{blue}{(} \textcolor{gray}{// Bit-Reverse Addressing}} \\
\verb|    |\texttt{\{ \{ TC13::OpCode::LD\_W, TC13::OperationFormat::DPBRA,} \\
\verb|        |\texttt{WIR\_RegisterParameter\textcolor{blue}{(} d, WIR\_Usage::def \textcolor{blue}{)},} \\
\verb|        |\texttt{WIR\_RegisterParameter\textcolor{blue}{(} p, WIR\_Usage::defuse \textcolor{blue}{)} \} \} \textcolor{blue}{)};} \\
\verb|  |\texttt{b.pushBackInstruction\textcolor{blue}{(} \textcolor{gray}{// Circular Addressing}} \\
\verb|    |\texttt{\{ \{ TC13::OpCode::LD\_W, TC13::OperationFormat::DPC10CA,} \\
\verb|        |\texttt{WIR\_RegisterParameter\textcolor{blue}{(} d, WIR\_Usage::def \textcolor{blue}{)},} \\
\verb|        |\texttt{WIR\_RegisterParameter\textcolor{blue}{(} p, WIR\_Usage::defuse \textcolor{blue}{)}, TC\_Const10\_Signed\textcolor{blue}{(} \textcolor{orange}{-512} \textcolor{blue}{)} \} \} \textcolor{blue}{)};} \\
\verb|  |\texttt{b.pushBackInstruction\textcolor{blue}{(} \textcolor{gray}{// Post-Increment Addressing}} \\
\verb|    |\texttt{\{ \{ TC13::OpCode::LD\_W, TC13::OperationFormat::DAC10PIA,} \\
\verb|        |\texttt{WIR\_RegisterParameter\textcolor{blue}{(} d, WIR\_Usage::def \textcolor{blue}{)},} \\
\verb|        |\texttt{WIR\_AddressingModeParameter\textcolor{blue}{(} TC13::AddressingMode::post \textcolor{blue}{)},} \\
\verb|        |\texttt{WIR\_RegisterParameter\textcolor{blue}{(} a, WIR\_Usage::defuse \textcolor{blue}{)}, TC\_Const10\_Signed\textcolor{blue}{(} \textcolor{orange}{-512} \textcolor{blue}{)} \} \} \textcolor{blue}{)};} \\
\verb|  |\texttt{b.pushBackInstruction\textcolor{blue}{(} \textcolor{gray}{// Pre-Increment Addressing}} \\
\verb|    |\texttt{\{ \{ TC13::OpCode::LD\_W, TC13::OperationFormat::DAC10PIA,} \\
\verb|        |\texttt{WIR\_RegisterParameter\textcolor{blue}{(} d, WIR\_Usage::def \textcolor{blue}{)},} \\
\verb|        |\texttt{WIR\_AddressingModeParameter\textcolor{blue}{(} TC13::AddressingMode::pre \textcolor{blue}{)},} \\
\verb|        |\texttt{WIR\_RegisterParameter\textcolor{blue}{(} a, WIR\_Usage::defuse \textcolor{blue}{)}, TC\_Const10\_Signed\textcolor{blue}{(} \textcolor{orange}{-512} \textcolor{blue}{)} \} \} \textcolor{blue}{)};}
}
\end{wirExample}

\gls{WIR} provides the unit test \texttt{WIR/arch/tricore/tests/tcAddressingModeParameter1.cc} that illustrates how to use and work with these TriCore-specific addressing modes.

\begin{lstlisting}[style=code,float=t,caption={Testing of Pipeline Properties of Class \texttt{TC13}},label=lst:tcpipetests]
// Determine whether an operation gets executed on the TriCore Integer Pipeline.
static @bool@ TC13::isIP@(@ const WIR_Operation æ&æo @)@;Ü\label{lst:tcpipetests1}Ü

// Determine whether an operation gets executed on the TriCore Loop Pipeline.
static @bool@ TC13::isLP@(@ const WIR_Operation æ&æo @)@;Ü\label{lst:tcpipetests2}Ü

// Determine whether an operation gets executed on the TriCore Load/Store Pipeline.
static @bool@ TC13::isLS@(@ const WIR_Operation æ&æo @)@;Ü\label{lst:tcpipetests3}Ü
\end{lstlisting}

The TriCore architecture has three different internal pipelines -- the standard integer pipeline, a load/store pipeline reserved for memory accesses and address arithmetics, and a loop pipeline realizing zero-overhead loops.
Under certain conditions~\cite{tricore03}, these pipelines can execute in parallel and the TriCore is able to issue two operations in parallel.
In order to be able to exploit this parallelism, a compiler needs to know on which of these pipelines an operation gets executed.
This is supported by the \texttt{TC13} processor model using the methods shown in \cref{lst:tcpipetests}.
They return \texttt{true} if a given operation is executed on the integer, load/store, or loop pipeline, resp.
\clearpage


% - - - - - - - - - - - - - - - - - - - - - - - - - - - - - - - - - - - -
\subsection{RISC-V RV32 Processor Models}
\label{ssec:rv32model}
% - - - - - - - - - - - - - - - - - - - - - - - - - - - - - - - - - - - -

Key feature of the RISC-V RV32 \gls{ISA}~\cite{riscv17} is the simplicity of the instruction set on the one hand, that resembles that of the MIPS model from \cref{ssec:mipsmodel}.
On the other hand, the RISC-V operations are organized in a highly modular fashion, allowing to include or not to include a couple of optional standard \gls{ISA} extensions.
This design has the effect that each and every RV32 implementation must include the mandatory base integer instruction set RV32I.
On top of this base \gls{ISA}, the extensions listed in \cref{tab:rv32isapackets} can be added.
Concrete RISC-V RV32 implementations are thus characterized by the prefix letters shown in \cref{tab:rv32isapackets} of the base \gls{ISA} plus all included extensions.

\begin{table}[h]
  \begin{center}
    \caption{Standard Extensions of the RISC-V RV32 Instruction Set}
    \label{tab:rv32isapackets}
    \begin{tabular}{|l|l|}
      \hline
      \textbf{Prefix} & \textbf{Extension} \\
      \hhline{|=|=|}
      I & Mandatory Base Integer Instruction Set \\
      \hhline{|=|=|}
      M & Integer Multiplication and Division \\
      \hline
      A & Atomic Operations \\
      \hline
      F & Single-Precision Floating-Point \\
      \hline
      D & Double-Precision Floating-Point \\
      \hline
      C & Compressed Operations \\
      \hline
    \end{tabular}
  \end{center}
\end{table}

Within the \gls{WIR} library, the base instruction set RV32I is modeled, as well as the standard extensions C, F, D, and M.
The extensibility of the \gls{ISA} is perfectly supported by the high modularity of the general \gls{WIR} hardware model.
This allows to provide a base class \texttt{RV32I} that models the base integer instruction set.
Using inheritance, standard extensions like, e.g., C or M are modeled leading to derived classes such as \texttt{RV32IC} or \texttt{RV32IM}.
Multiple inheritance is then used to simply derive all possible combinations of the base \gls{ISA} and potential standard extensions.
This is illustrated in \cref{fig:riscv32} showing the complete inheritance relationships between RISC-V related classes within \gls{WIR}.

\begin{figure}[ht]
  \begin{center}
    \includegraphics[width=\textwidth]{riscv.png}
  \end{center}
  \caption{Simplified Inheritance Diagram of RISC-V RV32 Processor Models}
  \label{fig:riscv32}
\end{figure}

The RV32I base integer architecture defines a total of 32 general-purpose integer registers \texttt{x0} -- \texttt{x31} that all have a width of 32 bits.
Thus, the \texttt{RV32I} processor model declares one register type \texttt{RV32I::RegisterType::reg} which is defined in file \texttt{WIR/arch/riscv/rv32iregistertypes.cc}.
If the \gls{ISA} is extended by some floating-point operations, i.e., by extensions F, D or Q from \cref{tab:rv32isapackets}, another set of 32 dedicated floating-point registers \texttt{f0} -- \texttt{f31} is added to the architecture.
Thus, \gls{WIR} class \texttt{RV32IF} introduces the second register type \texttt{fReg} (cf. file \texttt{rv32ifregistertypes.cc}) with a default width of 32 bits.
The \texttt{RV32IFD} model simply inherits these registers from its respective base class \texttt{RV32IF} and adjusts the width of the floating-point registers to 64 bits.

\begin{table}[t]
  \begin{center}
    \caption{Register Types, Classes and Methods of \texttt{RV32I}, \texttt{RV32IF} and \texttt{RV32IFD} Processor Models}
    \label{tab:rv32regtypes}
    \begin{tabular}{|c|c|c|l|l|l|l|}
      \hline
      \multirow{2}{*}[-1mm]{\textbf{\shortstack[c]{Register\\ Type}}} & \multirow{2}{*}[-1mm]{\textbf{\shortstack[c]{Bit\\ Width}}} & \multirow{2}{*}[-1mm]{\textbf{\shortstack[c]{Has\\ Childs}}} & \multicolumn{2}{c|}{\textbf{Physical Registers}} & \multicolumn{2}{c|}{\textbf{Virtual Registers}} \\
      & & & Class & Prefix & Class & Prefix \\
      \hhline{|=|=|=|=|=|=|=|}
      \texttt{reg} & 32 & & \texttt{RV\_RegP} & ``\texttt{x}'' & \texttt{RV\_RegV} & ``\texttt{x\_}'' \\
                   &    & & \multicolumn{2}{c|}{\texttt{x0()} -- \texttt{x31()}} & & \\
      \hhline{|=|=|=|=|=|=|=|}
      \texttt{fReg} & 32 or 64 & & \texttt{RV\_FRegP} & ``\texttt{f}'' & \texttt{RV\_FRegV} & ``\texttt{f\_}'' \\
                   &    & & \multicolumn{2}{c|}{\texttt{f0()} -- \texttt{f31()}} & & \\
      \hline
    \end{tabular}
  \end{center}
\end{table}

The key characteristics of these registers and their types are again summarized in \cref{tab:rv32regtypes}.
As can be seen, none of these register types are hierarchical.\TODO{Correct this!}
The classes implementing physical and virtual registers of these types are \texttt{RV\_RegP} and \texttt{RV\_RegV} (files \texttt{rvregphysical.h} and \texttt{rvregvirtual.h}) and \texttt{RV\_FRegP} and \texttt{RV\_FRegV} (files \texttt{rvfregphysical.h} and \texttt{rvfregvirtual.h}), resp.
As usual, the processor models \texttt{RV32I} and \texttt{RV32IF} provide methods \texttt{x0()} -- \texttt{x31()} and \texttt{f0()} -- \texttt{f31()} for the convenient access to the models' physical registers.

\begin{lstlisting}[style=code,float=h,caption={Testing of Physical Register Properties of Class \texttt{RV32I}},label=lst:rv32regtests]
// Determine whether a register is the RISC-V's null register x0.
static @bool@ RV32I::isX0@(@ const WIR_BaseRegister æ&ær @)@;Ü\label{lst:rv32regtests1}Ü

// Determine whether a register is the RISC-V's link register x1.
static @bool@ RV32I::isX1@(@ const WIR_BaseRegister æ&ær @)@;Ü\label{lst:rv32regtests2}Ü

// Determine whether a register is the RISC-V's stack pointer x2.
static @bool@ RV32I::isSP@(@ const WIR_BaseRegister æ&ær @)@;Ü\label{lst:rv32regtests3}Ü

// Determine whether a register is the RISC-V's frame pointer x8.
static @bool@ RV32I::isFP@(@ const WIR_BaseRegister æ&ær @)@;Ü\label{lst:rv32regtests4}Ü
\end{lstlisting}

Physical register \texttt{x0} is the null register that is always hardwired to the constant 0.
Register \texttt{x1} is the link register of the RV32 architecture, \texttt{x2} the stack pointer, and register \texttt{x8} serves as frame pointer.
In order to check whether a given register \texttt{r} is one of these particular RISC-V registers, class \texttt{RV32I} contains the methods shown in \cref{lst:rv32regtests}.
They return \texttt{true} if \texttt{r} is exactly the physical register to be identified, or if \texttt{r} is virtual and precolored accordingly.

As was the case for the previous \gls{WIR} processor models, also the RISC-V model comes with dedicated unit tests showing how to use the physical and virtual registers shown in \cref{tab:rv32regtypes}.
These tests can be found in files \texttt{WIR/arch/riscv/tests/rv*Reg*}.

The simplicity of the RISC-V's \gls{ISA} is reflected by the few different types of immediate parameters that had to be modeled within \gls{WIR}, see \cref{tab:rv32immtypes}.
For the base class \texttt{RV32I}, only three different kinds of immediates are required.
20-bit immediates of type \texttt{RV\_Const20\_Unsigned} are used to set the most-significant 20 bits of some register.
For the remaining least-significant bits, constants with a width of 12 bits are sufficient which are modeled by class \texttt{RV\_Const12\_Signed}.
With only these two kinds of immediates, any 32-bit value can be loaded into a register and all arithmetic operations of the RV32I \gls{ISA} exploiting immediates can be expressed.
As was the case already before for the MIPS and the TriCore \glspl{ISA}, specialized 5-bit immediates are required to encode shift amounts.

\begin{table}[t]
  \begin{center}
    \caption{Classes for Immediates of \texttt{RV32I} and \texttt{RV32IC} Processor Models}
    \label{tab:rv32immtypes}
    \begin{tabular}{|c|c|l|l|p{55mm}|}
      \hline
      \multirow{2}{*}[-1mm]{\textbf{\shortstack[c]{Immediate\\ Type}}} & \multirow{2}{*}[-1mm]{\textbf{\shortstack[c]{Bit\\ Width}}} & \multirow{2}{*}{\textbf{Class}} & \multirow{2}{*}{\textbf{Value Range}} & \multirow{2}{*}{\textbf{Comment}} \\
      & & & & \\
      \hhline{|=|=|=|=|=|}
      Unsigned 20-bit & 20 & \texttt{RV\_Const20\_Unsigned} & $[0$, $+2^{20} - 1]$ & \\
      \hline
      Signed 12-bit & 12 & \texttt{RV\_Const12\_Signed} & $[-2^{11}$, $+2^{11} - 1]$ & \\
      \hline
      Unsigned 5-bit & 5 & \texttt{RV\_Const5\_Unsigned} & $[0$, $+31]$ & Shift amounts for 32-bit architecture \\
      \hhline{|=|=|=|=|=|}
      Signed 10-bit & 10 & \texttt{RV\_Const10\_Signed} & $[-512$, $+496]$ & Non-zero stack pointer adjustment,\newline must be multiples of 16 \\
      \hline
      Unsigned 10-bit & 10 & \texttt{RV\_Const10\_Unsigned} & $[4$, $+1020]$ & Non-zero stack pointer-relative off-\newline sets, must be multiples of 4 \\
      \hline
      Unsigned 8-bit & 8 & \texttt{RV\_Const8\_Unsigned} & $[0$, $+255]$ & \\
      \hline
      Unsigned 7-bit & 7 & \texttt{RV\_Const7\_Unsigned} & $[-64$, $+60]$ & Pointer-relative offsets, must be\newline multiples of 4 \\
      \hline
      Signed 6-bit & 6 & \texttt{RV\_Const6\_Signed} & $[-32$, $+31]$ & \\
      \hline
      Unsigned 6-bit & 6 & \texttt{RV\_Const6\_Unsigned} & $[0$, $+63]$ & \\
      \hline
    \end{tabular}
  \end{center}
\end{table}

Due to the fact that the C \gls{ISA} extension defines compressed operations of only 16 bits width, the \texttt{RV32IC} model requires additional immediate types with shorter bit widths.
These immediate types are summarized in the lower part of \cref{tab:rv32immtypes}.
The 10-bit immediates dedicated to stack pointer handling have to be non-zero and must be multiples of either 16 or 4.
The 7-bit unsigned immediates for compact memory accesses have to be multiples of 4.
For all other immediate types, there exist no special provisions to constrain these value ranges any further.

The unit tests from files \texttt{WIR/arch/riscv/tests/rvConst*} provide code examples for all modeled RISC-V immediate classes.

\begin{table}[!b]
  \begin{center}
    \caption{Opcode Types of \texttt{RV32I} Processor Model}
    \label{tab:rv32iopctypes}
    \begin{tabular}{|c|p{95mm}|p{42mm}|}
      \hline
      \multirow{2}{*}[-1mm]{\textbf{\shortstack[c]{Opcode\\ Type}}} & \multirow{2}{*}{\textbf{Opcodes}} & \multirow{2}{*}{\textbf{Comment}} \\
      & & \\
      \hhline{|=|=|=|}
      \texttt{load} & \texttt{LB}, \texttt{LBU}, \texttt{LH}, \texttt{LHU}, \texttt{LW} & \\
      \hline
      \texttt{store} & \texttt{SB}, \texttt{SH}, \texttt{SW} & \\
      \hline
      \texttt{move} & \texttt{MOV} & Pseudo-opcode \\
      \hline
      \texttt{call} & \texttt{JAL} & \\
      \hline
      \texttt{ret} & \texttt{RET} & Pseudo-opcode \\
      \hline
      \texttt{condJump} & \texttt{BEQ}, \texttt{BGE}, \texttt{BGEU}, \texttt{BLT}, \texttt{BLTU}, \texttt{BNE} & \\
      \hline
      \texttt{uncondJump} & \texttt{J} & Pseudo-opcode \\
      \hline
      \texttt{indirectJump} & \texttt{JALR} & \\
      \hline
      \texttt{sideEffect} & \texttt{CSRRC}, \texttt{CSRRCI}, \texttt{CSRRS}, \texttt{CSRRSI}, \texttt{CSRRW}, \texttt{CSRRWI}, \texttt{EBREAK}, \texttt{ECALL} & \gls{CSR} handling, environment calls/breakpoints \\
      \hline
    \end{tabular}
  \end{center}
\end{table}

Out of the overall 47 operations defined by the base integer instruction set, all of them are modeled within class \texttt{RV32I} except the \texttt{FENCE} and \texttt{FENCEI} operations (see file \texttt{rv32iopcodes.cc}).
Their classification into opcode types is shown in \cref{tab:rv32iopctypes}.
The following three pseudo-opcodes have been added for convenience reasons and should consequently be used for moves, returns and unconditional jumps:
\begin{itemize}
  \item \texttt{J} represents an unconditional branch and gets substituted by \texttt{jal x0, \emph{label}}
  \item \texttt{MOV} denotes a register-register move that is replaced by \texttt{addi \emph{tgt}, \emph{src}, 0}
  \item \texttt{RET} realizes a return from a function call that boils down to \texttt{jalr x0, x1, 0}
\end{itemize}

Using the class \texttt{RV32IC} or its derivatives extends the processor model by 27 additional operations ranging from \texttt{CADD} to \texttt{CXOR} (cf. file \texttt{rv32icopcodes.cc}).
Yet another 27 operations (\texttt{FADD\_S} to \texttt{FSW}) for single-precision floating point computations are added by class \texttt{RV32IF} (file \texttt{rv32ifopcodes.cc}).
Likewise for file \texttt{rv32ifdopcodes.cc} defining the 27 double-precision floating-point operations \texttt{FADD\_D} to \texttt{FSUB\_D}.
Finally, all 8 operations \texttt{DIV} to \texttt{REMU} for integer division and multiplication (cf. file \texttt{rv32imopcodes.cc}) can be added by using class \texttt{RV32IM} or derivatives.
The classification of all these opcodes implementing the RISC-V standard extensions is summarized in \cref{tab:rv32icopctypes}.
Each of the two floating-point extensions \texttt{RV32IF} and \texttt{RV32IFD} specifies one pseudo-opcode representing moves of floating-point registers: \texttt{FMV\_S} and \texttt{FMV\_D}, resp.
As usual, operations that are not listed in \cref{tab:rv32iopctypes} or \cref{tab:rv32icopctypes} carry no particular attributes within \gls{WIR} and thus usually represent standard arithmetical operations.

\begin{table}[t]
  \begin{center}
    \caption{Opcode Types of \texttt{RV32IC}, \texttt{RV32IF} and \texttt{RV32IFD} Processor Models}
    \label{tab:rv32icopctypes}
    \begin{tabular}{|c|p{29mm}|p{29mm}|p{29mm}|p{42mm}|}
      \hline
      \multirow{2}{*}[-1mm]{\textbf{\shortstack[c]{Opcode\\ Type}}} & \multicolumn{3}{c|}{\textbf{Opcodes}} & \multirow{2}{*}{\textbf{Comment}} \\
      & \texttt{RV32IC} & \texttt{RV32IF} & \texttt{RV32IFD} & \\
      \hhline{|=|=|=|=|=|}
      \texttt{load} & \texttt{CLW}, \texttt{CLWSP} & \texttt{FLW} & \texttt{FLD} & \\
      \hline
      \texttt{store} & \texttt{CSW}, \texttt{CSWSP} & \texttt{FSW} & \texttt{FSD} & \\
      \hline
      \texttt{move} & \texttt{CMV} & \texttt{FMV\_S} & \texttt{FMV\_D} & Pseudo-opcodes \texttt{FMV\_S}\newline and \texttt{FMV\_D} \\
      \hline
      \texttt{call} & \texttt{CJAL} & & & \\
      \hline
      \texttt{indirectCall} & \texttt{CJALR} & & & \\
      \hline
      \texttt{condJump} & \texttt{CBEQZ}, \texttt{CBNEZ} & & & \\
      \hline
      \texttt{uncondJump} & \texttt{CJ} & & & \\
      \hline
      \texttt{indirectJump} & \texttt{CJR} & & & \\
      \hline
      \texttt{sideEffect} & \texttt{CEBREAK} & \texttt{FEQ\_S}, \texttt{FLE\_S},\newline \texttt{FLT\_S} & \texttt{FEQ\_D}, \texttt{FLE\_D},\newline \texttt{FLT\_D} & \\
      \hline
    \end{tabular}
  \end{center}
\end{table}

Base class \texttt{RV32I} makes use of a total of 17 operation formats that are all 32 bits wide (cf. file \texttt{rv32ioperation\-formats}).
The naming scheme used in the context of the RV32 processor model within \gls{WIR} is shown in the upper part of \cref{tab:rv32iopf}.
It is worth mentioning that two operation formats (\texttt{RV32I::OperationFormat::RSC5\_1} and \texttt{RSR\_1}) make use of unstructured string parameters (cf. \cref{ssec:strparams}).
These formats are exclusively used by the RISC-V's operations manipulating the \glspl{CSR}.
As the \gls{ISA} specification~\cite{riscv17} lacks a proper definition of available \glspl{CSR} and how they shall be named in assembly code, these \gls{CSR} names have to be provided as plain strings to the corresponding \texttt{RV32I} operations.

\begin{wirWarn}
Due to the fully unstructured nature of string parameters, the \gls{WIR} library has no means to distinguish between valid and invalid strings used to identify \glspl{CSR}.
Thus, it is easy to generate invalid \gls{WIR} operations here, it lies in the sole responsibility of application code to ensure that only valid \gls{CSR} names are used as string parameters.
\end{wirWarn}

The 16 additional operation formats added by class \texttt{RV32IC} (cf. file \texttt{rv32icoperationformats.cc}) are all 16 bits wide, and their names -- as usual within \gls{WIR} -- all start with the prefix ``\texttt{S}'' (cf. middle part of \cref{tab:rv32iopf}).
Additionally, the naming scheme is extended to also capture the 8- and 6-bits immediates required by the C \gls{ISA} extension.
Finally, extension \texttt{RV32IF} introduces 7 more 32-bit operation formats (file \texttt{rv32ifoperationformats.cc}), and the specifier ``\texttt{F}'' here denotes a floating-point register parameter (cf. bottom part of \cref{tab:rv32iopf}).

Examples how to deal with all modeled RISC-V opcodes and operation formats can be found in unit tests \texttt{WIR/arch/riscv/tests/rv32*Operations*}.
In analogy to the TriCore model, the unit tests ending with the number 1 (e.g., \texttt{rv32iOperations1.cc}) instantiate and verify all opcodes in all formats and test their individual attributes and properties.
Unit tests ending with the number 2 (e.g., \texttt{rv32imfdcOperations2.cc}) print out all modeled operations and pass them through the GNU RISC-V assembler.

\begin{table}[!t]
  \begin{center}
    \caption{Naming Scheme for Operation Formats of \texttt{RV32I}, \texttt{RV32IC} and \texttt{RV32IF} Processor Models}
    \label{tab:rv32iopf}
    \begin{tabular}{|c|p{136mm}|}
      \hline
      \multirow{2}{*}[-1mm]{\textbf{\shortstack[c]{Operation Format\\ Specifier}}} & \multirow{2}{*}{\textbf{Parameter Type}} \\
      & \\
      \hhline{|=|=|}
      ``\texttt{C12}'' & Signed 12-bit immediate constant \\
      \hline
      ``\texttt{C20}'' & Unsigned 20-bit immediate constant \\
      \hline
      ``\texttt{C5}'' & Unsigned 5-bit immediate constant \\
      \hline
      ``\texttt{L}'' & Label parameter \\
      \hline
      ``\texttt{NULL}'' & No parameters at all \\
      \hline
      ``\texttt{R}'' & Register parameter of type \texttt{RV32I::RegisterType::reg} \\
      \hline
      ``\texttt{S}'' & Unstructured string parameter \\
      \hhline{|=|=|}
      ``\texttt{C10}'' & Signed or unsigned 10-bit immediate constant \\
      \hline
      ``\texttt{C8}'' & Unsigned 8-bit immediate constant \\
      \hline
      ``\texttt{C7}'' & Unsigned 7-bit immediate constant \\
      \hline
      ``\texttt{C6}'' & Signed or unsigned 6-bit immediate constant \\
      \hline
      ``\texttt{S\dots}'' & Prefix for short 16-bit operation formats \\
      \hhline{|=|=|}
      ``\texttt{F}'' & Floating-point register parameter of type \texttt{RV32I::RegisterType::fReg} \\
      \hline
    \end{tabular}
  \end{center}
\end{table}

Since the RISC-V \gls{ISA} only knows one addressing mode, i.e., base + offset addressing, no addressing modes are explicitly modeled within \gls{WIR}.


% - - - - - - - - - - - - - - - - - - - - - - - - - - - - - - - - - - - -
\subsection{ARM Processor Models}
\label{ssec:armmodel}
% - - - - - - - - - - - - - - - - - - - - - - - - - - - - - - - - - - - -

The ARM-based \glspl{ISA}~\cite{arm05} are organized in a hierarchical manner with an ARMv4 instruction set as base (older \gls{ISA} versions v1, v2 or v3 are considered obsolete).
On top of this base, various more recent instruction set versions have been defined over the years.
This structure is properly reflected by the \gls{WIR} processor models which, again, exploit inheritance to realize a class hierarchy ranging from the base \texttt{ARMv4} model to the \texttt{ARMv6} processor model.
This class hierarchy is depicted in \cref{fig:arm}.

\begin{figure}[ht]
  \begin{center}
    \includegraphics[width=\textwidth]{arm.png}
  \end{center}
  \caption{Simplified Inheritance Diagram of ARM Processor Models}
  \label{fig:arm}
\end{figure}

The ARM architecture has a total of 16 general-purpose registers \texttt{r0} -- \texttt{r15} that are all 32 bits wide.
These registers are represented by register type \texttt{ARMv4::RegisterType::reg} that is defined in file \texttt{WIR/arch/arm/arm\-registertypes.cc}, see also \cref{tab:armregtypes}.
As usual, there exist implementations of classes \texttt{ARM\_RegP} and \texttt{ARM\_RegV} modelling both physical and virtual registers of this type (files \texttt{armregphysical.h} and \texttt{armrergvir\-tual.h}).
However, the ARM base class \texttt{ARMv4} does not provide any access methods for the quick access to the architecture's physical registers of type \texttt{reg}.
This is motivated by the fact that the ARM \gls{ISA} regularly partitions this general-purpose register file into high and low registers.
16-bit ARM THUMB operations can only access low registers, while the 32-bit ARM operations can access both low and high registers, i.e., the entire register file.
Due to this, the \texttt{ARMv4} processor model within \gls{WIR} furthermore defines the register types \texttt{hi} and \texttt{lo} as can be seen from \cref{tab:armregtypes}.
As their names suggest, these types stand for registers \texttt{r8} -- \texttt{r15} and \texttt{r0} -- \texttt{r7}, resp.
Class \texttt{ARMv4} now includes convenient access methods \texttt{R8()} -- \texttt{R15()} that return references to the physical \texttt{hi} registers, while methods \texttt{R0()} -- \texttt{R7()} provide access to the physical \texttt{lo} registers.
Physical and virtual \texttt{hi} registers are realized by classes \texttt{ARM\_HiRegP} and \texttt{ARM\_HiRegV} (files \texttt{armhiregphysical.h} and \texttt{armhiregvirtual.h}), resp.
Analogously, classes \texttt{ARM\_LoRegP} and \texttt{ARM\_LoRegV} from files \texttt{armloregphysical.h} and \texttt{armloregvirtual.h} model the ARM's \texttt{lo} registers, resp.
None of the register types defined for the \texttt{ARMv4} model is hierarchical, register \texttt{R13} is the ARM's stack pointer.

Hierarchical registers are introduced beginnig with the \texttt{ARMv5TE} model.
These extended ARM registers group neighboring pairs of 32-bit general-purpose registers together in order to form 64-bit registers.
For example, standard registers \texttt{R4} and \texttt{R5} form the extended 64-bit register \texttt{P4}.
Thus, the \texttt{ARMv5TE} model introduces the hierarchical register type \texttt{pReg}.
Physical registers of this type are implemented by class \texttt{ARMv5TE\_PRegP} (file \texttt{armv5tepregphysical.h}), while class \texttt{ARv5TE\_PregV} relizes virtual extended registers (file \texttt{armv5tepregvir\-tual.h}).
Consequently, class \texttt{ARMv5TE} provides access methods \texttt{P0()} -- \texttt{P14()} for the easy access to these physical extended registers.

Unit tests \texttt{WIR/arch/arm/tests/armPhRegs1.cc} and \texttt{armVirtualRegister1.cc} provide code examples for the usage of all \gls{WIR} classes modeling both physical and virtual ARM registers.

\begin{table}[!b]
  \begin{center}
    \caption{Register Types, Classes and Methods of \texttt{ARMv4} Processor Model}
    \label{tab:armregtypes}
    \begin{tabular}{|c|c|c|l|l|l|l|}
      \hline
      \multirow{2}{*}[-1mm]{\textbf{\shortstack[c]{Register\\ Type}}} & \multirow{2}{*}[-1mm]{\textbf{\shortstack[c]{Bit\\ Width}}} & \multirow{2}{*}[-1mm]{\textbf{\shortstack[c]{Has\\ Childs}}} & \multicolumn{2}{c|}{\textbf{Physical Registers}} & \multicolumn{2}{c|}{\textbf{Virtual Registers}} \\
      & & & Class & Prefix & Class & Prefix \\
      \hhline{|=|=|=|=|=|=|=|}
      \texttt{reg} & 32 & & \texttt{ARM\_RegP} & ``\texttt{r}'' & \texttt{ARM\_RegV} & ``\texttt{r\_}'' \\
                   &    & & \multicolumn{2}{c|}{} & & \\
      \hline
      \texttt{hi} & 32 & & \texttt{ARM\_HiRegP} & ``\texttt{r}'' & \texttt{ARM\_HiRegV} & ``\texttt{h\_}'' \\
                   &    & & \multicolumn{2}{c|}{\texttt{R8()} -- \texttt{R15()}} & & \\
      \hline
      \texttt{lo} & 32 & & \texttt{ARM\_LoRegP} & ``\texttt{r}'' & \texttt{ARM\_LoRegV} & ``\texttt{l\_}'' \\
                   &    & & \multicolumn{2}{c|}{\texttt{R0()} -- \texttt{R7()}} & & \\
      \hhline{|=|=|=|=|=|=|=|}
      \texttt{pReg} & 64 & \checkmark   & \texttt{ARMv5TE\_PRegP} & ``\texttt{p}'' & \texttt{ARMv5TE\_PRegV} & ``\texttt{p\_}'' \\
                    &    & \texttt{reg} & \multicolumn{2}{c|}{\texttt{P0()}, \texttt{P2()}, \dots, \texttt{P14()}} & & \\
      \hline
    \end{tabular}
  \end{center}
\end{table}

The ARM \gls{ISA} is characterized by a rather small and compact base of opcodes provided by class \texttt{ARMv4}, resembling a classical \gls{RISC} design.
Similar to the modular approach of the RISC-V architecture (cf. \cref{ssec:rv32model}), this base \gls{ISA} can be extended by additional operations.
The main difference of ARM operations compared to previous architectures, however, lies in the enormous complexity of the operations' encodings themselves.
The ARM \gls{ISA} allows to generate a huge number of different variants of one and the same machine operation, which resembles more a classical \gls{CISC} architecture design.
This is mostly due to the various highly complex addressing modes that the ARM \gls{ISA} offers, and the various operation encodings have the effect that a rather large number of different immediate types are required.

\begin{table}[t]
  \begin{center}
    \caption{Classes for Immediates of \texttt{ARMv4}, \texttt{ARMv4T}, \texttt{ARMv5T} and \texttt{ARMv6} Processor Models}
    \label{tab:armimmtypes}
    \begin{tabular}{|c|c|l|l|p{39mm}|}
      \hline
      \multirow{2}{*}[-1mm]{\textbf{\shortstack[c]{Immediate\\ Type}}} & \multirow{2}{*}[-1mm]{\textbf{\shortstack[c]{Bit\\ Width}}} & \multirow{2}{*}{\textbf{Class}} & \multirow{2}{*}{\textbf{Value Range}} & \multirow{2}{*}{\textbf{Comment}} \\
      & & & & \\
      \hhline{|=|=|=|=|=|}
      Unsigned 24-bit & 24 & \texttt{ARM\_Const24\_Unsigned} & $[0$, $+2^{24} - 1]$ & \\
      \hline
      Unsigned 12-bit & 12 & \texttt{ARM\_Const12\_Unsigned} & $[0$, $+2^{12} - 1]$ & \\
      \hline
      Unsigned 10-bit & 10 & \texttt{ARM\_Const10\_CoprocessorOffset} & $[0$, $+1023]$ & Coprocessor offsets,\newline must be multiples of 4 \\
      \hline
      Unsigned 8-bit & 8 & \texttt{ARM\_Const8\_Unsigned} & $[0$, $+255]$ & \\
      \hline
      Unsigned 6-bit & 6 & \texttt{ARM\_Const6\_Unsigned0} & $[+1$, $+32]$ & Value 0 not included \\
      \hline
      Unsigned 5-bit & 5 & \texttt{ARM\_Const5\_Unsigned} & $[0$, $+31]$ & \\
      \hline
      Unsigned 5-bit & 5 & \texttt{ARM\_Const5\_Unsigned0} & $[+1$, $+31]$ & Value 0 not included \\
      \hline
      Unsigned 5-bit & 5 & \texttt{ARM\_Const5\_RotateAmount} & $[0$, $+30]$ & Rotate amounts, must be multiples of 2 \\
      \hline
      Unsigned 4-bit & 4 & \texttt{ARM\_Const4\_CoprocessorOpcode} & $[0$, $+15]$ & Coprocessor opcodes \\
      \hline
      Unsigned 3-bit & 3 & \texttt{ARM\_Const3\_CoprocessorOpcode} & $[0$, $+7]$ & Coprocessor opcodes \\
      \hhline{|=|=|=|=|=|}
      Unsigned 9-bit & 9 & \texttt{ARM\_Const9\_Unsigned} & $[0$, $+511]$ & Must be multiples of 4 \\
      \hline
      Unsigned 7-bit & 7 & \texttt{ARM\_Const7\_Unsigned} & $[0$, $+127]$ & Must be multiples of 4 \\
      \hline
      Unsigned 6-bit & 6 & \texttt{ARM\_Const6\_Unsigned} & $[0$, $+63]$ & Must be multiples of 2 \\
      \hline
      Unsigned 3-bit & 3 & \texttt{ARM\_Const3\_Unsigned} & $[0$, $+7]$ & \\
      \hhline{|=|=|=|=|=|}
      Unsigned 16-bit & 16 & \texttt{ARM\_Const16\_Unsigned} & $[0$, $+2^{16} - 1]$ & \\
      \hhline{|=|=|=|=|=|}
      Unsigned 5-bit & 5 & \texttt{ARM\_Const5\_SatPos} & $[+1$, $+16]$ & Saturation positions \\
      \hline
    \end{tabular}
  \end{center}
\end{table}

All \gls{WIR} immediate classes used in the context of the ARM processor models are summarized in \cref{tab:armimmtypes}.
For the sake of brevity, this user's guide will not discuss all the various classes and their rationale in detail, the interested reader is referred to~\cite{arm05} instead.
Instead, only a few comments on very particular immediate classes will be given in the following.

\begin{table}[!b]
  \begin{center}
    \caption{Opcode Types of \texttt{ARMv4} Processor Model}
    \label{tab:armv4opctypes}
    \begin{tabular}{|c|p{95mm}|p{42mm}|}
      \hline
      \multirow{2}{*}[-1mm]{\textbf{\shortstack[c]{Opcode\\ Type}}} & \multirow{2}{*}{\textbf{Opcodes}} & \multirow{2}{*}{\textbf{Comment}} \\
      & & \\
      \hhline{|=|=|=|}
      \texttt{load} & \texttt{LDC}, \texttt{LDM}, \texttt{LDR}, \texttt{LDRB}, \texttt{LDRBT}, \texttt{LDRH}, \texttt{LDRSB}, \texttt{LDRSH}, \texttt{LDRT}, \texttt{SWP}, \texttt{SWPB} & \\
      \hline
      \texttt{store} & \texttt{STC}, \texttt{STM}, \texttt{STR}, \texttt{STRB}, \texttt{STRBT}, \texttt{STRH}, \texttt{STRT}, \texttt{SWP}, \texttt{SWPB} & \\
      \hline
      \texttt{call} & \texttt{BL} & \\
      \hline
      \texttt{indirectCall} & \texttt{SWI} & \\
      \hline
      \texttt{condJump} & \texttt{B} & \\
      \hline
      \texttt{uncondJump} & \texttt{B} & \\
      \hline
    \end{tabular}
  \end{center}
\end{table}

The ARMv4 \gls{ISA} inherently supports to communicate with a coprocessor, requiring the specification of coprocessor offsets and opcodes.
These elements are modeled by classes \texttt{ARM\_Const10\_CoprocessorOffset}, \texttt{ARM\_Const4\_CoprocessorOpcode} and \texttt{ARM\_Const3\_CoprocessorOffset}, where such coprocessor offsets have to be multiples of 4.
Class \texttt{ARM\_Const6\_Unsigned0} is designed to cover the value range from 1 to 32, excluding the value 0.
Similarly, class \texttt{ARM\_Const5\_Unsigned0} ranges from 1 to 31, excluding 0.
Some addressing modes support the automatic rotation of some argument, which requires a rotate amount as parameter.
Such rotate amounts are realized by class \texttt{ARM\_Const5\_RotateAmount} which covers the value range from 0 to 30, and rotate amounts are enforced to be even numbers.

The \texttt{ARMv4T} model introduces 16-bit thumb operations which again require additional immediate types as can be seen from the middle part of \cref{tab:armimmtypes}.
New immediates with widths of 9, 7, 6 and 3 bits are realized by the according classes, resp.
Actual values for \texttt{ARM\_Const9\_Unsigned} and \texttt{ARM\_Const7\_Unsigned} must be multiples of 4, and \texttt{ARM\_Const6\_Unsigned} only accepts even values.

\enlargethispage{5mm}

The \texttt{ARMv5T} processor model adds yet another class \texttt{ARM\_Const16\_Unsigned}, and immediate type \texttt{ARM\_\-Const5\_SatPos} is introduced by the \texttt{ARMv6} processor model and reflects saturation positions ranging from 1 to 16 (cf. bottom of \cref{tab:armimmtypes}).

Code examples for all ARM immediate classes can be found in unit tests \texttt{WIR/arch/arm/tests/armConst*}.

\begin{table}[t]
  \begin{center}
    \caption{Opcode Types of \texttt{ARMv4T}, \texttt{ARMv5T}, \texttt{ARMv5TE}, \texttt{ARMv5TEJ} and \texttt{ARMv6} Processor Models}
    \label{tab:armv4topctypes}
    \begin{tabular}{|c|p{20mm}|p{20mm}|p{20mm}|p{20mm}|p{20mm}|p{20mm}|}
      \hline
      \multirow{2}{*}[-1mm]{\textbf{\shortstack[c]{Opcode\\ Type}}} & \multicolumn{5}{c|}{\textbf{Opcodes}} & \multirow{2}{*}{\textbf{Comment}} \\
      & \texttt{ARMv4T} & \texttt{ARMv5T} & \texttt{ARMv5TE} & \texttt{ARMv5TEJ} & \texttt{ARMv6} & \\
      \hhline{|=|=|=|=|=|=|=|}
      \texttt{load} & \texttt{LDMIA}, \texttt{POP} & \texttt{LDC2} & \texttt{LDRD} & & \texttt{LDREX} & \\
      \hline
      \texttt{store} & \texttt{PUSH}, \texttt{STMIA} & \texttt{STC2} & \texttt{STRD} & & \texttt{SRS}, \texttt{STREX} & \\
      \hline
      \texttt{move} & & & & & \texttt{CPY} & \\
      \hline
      \texttt{call} & & \texttt{BLX} & & & & \\
      \hline
      \texttt{ret} & & & & & \texttt{RFE} & \\
      \hline
      \texttt{indirectJump} & \texttt{BX} & \texttt{BKPT} & & \texttt{BXJ} & & \\
      \hline
    \end{tabular}
  \end{center}
\end{table}

The \texttt{ARMv4} model specifies all 26 opcodes of the base \gls{ISA}, beginning with \texttt{ADC} and ranging to \texttt{UMULL}.
As can be seen from \cref{tab:armv4opctypes}, many of these opcodes are characterized within \gls{WIR} either as memory load or store operations.
Three more opcodes are furthermore modeled to be branches or function calls, resp.
None of the opcodes modeled within \gls{WIR} is a pseudo-opcode.

\cref{tab:armv4topctypes} summarizes the opcode characterizations of all the ARM processor models derived from the base model \texttt{ARMv4}.
This table again shows that most of these new opcodes are characterized as loads or stores.
A couple of other opcodes being register-register moves, function calls and returns, and indirect branches are introduced by the derived processor models as well.

Due to the complex different encodings of ARM operations, one opcode can play different roles when it comes to its characterization within \gls{WIR}.
For this reason, it might not be appropriate to simply classify, e.g., opcode \texttt{B} as unconditional branch, since it might also be a conditional branch depending on its condition field.
For this reason, the \texttt{ARMv4} model contains dedicated overloaded methods for opcode handling as shown in \cref{lst:armopcodes}.
These methods inspect an opcode in its specific context of some \gls{WIR} operation and return whether the opcode definitely realizes a register-register move, an indirect call, a function return, a conditional, an unconditional or an indirect jump, resp.

The complexity of the ARM \gls{ISA} can be seen from the fact that a total of 212 different operation formats are required within \gls{WIR} to capture all possible encodings resulting from the only 26 opcodes of base class \texttt{ARMv4}.
The majority of 128 of all these formats are solely required to model the two opcodes \texttt{LDM} and \texttt{STM}.
These two opcodes allow to load or store multiple registers in one machine operation, where the list of registers to be loaded or stored can contain 1 up to 16 registers.
For the proper data flow and def-use analysis within \gls{WIR}, all possible lengths of such register lists need to be made explicit, causing this observed effect.
The \gls{WIR} processor models \texttt{ARMv4T}, \texttt{ARMv5T}, \texttt{ARMv5TE} and \texttt{ARMv6} introduce additional 81, 10, 16 and 24 operation formats, resp.
The naming scheme used for all ARM operation formats is summarized in \cref{tab:armv4opf} and follows the usual approach already introduced in the previous sections.

\begin{lstlisting}[style=code,float=!t,caption={Specific Opcode Handling of Class \texttt{ARMv4}},label=lst:armopcodes]
class ARMv4 : public WIR_Processor<ARMv4>
{
  public:

    [...]

    class OpCode : public WIR_Processor::OpCode
    {
      public:

        [...]

        // Determine whether an ARM opcode is a register-register move.
        virtual @bool@ isMove@(@ const WIR_Operation æ&æo @)@ const;Ü\label{lst:armopcodes1}Ü

        // Determine whether an ARM opcode indirectly calls a function.
        virtual @bool@ isIndirectCall@(@ const WIR_Operation æ&æo @)@ const;Ü\label{lst:armopcodes2}Ü

        // Determine whether an ARM opcode returns from a function call.
        virtual @bool@ isReturn@(@ const WIR_Operation æ&æo @)@ const;Ü\label{lst:armopcodes3}Ü

        // Determine whether an ARM opcode performs a conditional jump.
        virtual @bool@ isConditionalJump@(@ const WIR_Operation æ&æo @)@ const;Ü\label{lst:armopcodes4}Ü

        // Determine whether an ARM opcode performs an unconditional jump.
        virtual @bool@ isUnconditionalJump@(@ const WIR_Operation æ&æo @)@ const;Ü\label{lst:armopcodes5}Ü

        // Determine whether an ARM opcode indirectly performs a jump.
        virtual @bool@ isIndirectJump@(@ const WIR_Operation æ&æo @)@ const;Ü\label{lst:armopcodes6}Ü

        [...]
    };
};
\end{lstlisting}

Unit tests \texttt{WIR/arch/arm/tests/arm*Operations*} illustrate the construction and use all of modeled ARM opcodes in all existing formats.
As for the previous processor models, these unit tests once create and validate all opcodes and their properties, and once print them out and feed them through the GNU ARM assembler.

The \gls{WIR} processor models feature all addressing modes defined by the ARM \gls{ISA}.
This covers, e.g., all variants of shifts and rotations, address additions or subtractions, pre- or post-indexed addressing, increment or decrement before and after, full ascending/descending etc., as well as all coprocessor addressing modes.
In total, the \texttt{ARMv4} processor model features 49 different addressing modes whose implementation can be found in file \texttt{armaddressingmodes.cc}.

The ARM architecture furthermore defines a total of 15 condition codes that are used for the predicated execution of machine operations.
All these condition codes are properly modeled within \gls{WIR}, ranging from ``\texttt{eq}'' that stands for a test on equality, up to ``\texttt{al}'' that denotes to execute an operation always.
The realization of all these \texttt{ARMv4} condition codes can be found in file \texttt{armconditions.cc}.

\begin{table}[ht]
  \begin{center}
    \caption{Naming Scheme for Operation Formats of \texttt{ARMv4}, \texttt{ARMv4T}, \texttt{ARMv5T}, \texttt{ARMv5TE} and \texttt{ARMv6} Processor Models}
    \label{tab:armv4opf}
    \begin{tabular}{|c|p{136mm}|}
      \hline
      \multirow{2}{*}[-1mm]{\textbf{\shortstack[c]{Operation Format\\ Specifier}}} & \multirow{2}{*}{\textbf{Parameter Type}} \\
      & \\
      \hhline{|=|=|}
      ``\texttt{A}'' & Addressing mode parameter \\
      \hline
      ``\texttt{C}'' & Condition field parameter \\
      \hline
      ``\texttt{C24}'' & Unsigned 24-bit immediate constant \\
      \hline
      ``\texttt{C12}'' & Unsigned 12-bit immediate constant \\
      \hline
      ``\texttt{C10}'' & Unsigned 10-bit coprocessor offset \\
      \hline
      ``\texttt{C8}'' & Unsigned 8-bit immediate constant \\
      \hline
      ``\texttt{C60}'' & Unsigned 6-bit immediate constant without 0 \\
      \hline
      ``\texttt{C5}'' & Unsigned 5-bit immediate constant \\
      \hline
      ``\texttt{C50}'' & Unsigned 5-bit immediate constant without 0 \\
      \hline
      ``\texttt{L}'' & Label parameter \\
      \hline
      ``\texttt{R}'' & Register parameter of type \texttt{ARMv4::RegisterType::reg} \\
      \hline
      ``\texttt{R\emph{n}}'' & A series of \texttt{\emph{n}} register parameters of type \texttt{ARMv4::RegisterType::reg} \\
      \hline
      ``\texttt{RA}'' & Unsigned 5-bit rotate amount \\
      \hline
      ``\texttt{O}'' & Unsigned 3-bit or 4-bit coprocessor opcode \\
      \hline
      ``\texttt{S}'' & Unstructured string parameter \\
      \hhline{|=|=|}
      ``\texttt{C10}'' & Unsigned 10-bit coprocessor offset \\
      \hline
      ``\texttt{C9}'' & Unsigned 9-bit immediate constant \\
      \hline
      ``\texttt{C7}'' & Unsigned 7-bit immediate constant \\
      \hline
      ``\texttt{C3}'' & Unsigned 3-bit immediate constant \\
      \hline
      ``\texttt{LR}'' & Physical link register parameter \texttt{R14} \\
      \hline
      ``\texttt{PC}'' & Physical program counter register parameter \texttt{R15} \\
      \hline
      ``\texttt{SP}'' & Physical stack pointer register parameter \texttt{R13} \\
      \hline
      ``\texttt{T\dots}'' & Prefix for 16-bit thumb operation formats \\
      \hhline{|=|=|}
      ``\texttt{C16}'' & Unsigned 16-bit immediate constant \\
      \hhline{|=|=|}
      ``\texttt{P}'' & Extended register parameter of type \texttt{ARMv5TE::RegisterType::pReg} \\
      \hhline{|=|=|}
      ``\texttt{C5SP}'' & Unsigned 5-bit saturation position \\
      \hline
    \end{tabular}
  \end{center}
\end{table}

How to deal with the ARM-specific condition field parameters can be seen in unit test \texttt{WIR/arch/arm/tests/} \texttt{armConditionFieldParameter1.cc}.


%-------------------------------------------------------------------------
\section{Physical Memory Regions}
\label{sec:memoryregions}
%-------------------------------------------------------------------------

While processor models as described in the previous \cref{sec:processors} play a critical role within \gls{WIR}, they are by far not the only components of modern computing systems that require careful modelling.
Due to the fact that physical memories have a direct impact on an entire system's performance, it is crucial to also consider the properties of memory regions in \gls{WIR} system models.
This is precisely the purpose of the class \texttt{WIR\_MemoryRegion} described in the following.
\gls{WIR} uses the notion of memory regions instead of entire memories in order to enable the partitioning of complete memories into different chunks that might carry different attributes.
Each such chunk of a full memory is then represented by a dedicated \gls{WIR} memory region.

\begin{lstlisting}[style=code,float=ht,caption={Name, Address and Size Handling of Class \texttt{WIR\_MemoryRegion}},label=lst:regaddr]
// Get the region name.
ästd::stringä WIR_MemoryRegion::getName@( void )@ const;Ü\label{lst:regaddr1}Ü

// Get a region's base address.
WIR_MemoryAddress WIR_MemoryRegion::getBaseAddress@( void )@ const;Ü\label{lst:regaddr2}Ü

// Get a region's length in bytes.
ästd::size_tä WIR_MemoryRegion::getLength@( void )@ const;Ü\label{lst:regaddr3}Ü
\end{lstlisting}

The most fundamental attributes of a region are its name, start address and size.
They can be retrieved using the methods \texttt{getName}, \texttt{getBaseAddress} and \texttt{getLength} shown in \crefrange{lst:regaddr1}{lst:regaddr3} of \cref{lst:regaddr}.

\begin{lstlisting}[style=code,float=!b,caption={Section Handling of Class \texttt{WIR\_MemoryRegion}},label=lst:regsec]
// Get the set of all sections attached to a region.
const WIR_SectionSet æ&æWIR_MemoryRegion::getSections@( void )@ const;Ü\label{lst:regsec1}Ü

// Determine whether a region contains a given section.
@bool@ WIR_MemoryRegion::containsSection@(@ const WIR_Section æ&æ @)@ const;Ü\label{lst:regsec2}Ü

// Get a region's free space in bytes.
ästd::size_tä WIR_MemoryRegion::getFreeSpace@( void )@;Ü\label{lst:regsec3}Ü
\end{lstlisting}

As already introduced in \cref{sssec:sections}, logical sections of an executable are attached to a physical memory region (cf. \cref{lst:sctconstruction}).
Regions thus act as containers for sections within \gls{WIR}.
The information which sections are actually associated with a region can be retrieved using method \texttt{getSections} shown in \cref{lst:regsec1} of \cref{lst:regsec} which returns a reference to a section set.
Method \texttt{containsSection} (\cref{lst:regsec2}) can be used to determine whether some logical section is actually associated with a memory region or not.
Furthermore, \texttt{getFreeSpace} (\cref{lst:regsec3}) computes how many bytes are still free and unoccupied within a region.
This free space depends on the sizes of all the sections associated with a region, their respective contents and alignments, etc.

\begin{lstlisting}[style=code,float=!t,caption={Characteristic Attributes of Class \texttt{WIR\_MemoryRegion}},label=lst:regattribs]
// Determine whether a region is dynamically allocated during runtime.
@bool@ WIR_MemoryRegion::isDynamicallyAllocated@( void )@ const;Ü\label{lst:regattribs1}Ü

// Get a bit-field encoding memory region attributes.
@unsigned long@ WIR_MemoryRegion::getAttributes@( void )@ const;Ü\label{lst:regattribs2}Ü

// Get the name of the memory type that a region belongs to.
ästd::stringä WIR_MemoryRegion::getMemoryType@( void )@ const;Ü\label{lst:regattribs3}Ü

// Get the ratio of the memory region's clock frequency compared to the processor's clock frequency.
@long double@ WIR_MemoryRegion::getClockRatio@( void )@ const;Ü\label{lst:regattribs4}Ü

// Get the minimum number of clock cycles for which a region can delay a memory access.
@unsigned int@ WIR_MemoryRegion::getMinDelay@( void )@ const;Ü\label{lst:regattribs5}Ü

// Get the maximum number of clock cycles for which a region can delay a memory access.
virtual @unsigned int@ WIR_MemoryRegion::getMaxDelay@(@ const @unsigned int@ d @=@ §0§ @)@ const;Ü\label{lst:regattribs6}Ü

// Get the minimum burst access delay of a region in clock cycles.
@unsigned long@ WIR_MemoryRegion::getMinBurstDelay@( void )@ const;Ü\label{lst:regattribs7}Ü

// Get the maximum burst access delay of a region in clock cycles.
@unsigned long@ WIR_MemoryRegion::getMaxBurstDelay@( void )@ const;Ü\label{lst:regattribs8}Ü

// Get the byte width of the memory bus that connects a region.
@unsigned long@ WIR_MemoryRegion::getBusWidth@( void )@ const;Ü\label{lst:regattribs9}Ü

// Get the hierarchy of all system components residing in front of this memory region.
const ästd::setäæ<æästd::listäæ<æästd::reference_wrapperäæ<æWIR_SystemComponentæ>>> &æWIR_MemoryRegion::getHierarchy@( void )@ const;Ü\label{lst:regattribs10}Ü

// Enum WIR_MemoryRegionAttributes represents different attributes of WIR memory regions.Ü\label{lst:regattribs11}Ü
enum class WIR_MemoryRegionAttributes : @unsigned char@
{
  none = §0§,                   // A memory region without any specific attributes.
  read = §1§ << §0§,              // A memory region that can be read.
  write = §1§ << §1§,             // A memory region that can be written.
  execute = §1§ << §2§,           // A memory region that can be executed.
  allocated = §1§ << §3§,         // A memory region that is allocated.
  initialized = §1§ << §4§,       // A memory region that is initialized.
  cached = §1§ << §5§,            // A memory region that is cached.
  all = ( cached << §1§ ) - §1§   // Covers all existing attributes.
};Ü\label{lst:regattribs12}Ü
\end{lstlisting}

Within \gls{WIR}, a physical memory region carries a couple of additional properties which can be queried using the methods shown in \cref{lst:regattribs}.
Firstly, method \texttt{isDynamicallyAllocated} (cf. \cref{lst:regattribs1}) can be used to determine whether the amount of used space within a memory region may vary dynamically during runtime.
This can, e.g., happen when a system's runtime stack is mapped to some region.
Depending on whether a region's contents is read-only or executable etc., it is furthermore characterized by additional attributes that model its ``access policy''.
These attributes can be obtained via method \texttt{getAttributes} (\cref{lst:regattribs2}).
This method returns an \texttt{unsigned long} value that forms a bit-field whose individual bits represent specific access attributes.
These bit-level attributes are modeled by the \gls{WIR} enumeration type \texttt{WIR\_MemoryRegionAttributes} shown in \crefrange{lst:regattribs11}{lst:regattribs12}.
Besides their very own names, regions can also be given a human-readable name standing for the memory type on which a region bases like, e.g., ``FLASH'' or ``ROM''.
This memory type name can be retrieved via method \texttt{getMemoryType} (\cref{lst:regattribs3}).

The next couple of methods is useful to obtain performance figures related to a memory region.
First, method \texttt{getClockRatio} (\cref{lst:regattribs4}) returns a double-precision floating-point value that characterizes the relationship between the region's and the processor core's clock frequencies.
\texttt{getMinDelay} and \texttt{getMaxDelay} (\cref{lst:regattribs5,lst:regattribs6}) return a region's minimum and maximum access delays in clock cycles, resp.
Argument \texttt{d} of \texttt{getMaxDelay} intends to represent the additional delay induced by other components in a system's memory hierarchy.
Since a memory region by construction is the very last component in a memory hierarchy, this argument is simply ignored here.
Memories usually exhibit different access latencies depending on whether single accesses or burst accesses are considered.
Methods \texttt{getMinDelay} and \texttt{getMaxDelay} provide the latencies for one single access.
Accordingly, methods \texttt{getMinBurstDelay} and \texttt{getMaxBurstDelay} (\cref{lst:regattribs7,lst:regattribs8}) provide the minimum and maximum access delays for bursts, resp.

\gls{WIR} allows to model systems with complex memory hierarchies consisting of buses (see \cref{sec:buses}), various levels of caches (cf. \cref{sec:caches}), and memories of different types.
When considering the position of a memory region in such a hierarchy, the bit width of the interface that connects the region to its bus is yet another important attribute.
This bit width can be obtained via method \texttt{getBusWidth} from \cref{lst:regattribs9}.
Finally, method \texttt{getHierarchy} (\cref{lst:regattribs10}) can be used to obtain the full hierarchy of system components that reside ``in front of'' a region.
Here, ``in front'' refers to all components that are closer to a processor core than the considered region itself.
\texttt{getHierarchy} returns a set of lists of such system components.
Each of these lists only contains buses or caches, and the component at a list's head is closest to the processor core, while the component at a list's tail is the farthest away from the processor core, i.e., the closest to the currently considered memory region.
In real-life system architectures, there can be multiple paths from a processor core through the memory hierarchy to one and the same physical memory region.
This can happen if, e.g., a certain part of a region's address space is covered by a cache, while other parts of the same region's address space are uncached.
For this reason, method \texttt{getHierarchy} returns a set potentially containing several of these lists described earlier.
Each list within the set returned by \texttt{getHierarchy} represents exactly one unique path through a system's memory hierarchy from a processor core to the region ifself.


%--------------------------------------------------------------------------
\section{Caches}
\label{sec:caches}
%-------------------------------------------------------------------------

Modern processor architectures feature sophisticated memory hierarchies that not only consist of physical memory regions as described in the previous section, but that also use caches.
Thus, class \texttt{WIR\_Cache} is derived from base class \texttt{WIR\_SystemComponent} and provides the representation of various flavors of caches within the \gls{WIR} library.

\begin{lstlisting}[style=code,float=t,caption={Architectural Attributes of Class \texttt{WIR\_Cache}},label=lst:cachearch]
// Determine whether a cache is enabled and active.
@bool@ WIR_Cache::isEnabled@( void )@ const;Ü\label{lst:cachearch1}Ü

// Get a cache's type.
WIR_Cache::CacheType WIR_Cache::getCacheType@( void )@ const;Ü\label{lst:cachearch2}Ü

// Get a cache's position in a system's memory hierarchy.
WIR_Cache::CacheLevel WIR_Cache::getLevel@( void )@ const;Ü\label{lst:cachearch3}Ü

// Determine whether a cache is shared between multiple processor cores.
@bool@ WIR_Cache::isShared@( void )@ const;Ü\label{lst:cachearch4}Ü

// Get the time in clock cycles that a cache hit needs to complete in the best case.
@unsigned int@ WIR_Cache::getHitDelay@( void )@ const;Ü\label{lst:cachearch5}Ü

// Get the time in clock cycles that a cache miss needs to complete in the worst case.
@unsigned int@ WIR_Cache::getMissDelay@( void )@ const;Ü\label{lst:cachearch6}Ü

// Get the maximum number of clock cycles for which a cache can delay a memory access.
virtual @unsigned int@ WIR_Cache::getMaxDelay@(@ const @unsigned int@ d @)@ const;Ü\label{lst:cachearch7}Ü

// Get the byte width of the memory bus that connects a cache.
@unsigned long@ WIR_Cache::getBusWidth@( void )@ const;Ü\label{lst:cachearch8}Ü

// Enum CacheType represents the different types a cache can have.
enum class WIR_Cache::CacheType : @char@Ü\label{lst:cachearch9}Ü
{
  I,            // An instruction cache.
  D,            // A data cache.
  U             // A unified cache.
};

// Enum CacheLevel represents a cache's position in the memory hierarchy.
enum class WIR_Cache::CacheLevel : @char@Ü\label{lst:cachearch10}Ü
{
  NOTSET @=@ §0§,   // A cache whose level in the memory hierarchy has not been set.
  L1,           // A level-1 cache.
  L2            // A level-2 cache.
};
\end{lstlisting}

The interface of class \texttt{WIR\_Cache} to query all system-level architectural attributes of a cache without considering its internal organization is shown in \cref{lst:cachearch}.
Caches can regularly be enabled or disabled by setting system-specific status bits.
The current state of a cache can be retrieved with method \texttt{isEnabled} (cf. \cref{lst:cachearch1}).
Caches are regularly classified as instruction, data or unified caches.
This classification is modeled by the enumeration class \texttt{CacheType} (cf. \cref{lst:cachearch9}), and a cache's type can be obtained using method \texttt{getCacheType} (\cref{lst:cachearch2}).
The position of a cache within a system's memory hierarchy, i.e., the cache's level, is modeled by enumeration \texttt{CacheLevel} (\cref{lst:cachearch10}), and it can be determined using method \texttt{getLevel} (\cref{lst:cachearch3}).
In multi-core systems, a cache can be shared between several processor cores, and this information is accessible via method \texttt{isShared} (\cref{lst:cachearch4}).

The best-case delay in clock cycles that a cache can cause for a cache hit can be obtained with method \texttt{getHit\-Delay} (\cref{lst:cachearch5}).
Similarly, \texttt{getMissDelay} (\cref{lst:cachearch6}) returns the worst-case delay in clock cycles for the case of a cache miss.
Note that this miss delay only refers to the extra time required by the cache itself, excluding the delays caused by accesses to parts of the memory hierarchy lying behind the missing cache.
The maximal delay of a memory cache that can be caused by a cache including that of subsequent parts of the memory hierarchy can be obtained with the virtual method \texttt{getMaxDelay} (\cref{lst:cachearch7}).
It expects an unsigned integer argument \texttt{d} that represents the access delay caused by the tail of the memory hierarchy.
\texttt{getMaxDelay} assumes a worst-case scenario where an evicted cache line needs to be written back and a new cache line is loaded.
It computes how many bus transfers are required for these two transactions, multiplies this by delay \texttt{d}, and finally adds the cache's very own intrinsic miss delay.
This default implementation of \texttt{getMaxDelay} can be overridden for specialized classes inheriting from \texttt{WIR\_Cache}.
The number of bus transfers initiated by a cache naturally depends on the width of the interface connecting the cache with a system's bus.
This bus width in bytes can be retrieved using method \texttt{getBusWidth} (\cref{lst:cachearch8}).

\begin{lstlisting}[style=code,float=t,caption={Geometry-Related Attributes of Class \texttt{WIR\_Cache}},label=lst:cachegeom]
// Get the overall byte size of the entire data memory of a cache.
@unsigned int@ WIR_Cache::getSize@( void )@ const;Ü\label{lst:cachegeom1}Ü

// Get a cache's associativity.
@unsigned int@ WIR_Cache::getAssociativity@( void )@ const;Ü\label{lst:cachegeom2}Ü

// Get the number of sets (i.e., rows) of a cache.
@unsigned int@ WIR_Cache::getNumberOfSets@( void )@ const;Ü\label{lst:cachegeom3}Ü

// Get a cache's line size in bytes.
@unsigned int@ WIR_Cache::getLineSize@( void )@ const;Ü\label{lst:cachegeom4}Ü

// Determine whether a cache's size is given as percentage, and not as an absolute byte size.
@bool@ WIR_Cache::isSizedInPercent@( void )@ const;Ü\label{lst:cachegeom5}Ü

// Set a cache's size to a percentage of a given task size.
@void@ WIR_Cache::setRelativeSize@( unsigned int@ t @)@;Ü\label{lst:cachegeom6}Ü

// Get the number of offset bits that memory addresses require for this cache.
@int@ WIR_Cache::getOffsetBits@( void )@ const;Ü\label{lst:cachegeom7}Ü

// Get the number of index bits that memory addresses require for this cache.
@int@ WIR_Cache::getIndexBits@( void )@ const;Ü\label{lst:cachegeom8}Ü

// Get the number of tag bits that memory addresses require for this cache.
@int@ WIR_Cache::getTagBits@( void )@ const;Ü\label{lst:cachegeom9}Ü

// Determine whether a cache implements the write-through policy. If not, the cache is a write-back cache.
@bool@ WIR_Cache::isWriteThrough@( void )@ const;Ü\label{lst:cachegeom10}Ü

// Determine whether the cache allocates on writes.
@bool@ WIR_Cache::isWriteAllocate@( void )@ const;Ü\label{lst:cachegeom11}Ü
\end{lstlisting}

All attributes related to the internal geometry of a cache can be retrieved with the methods shown in \cref{lst:cachegeom}.
Most relevantly, this includes the sizes of a cache and its numbers of rows and columns.
Thus, the overall byte size of a cache can be obtained using method \texttt{getSize} (cf. \cref{lst:cachegeom1}).
The returned size refers solely to the cache's data memory, i.e., that memory storing its payload in the form of cached code or data memory lines.
\texttt{getSize} does not consider the size of the tag memory that a cache requires for its hit/miss checks.
The associativity of a cache can be queried via method \texttt{getAssociativity} (\cref{lst:cachegeom2}).
Typical values might range from 1 (direct-mapped cache) to 8 (set-associative cache).
As \texttt{getAssociativity} in other words returns the number of internal columns of a cache, method \texttt{getNumberOfSets} (\cref{lst:cachegeom3}) provides the total number of sets, i.e., rows that a cache features.
Class \texttt{WIR\_Cache} enforces that a cache's size, associativity and its number of sets are all powers of 2.
Each single element inside a cache in a certain row and column is called a cache line.
The byte size of such a cache line can be retrieved by method \texttt{getLineSize} (\cref{lst:cachegeom4}).

Based on the above geometrical attributes, class \texttt{WIR\_Cache} ensures that the following relationship always holds:
\begin{displaymath}
  \mathit{Size}_\mathit{Cache} := \mathit{Associativity} \ast \mathit{\#Sets} \ast \mathit{Size}_\mathit{Line}
\end{displaymath}
By multiplying associativity and the number of available sets, the dimension of a cache's internal memory matrix is determined.
The final multiplication by a single element's byte size then yiels the absolute byte size of the full cache (excluding tag memories).

Besides specifying a cache's size as absolute byte size according to some real processor architecture, it might occasionally be useful to characterize a cache size in a relative fashion, as a percentage of the size of some software task executed on a system.
Typical use cases for such relative cache sizes are, e.g., benchmarking campaigns that evaluate a cache's performance depending on the sizes of the used benchmarks.
Whether or not the size of a \gls{WIR} cache is given as such a percentage can be checked with method \texttt{isSizedInPercent} (cf. \cref{lst:cachegeom5} of \cref{lst:cachegeom}).
The absolute byte size of such a benchmark task which is considered to be 100\% can be specified using method \texttt{setRelativeSize} (\cref{lst:cachegeom6}).

In order to locate an entry at some given memory address inside a cache, its row and column positions within the cache's internal matrix need to be determined.
For this purpose, the bits of a memory address are partitioned into its offset (least-significant part), index (middle part), and tag (most-significant part).
How many bits are actually needed for a cache with a given geometry can be determined with the methods \texttt{getOffsetBits}, \texttt{getIndexBits} and \texttt{getTagBits} (cf. \crefrange{lst:cachegeom7}{lst:cachegeom9}).
Here, the \gls{WIR} library always adheres to the following relationships:
\begin{displaymath}
  \begin{array}{rcl}
    \mathit{Size}_\mathit{Line} & := & 2^\mathit{\#Offset Bits} \\
    \mathit{\#Sets} & := & 2^\mathit{\#Index Bits} \\
    \mathit{\#Tag Bits} & := & \mathit{AddrWidth} - \mathit{\#Index Bits} - \mathit{\#Offset Bits}
  \end{array}
\end{displaymath}
where $\mathit{AddrWidth}$ denotes the bit width of a \gls{WIR} memory address.

Besides looking up an entry inside a cache, various policies exist whenever a cache entry is overwritten.
The most simple write policy is write-through, where the overwritten cache line is directly transferred to the rear parts of the memory hierarchy coming after the involved cache.
The write-back policy in contrast delays such costly transfers as long as possible.
Which of these two write policies is implemented by a \gls{WIR} cache can be checked using method \texttt{isWriteThrough} (\cref{lst:cachegeom10}).
This method returns \texttt{true} for write-through caches and \texttt{false} for write-back caches.
In the case of a write miss inside a cache, the write allocate or the no-write allocate policies come into play.
The former allocates space in the missing cache for the new data to be written, while the latter simply forwards the new data to the tail of the memory hierarchy without changing the cache's internal state.
Method \texttt{isWriteAllocate} (\cref{lst:cachegeom11}) finally returns \texttt{true} if write allocate is active and \texttt{false} if a cache implements no-write allocate.

The functional behavior of caches heavily depends on memory addresses and address ranges that serve as input to a cache.
As already described on \cpageref{page:memaddr}, addresses are represented by type \texttt{WIR\_MemoryAddress}.
Based on this type, address ranges are modeled within \gls{WIR} as intervals of addresses having a well-defined lower and upper bound.
Thus, type \texttt{WIR\_AddressRange} is defined as \texttt{WIR\_Interval<WIR\_MemoryAddress>}.
The \gls{WIR}-internal template class \texttt{WIR\_Interval} (cf. file \texttt{analyses/generic/wirinterval.h}) provides the simple construction of and access to such intervals with lower and upper bounds, as well as all standard arithmetical and comparison operations for intervals.

In order to support static analyses that aim to estimate how \gls{WIR} code consisting of sequences of operations changes a cache's internal state, it is necessary to relate addresses of memory accesses issued by such operations to the proper locations inside the cache which might be affected.
For this purpose, class \texttt{WIR\_Cache} provides the \gls{API} depicted in \cref{lst:cacheaddr}.

\begin{lstlisting}[style=code,float=t,caption={Address and Address Range Handling of Class \texttt{WIR\_Cache}},label=lst:cacheaddr]
// Local type definition to represent cache lines (tag + index).
using WIR_Cache::CacheLineID = WIR_MemoryAddress;Ü\label{lst:cacheaddr1}Ü

// Local type definition to represent cache index bits.
using WIR_Cache::CacheIndex = WIR_MemoryAddress;Ü\label{lst:cacheaddr2}Ü

// Local type definition to represent cache tag bits.
using WIR_Cache::CacheTag = WIR_MemoryAddress;Ü\label{lst:cacheaddr3}Ü

// Get the cache line's ID (tag and index bits) targeted by an access to a given address.
WIR_Cache::CacheLineID WIR_Cache::getLineID@(@ const WIR_MemoryAddress æ&æa @)@ const;Ü\label{lst:cacheaddr4}Ü

// Get the index bits from a given memory address.
WIR_Cache::CacheIndex WIR_Cache::getIndex@(@ const WIR_MemoryAddress æ&æa @)@ const;Ü\label{lst:cacheaddr5}Ü

// Get the tag bits from a given memory address.
WIR_Cache::CacheTag WIR_Cache::getTag@(@ const WIR_MemoryAddress æ&æa @)@ const;Ü\label{lst:cacheaddr6}Ü

// Get the range of cache line IDs that may be affected by an access to a given address range.
WIR_Interval<WIR_Cache::CacheLineID> WIR_Cache::getAffectedLineIDs@(@ const WIR_AddressRange æ&ær @)@ const;Ü\label{lst:cacheaddr7}Ü

// Get the range of cache sets that may be affected by an access to a given address range.
WIR_Interval<WIR_Cache::CacheIndex> WIR_Cache::getAffectedSets@(@ const WIR_AddressRange æ&ær @)@ const;Ü\label{lst:cacheaddr8}Ü

// Get the range of cache tags that may be affected by an access to a given address range.
WIR_Interval<WIR_Cache::CacheTag> WIR_Cache::getAffectedTags@(@ const WIR_AddressRange æ&ær @)@ const;Ü\label{lst:cacheaddr9}Ü
\end{lstlisting}

Methods \texttt{getLineID}, \texttt{getIndex} and \texttt{getTag} (cf. \crefrange{lst:cacheaddr4}{lst:cacheaddr6}) return which cache line, set and tag is affected by a given memory address, resp.
As described earlier, the internal geometry of a cache determines which and how many bits of an address represent tags, indices and offsets.
The aforementioned methods then extract the tag and index bits from the given address and basically return these bits.
For the sake of type safeness, class \texttt{WIR\_Cache} defines its own data types representing cache line IDs, index and tag bits.
These local types finally simply boil down to \gls{WIR} addresses as can be seen from \crefrange{lst:cacheaddr1}{lst:cacheaddr3} of \cref{lst:cacheaddr}.

During static cache analysis, not only the effect of one single address on a cache's internal state needs to be estimated, but also that of entire address ranges.
Consequently, methods \texttt{getAffectedLineIDs}, \texttt{getAffected\-Sets} and \texttt{getAffectedTags} (cf. \crefrange{lst:cacheaddr7}{lst:cacheaddr9}) get an interval \texttt{r} of addresses as argument.
They then determine which ranges of cache lines, sets or tags, resp., are potentially affected by memory accesses within the given address range \texttt{r}.

\begin{lstlisting}[style=code,float=t,caption={Event Counting of Class \texttt{WIR\_Cache}},label=lst:cachecnt]
// Get the number of counted cache hits.
@unsigned int@ WIR_Cache::getHitCount@( void )@ const;Ü\label{lst:cachecnt1}Ü

// Increase the number of counted cache hits by one.
@void@ WIR_Cache::increaseHitCount@( void )@ const;Ü\label{lst:cachecnt2}Ü

// Clear the number of counted cache hits.
@void@ WIR_Cache::clearHitCount@( void )@ const;Ü\label{lst:cachecnt3}Ü

// Get the number of counted cache misses.
@unsigned int@ WIR_Cache::getMissCount@( void )@ const;Ü\label{lst:cachecnt4}Ü

// Increase the number of counted cache misses by one.
@void@ WIR_Cache::increaseMissCount@( void )@ const;Ü\label{lst:cachecnt5}Ü

// Clear the number of counted cache misses.
@void@ WIR_Cache::clearMissCount@( void )@ const;Ü\label{lst:cachecnt6}Ü

// Get the number of counted unknown cache hits/misses.
@unsigned int@ WIR_Cache::getUnknownCount@( void )@ const;Ü\label{lst:cachecnt7}Ü

// Increase the number of counted unknown cache hits/misses by one.
@void@ WIR_Cache::increaseUnknownCount@( void )@ const;Ü\label{lst:cachecnt8}Ü

// Clear the number of counted unknown cache hits/misses.
@void@ WIR_Cache::clearUnknownCount@( void )@ const;Ü\label{lst:cachecnt9}Ü
\end{lstlisting}

A static cache analysis usually classifies each individual memory access as either a definite cache hit, a definite cache miss, or -- in the case of uncertainty -- as unknown cache hit/miss.
To support such counting of events, class \texttt{WIR\_Cache} provides the methods shown in \cref{lst:cachecnt}.
For each of the three event classes, there is a \texttt{get} (cf. \cref{lst:cachecnt1,lst:cachecnt4,lst:cachecnt7}), an \texttt{increase} (\cref{lst:cachecnt2,lst:cachecnt5,lst:cachecnt8}) and a \texttt{clear} (\cref{lst:cachecnt3,lst:cachecnt6,lst:cachecnt9}) method.


%-------------------------------------------------------------------------
\section{Buses}
\label{sec:buses}
%-------------------------------------------------------------------------

The very last system component that finally completes the description of the \gls{WIR} hardware model in this chapter are buses.
Their purpose is to connect a number of processor cores with various caches and/or physical memory regions.

\begin{wirFutureWork}
At the time of writing of this document, class \texttt{WIR\_Bus} is not yet designed and implemented in full detail, since it is not yet actively exploited by actual code analyses or optimizations.
The bus model that is currently implemented within \gls{WIR} actually boils down to the specification of bus arbitration policies and their configuration parameters.
Other relevant aspects like, e.g., arbitration and access delays, the distinction between single and burst accesses, or hierarchical buses with bridges in between are currently not modeled.
This is left up to future work.
\end{wirFutureWork}

\begin{lstlisting}[style=code,float=t,caption={Characteristic Attributes of Class \texttt{WIR\_Bus}},label=lst:busattribs]
// Get the maximum number of clock cycles for which a bus can delay a memory access.
virtual @unsigned int@ WIR_Bus::getMaxDelay@(@ const @unsigned int@ d æ=æ §0§ @)@ const;Ü\label{lst:busattribs1}Ü

// Get the maximum access time for a single memory access via this bus.
@unsigned int@ WIR_Bus::getMaxAccessTime@( void )@ const;Ü\label{lst:busattribs2}Ü

// Get a buses' arbitration strategy.
WIR_BusArbitrationType WIR_Bus::getArbitrationType@( void )@ const;Ü\label{lst:busattribs3}Ü

// Get a buses' arbitration configuration.
const WIR_BusArbitration æ&æWIR_Bus::getBusArbitration@( void )@ const;Ü\label{lst:busattribs4}Ü

// Enum BusArbitrationType represents the different arbitration strategies of buses.
enum class WIR_BusArbitrationType : @char@Ü\label{lst:busattribs5}Ü
{
  fp,           // Fixed-Priority arbitration.
  pd,           // Priority-Division arbitration.
  rr,           // Round-Robin arbitration.
  tdma          // Time-Division Multiple Access (TDMA) arbitration.
};
\end{lstlisting}

Similar to caches, class \texttt{WIR\_Bus} comes with a method \texttt{getMaxDelay} (see \cref{lst:busattribs1} of \cref{lst:busattribs}) that determines the maximal delay of a bus including that of subsequent parts of the memory hierarchy following behind the bus.
Argument \texttt{d} of this method denotes the maximal delay that can be caused by system components that are connected to the bus and that form the tail of a system's memory hierarchy.
Method \texttt{getMaxAccessTime} (\cref{lst:busattribs2}) returns the maximum number of clock cycles for a single memory access via a bus.
This value does not include the arbitration time potentially required by a bus ifself, it only reflects the delay stemming from an access to some memory module connected to the bus.
This method can be used to determine whether a pending memory access via a bus still has sufficient time left to complete within the currently running arbitration slot of the bus, or whether the access has to be carried over into some future slot.

\gls{WIR} currently supports four different bus arbitration strategies: Fixed-Priority, Priority-Division, Round-Robin, and \acrfull{TDMA}.
These arbitration policies are modeled by enumeration type \texttt{WIR\_BusAr\-bitrationType} (cf. \cref{lst:busattribs5} of \cref{lst:busattribs}), and a buses' actual arbitration policy can be determined using method \texttt{getArbitrationType} (cf. \cref{lst:busattribs3}).

Finally, method \texttt{getBusArbitration} (cf. \cref{lst:busattribs4}) provides access to the actual configuration parameters of the arbitration policy of a bus.
This method returns a reference to an object of class \texttt{WIR\_BusArbitration} which is simply an abstract base class from which specialized classes implementing, e.g., fixed-priority arbitration are inherited.
\cref{fig:wirbusarbitration} depicts the inheritance diagram of all currently supported \gls{WIR} bus arbitration policies.
These policies are described in more detail in the following \crefrange{ssec:fparbitration}{ssec:pdarbitration}.

\begin{figure}[ht]
  \begin{center}
    \includegraphics[width=0.82\textwidth]{busarbitration.png}
  \end{center}
  \caption{Inheritance Diagram of WIR Bus Arbitration Policies}
  \label{fig:wirbusarbitration}
\end{figure}


% - - - - - - - - - - - - - - - - - - - - - - - - - - - - - - - - - - - -
\subsection{Fixed-Priority Bus Arbitration}
\label{ssec:fparbitration}
% - - - - - - - - - - - - - - - - - - - - - - - - - - - - - - - - - - - -

Class \texttt{WIR\_FixedPriorityBusArbitration} realizes fixed-priority bus arbitration where each processor core has a fixed priority assigned.
Whenever a conflict occurs where several cores compete for access to the shared bus simultaneously, always that core with the highest fixed priority wins.
Currently ongoing bus transactions are, however, not preempted so that a bus transfer by some low-priority core is not interrupted if a request is made by a high-priority core.

Arbitration priorities are modeled using data type \texttt{Priority} which in turn goes back to the standard data type \texttt{unsigned int} (see \cref{lst:fpbus1} of \cref{lst:fpbus}).
The higher the numerical value represented by type \texttt{Priority} is, the higher a core's arbitration policy is.
Assuming that a total of $\mathit{n}$ cores are attached to a fixed-priority bus, the priorities can be from the interval $[0, \mathit{n}]$.
A priority of 0 can be used for cores that shall not be considered by the arbiter of a bus, which might be useful for special arbitration policies like, e.g., priority-division (cf. \cref{ssec:pdarbitration}).
Each processor core that is actively managed by a fixed-priority arbiter must carry its exclusive priority so that one priority value cannot be assigned to several processor cores (except priority 0).

\begin{lstlisting}[style=code,float=!b,caption={Characteristic Attributes of Class \texttt{WIR\_FixedPriorityBusArbitration}},label=lst:fpbus]
// Local type definition to represent fixed priorities of processor cores.
using WIR_FixedPriorityBusArbitration::Priority = @unsigned int@;Ü\label{lst:fpbus1}Ü

// Get a processor core's fixed arbitration priority.
Priority WIR_FixedPriorityBusArbitration::getPriority@(@ const WIR_BaseProcessor æ&æc @)@ const;Ü\label{lst:fpbus2}Ü

// Get all processor cores' fixed arbitration priorities for this bus.
ästd::mapäæ<æWIR_id_t, Priorityæ> &æWIR_FixedPriorityBusArbitration::getPriorities@( void )@ const;Ü\label{lst:fpbus3}Ü

// Get the total number of processor cores attached to this fixed-priority bus.
@unsigned int@ WIR_FixedPriorityBusArbitration::getNumberOfCores@( void )@ const;Ü\label{lst:fpbus4}Ü

// Get the total number of processor cores with non-zero arbitration priority.
virtual @unsigned int@ WIR_FixedPriorityBusArbitration::getNumberOfArbitratedCores@( void )@ const;Ü\label{lst:fpbus5}Ü

// Get the vector of cores managed by this bus arbiter, sorted by priority (highest priority first).
const ästd::vectoräæ<æWIR_id_tæ> &æWIR_FixedPriorityBusArbitration::getCoreOrder@( void )@ const;Ü\label{lst:fpbus6}Ü
\end{lstlisting}

Method \texttt{getPriority} (cf. \cref{lst:fpbus2}) provides the priority assigned to some actual processor core \texttt{c}.
A map with all priorities over all processor cores can be obtained using method \texttt{getPriorities} (\cref{lst:fpbus3}).
In the returned map, the unique identifiers of the processor cores (cf. \cref{sec:ids}) serve as keys.
The total number of cores that are attached to a fixed-priority bus can be retrieved via method \texttt{getNumberOfCores} (\cref{lst:fpbus4}).
This includes cores with arbitration priority 0, i.e., which are not actively managed by this bus.
The number of cores that effectively are managed by fixed priorities can be obtained with method \texttt{getNumberOfArbitratedCores} (\cref{lst:fpbus5}).
Finally, method \texttt{getCoreOrder} (\cref{lst:fpbus6}) can be used to obtain a vector with the unique identifiers of all processor cores, where this vector is sorted from highest priority cores at its head to lowest priorities at its tail.


% - - - - - - - - - - - - - - - - - - - - - - - - - - - - - - - - - - - -
\subsection{Round-Robin Bus Arbitration}
\label{ssec:rrarbitration}
% - - - - - - - - - - - - - - - - - - - - - - - - - - - - - - - - - - - -

A round-robin arbiter cyclically rotates the bus access among all cores, where each core possesses a unique number from the interval $[0, \mathit{n}-1]$ and $\mathit{n}$ denotes the total number of cores attached to a bus.
Whenever the currently active core $\mathit{n}_\mathit{c}$ releases the bus, $\mathit{n}_\mathit{c}$ is incremented to the next core currently requesting the bus.
If $\mathit{n}_\mathit{c}$ is advanced beyond the interval's maximal value $\mathit{n}-1$, a wrap around takes place and $\mathit{n}_\mathit{c}$ is set to 0.

This arbitration policy obviously does not require any configurable parameters at all so that the \gls{API} of the corresponding class \texttt{WIR\_RoundRobinBusArbitration} basically is empty.


% - - - - - - - - - - - - - - - - - - - - - - - - - - - - - - - - - - - -
\subsection{Time-Division Multiple Access Bus Arbitration}
\label{ssec:tdmaarbitration}
% - - - - - - - - - - - - - - - - - - - - - - - - - - - - - - - - - - - -

\gls{TDMA} arbitration realizes a repetitive cyclic schedule over time that consists of a fixed sequence of $\mathit{n}_\mathit{TDMA}$ \gls{TDMA} slots.
Each slot $\mathit{s} \in [0, \mathit{n}_\mathit{TDMA}-1]$ is characterized by its length $\mathit{l}_\mathit{s}$ over time as well as by a processor core $\mathit{o}_\mathit{s} \in [0, \mathit{n}-1]$ that is the exclusive owner of that \gls{TDMA} slot.
Again, $\mathit{n}$ represents the total number of processor cores connected to a \gls{TDMA}-arbitrated bus.
The total length $\mathit{l}_\mathit{TDMA}$ of one complete cycle of the schedule over all slots is $\mathit{l}_\mathit{TDMA} = \sum_{\mathit{s}=0}^{\mathit{n}_\mathit{TDMA}-1}\mathit{l}_\mathit{s}$.
The current position $\mathit{t}_\mathit{curr}$ in this cyclic schedule is determined by taking a global clock's time $\mathit{T}$ and computing its modulus to the schedule length: $\mathit{t}_\mathit{curr} = \mathit{T}\ \textrm{mod}\ \mathit{l}_\mathit{TDMA}$.
The use of the modulo operator here has the effect that once the current time reaches the end of the \gls{TDMA} schedule, it is reset to 0 so that the entire schedule restarts from the beginning, implementing the cyclic nature of \gls{TDMA} over time.
Based on these concepts, bus access is exclusively granted to that core $\mathit{o}_\mathit{s}$ that owns the \gls{TDMA} slot $\mathit{s}$ that is active at the current \gls{TDMA} time $\mathit{t}_\mathit{curr}$.
Simultaneous bus requests by other cores are rejected so that these other cores have to wait until the current time $\mathit{t}_\mathit{curr}$ has progressed to their very own dedicated \gls{TDMA} slot.

\begin{lstlisting}[style=code,float=t,caption={Characteristic Attributes of Class \texttt{WIR\_TDMABusArbitration}},label=lst:tdmabus]
// Type definition to represent relative times within cyclic TDMA schedules.
using WIR_TDMAOffset = ästd::uint16_tä;Ü\label{lst:tdmabus1}Ü

// Get the total number of cores arbitrated by this TDMA bus.
virtual @unsigned int@ WIR_TDMABusArbitration::getNumberOfArbitratedCores@( void )@ const;Ü\label{lst:tdmabus2}Ü

// Get properties of all TDMA slots of this bus.
const ästd::vectoräæ<æSlotInfoæ> &æWIR_TDMABusArbitration::getSlots@( void )@ const;Ü\label{lst:tdmabus3}Ü

// Get properties of all TDMA slots assigned to a given processor core.
const ästd::listäæ<æSlotInfoæ>æ WIR_TDMABusArbitration::getSlots@(@ const WIR_BaseProcessor æ&æc @)@ const;Ü\label{lst:tdmabus4}Ü

// Get the total number of TDMA slots of this bus.
@unsigned int@ WIR_TDMABusArbitration::getNumberOfSlots@( void )@ const;Ü\label{lst:tdmabus5}Ü

// Get the total length of the configured cyclic TDMA schedule.
WIR_TDMAOffset WIR_TDMABusArbitration::getScheduleLength@( void )@ const;Ü\label{lst:tdmabus6}Ü
\end{lstlisting}

\gls{WIR} class \texttt{WIR\_TDMABusArbitration} models this behavior of a \gls{TDMA}-arbitrated bus, and its key attributes can be inspected using the methods shown in \cref{lst:tdmabus}.
Method \texttt{getNumberOfArbitratedCores} (cf. \cref{lst:tdmabus2}) provides the total number $\mathit{n}$ of processor cores that are connected to and arbitrated by a \gls{TDMA} bus.
The specific information about all slots making up the entire cyclic \gls{TDMA} schedule can be retrieved via method \texttt{getSlots} from \cref{lst:tdmabus3}.
This variant of \texttt{getSlots} returns a reference to a vector containing the information about each \gls{TDMA} slot.
This slot information is represented by class \texttt{WIR\_TDMABusArbitration::SlotInfo} described below in more detail.
In order to obtain information about all \gls{TDMA} slots associated with a given processor core, the \texttt{getSlots} variant from \cref{lst:tdmabus4} can be used.
It expects a reference \texttt{c} to the processor core under consideration as argument and returns a list containing only the information of all slots exclusively owned by core \texttt{c}.
The total number $\mathit{n}_\mathit{TDMA}$ of \gls{TDMA} slots can be obtained via method \texttt{getNumberOfSlots} (\cref{lst:tdmabus5}).
Finally, method \texttt{getScheduleLength} (\cref{lst:tdmabus6}) provides the overall duration $\mathit{l}_\mathit{TDMA}$ of one full \gls{TDMA} cycle.
Such relative times within a \gls{TDMA} schedule are represented by type \texttt{WIR\_TDMAOffset}, which in turn is declared to be \texttt{std::uint16\_t} (cf. \cref{lst:tdmabus1}).

Class \texttt{SlotInfo} is a sub-class of \texttt{WIR\_TDMABusArbitration} modelling all relevant properties of a \gls{TDMA} slot.
These properties can be obtained using the methods from \cref{lst:tdmaslot}.
Method \texttt{getIndex} returns the position of a slot within the full \gls{TDMA} schedule, i.e., its index from the interval $[0, \mathit{n}_\mathit{TDMA}-1]$.
A reference to a slot's exclusively owning processor core is returned by method \texttt{getOwner} (\cref{lst:tdmaslot2}).
The point in time where a slot starts relative to the beginning of the full \gls{TDMA} schedule can be obtained via method \texttt{getStart} (\cref{lst:tdmaslot3}).
Finally, method \texttt{getLength} (\cref{lst:tdmaslot4}) provides a slot's length $\mathit{l}_\mathit{s}$.

\begin{lstlisting}[style=code,float=t,caption={Characteristic Attributes of Class \texttt{WIR\_TDMABusArbitration::SlotInfo}},label=lst:tdmaslot]
// Get the position of a slot within the entire cyclic TDMA schedule.
@unsigned int@ WIR_TDMABusArbitration::SlotInfo::getIndex@( void )@ const;Ü\label{lst:tdmaslot1}Ü

// Get the processor core owning a TDMA slot.
const WIR_BaseProcessor æ&æWIR_TDMABusArbitration::SlotInfo::getOwner@( void )@ const;Ü\label{lst:tdmaslot2}Ü

// Get a TDMA slot's start time, counted from the beginning of the entire TDMA schedule.
WIR_TDMAOffset WIR_TDMABusArbitration::SlotInfo::getStart@( void )@ const;Ü\label{lst:tdmaslot3}Ü

// Get a TDMA slot's length in clock cycles.
WIR_TDMAOffset WIR_TDMABusArbitration::SlotInfo::getLength@( void )@ const;Ü\label{lst:tdmaslot4}Ü
\end{lstlisting}


% - - - - - - - - - - - - - - - - - - - - - - - - - - - - - - - - - - - -
\subsection{Priority-Division Bus Arbitration}
\label{ssec:pdarbitration}
% - - - - - - - - - - - - - - - - - - - - - - - - - - - - - - - - - - - -

\gls{TDMA} is criticized in the literature for its inefficiency, because a slot where its exclusively owning processor core does not request the bus, remains unused, thus leading to idle periods.
Priority-division arbitration is a generalization of \gls{TDMA}, it can be seen as a combination of \gls{TDMA} and fixed-priority arbitration.
Under a priority-division policy, all slots are either exclusively owned by one processor core (as is the case for \gls{TDMA}, cf. \cref{ssec:tdmaarbitration}), or priorities over various cores are defined (like in fixed-priority arbitration, cf. \cref{ssec:fparbitration}).
When it comes to bus requests and arbitration decisions, the arbiter again considers that slot $\mathit{s}$ that is currently active at time $\mathit{t}_\mathit{curr}$ relative to the cyclic schedule's very beginning.
If this slot is configured to have one exclusive owner $\mathit{o}_\mathit{s}$ and this owner currently requests the bus, bus access is granted to $\mathit{o}_\mathit{s}$.
If owner $\mathit{o}_\mathit{s}$ currently does not request the bus, the current slot $\mathit{s}$ remains idle.
However, if the current slot is configured to use priorities, bus access is granted to that processor core having the highest priority among all competing cores.

This arbitration scheme is realized by class \texttt{WIR\_PriorityDivisionBusArbitration} within \gls{WIR}.
Multiple inheritance from classes \texttt{WIR\_FixedPriorityBusArbitration} and \texttt{WIR\_TDMABusArbitration} realizes the combination of fixed-priority and \gls{TDMA} arbitration mentioned above (cf. \cref{fig:wirbusarbitration}).
Since the \glspl{API} of both base classes are already described earlier in \cref{ssec:fparbitration,ssec:tdmaarbitration}, there remain only two details for priority-division arbitration.

\begin{lstlisting}[style=code,float=h,caption={Characteristic Attributes of Class \texttt{WIR\_PriorityDivisionBusArbitration}},label=lst:pdbus]
// Get the total number of cores arbitrated by this priority-division bus.
virtual @unsigned int@ WIR_PriorityDivisionBusArbitration::getNumberOfArbitratedCores@( void )@ const override;Ü\label{lst:pdbus1}Ü

// Get the operation modes per bus slot.
const ästd::vectoräæ<æWIR_BusArbitrationTypeæ> &æWIR_PriorityDivisionBusArbitration::getSlotModes@( void )@ const;Ü\label{lst:pdbus2}Ü
\end{lstlisting}

First, method \texttt{getNumberOfArbitratedCores} (cf. \cref{lst:pdbus1} of \cref{lst:pdbus}) already known from both fixed-priority and \gls{TDMA} arbitration is overridden such that it returns the number of cores that exclusively own some bus slot, or that have non-zero fixed priorities.
Second, priority-division introduces the novel method \texttt{getSlotModes} (cf. \cref{lst:pdbus2}).
This method returns a reference to a vector storing how each slot is actually behaving -- either as fixed-priority or as \gls{TDMA} slot.
The elements of the returned vector follow the same slot order as provided by \texttt{getSlots} from \cref{lst:tdmabus3} from \cref{lst:tdmabus}.
The information which policy is actually used per slot is encoded using the enumeration type \texttt{WIR\_BusArbitrationType} already introduced in \cref{lst:busattribs5} of \cref{lst:busattribs}.

\cleardoubleoddpage
